\section{1. Языки и грамматики}\label{1.}
\begin{frame}
\frametitle{1. Языки и грамматики }

Начинаем мы с самого заслуженного и 
до сих пор самого популярного средства описания синтаксиса формальных языков, имеющих текстовое представление~---
 \odmkey{теории формальных
грамматик}{Теория!формальных!грамматик}.


\medskip
Для особенно простых, но часто применяемых в программировании
формальных языков, которые называются регулярными,
имеется средство описания, которое называется
\odmkey{регулярные выражения}{Выражение!регулярное}.

\medskip
В этой главе приводятся основные определения, факты и алгоритмы из
теории формальных грамматик, излагается техника регулярных выражений
и рассматривается применение описанного аппарата для работы
с языками программирования, в частности,
с языками предметных областей.

\end{frame}


\subsection{1.1. Порождающие грамматики}
\label{1.1.}
\begin{frame}
\frametitle{1.1. Порождающие грамматики (1/3)}
Первым рассматривается простейший случай, когда система знаков~--- это
конечный алфавит символов, а способ
комбинирования~--- это составление линейных последовательностей
символов, то есть {\em текстов}.

\bigskip
\begin{digress}
Исторически идеи порождающих грамматик возникли 
применительно к текстам, но в настоящее время
развились до применения к изображениям, звукам, жестам и многим 
другим символам различной природы. 
Некоторые примеры рассматриваются в главе 3.  
\end{digress}
\end{frame}

\begin{frame}
\frametitle{1.1. Порождающие грамматики (2/3)}

\begin{tabular}{p{10mm}|p{40mm}|p{95mm}}
  № & Название параграфа
  & Ключевые термины и обозначения \\
  \hline
 \DMref{1.1.1.}{1.1.1.} &  Алфавит, слово и язык
  & {  
  алфавит, буква, символ, слово, цепочка, строка, вхождение, пустое слово $\varepsilon$,
  вхождение, вырезка, длина слова, конкатенация, сцепление, степень, префикс, суффикс,
  пересечение слов, язык, лексикографический порядок }
   \\
\DMref{1.1.2.}{1.1.2.} & Формальные грамматики 
  & 
{ порождающая грамматика,
  нетерминал, переменная, терминал, константа,
  правило вывода, продукция, аксиома грамматики,
  вывод, непосредственный вывод, бесповторный вывод, 
  порождаемый язык $\frak{L}$,  допуск слова,
  эквивалентность грамматик  
}
  \\
\DMref{1.1.3.}{1.1.3.}  & Перечислимость языков &
\\
\DMref{1.1.4.}{1.1.4.}  & Классификация грамматик и языков &
{Ноам Хомский, грамматика: с фразовой структу\-рой, непосредственных составляющих,
 контекстно-зависи\-мая, контекстно-свободная,
 регулярная, $\frak{T}_i$
 } 
\end{tabular}

\end{frame}


\begin{frame}
\frametitle{1.1. Порождающие грамматики (3/3)}

\begin{tabular}{p{10mm}|p{40mm}|p{95mm}}
  № & Название параграфа
  & Ключевые термины и обозначения \\
  \hline
\DMref{1.1.5.}{1.1.5.} & Неукорачивающие грамматики &
{разрешимый, полуразрешимый, легко разрешимый язык, распознающий алгоритм}\\   
\DMref{1.1.6.}{1.1.6.} & Контекстно-свободные грамматики &
{левосторонний вывод, правосторонний вывод, дерево разбора, дерево вывода,
дерево порождения, неоднозначная грамматика}\\   
\DMref{1.1.7.}{1.1.7.} & Линейные, регулярные и автоматные грамматики &
{правило: начальное, заключительное, леволинейное, праволинейное, 
порождающие и недостижимые символы, цепное правило, 
$\varepsilon$-правило, $S$-правило, приведённая грамматика }\\   
\DMref{1.1.8.}{1.1.8.} & Языки записи грамматик &
{ метасимвол, нормальная форма Бэкуса--Наура,
металингвистическая переменная, металингвистическая формула }\\   
\DMref{1.1.9.}{1.1.9.} & Альтернативные средства задания синтаксиса &
{Никлаус Вирт, диаграмма, тривиальная диаграмма, элементарная диаграмма, 
язык разметки   }\\   

\end{tabular}

\end{frame}


\subsubsection{1.1.1. Алфавит, слово и язык}\label{1.1.1}
\begin{frame}
\frametitle{1.1.1. Алфавит, слово и язык (1/9)}
Пусть задано конечное множество $A =\{\Tuple{a}{1}{n}\}$, которое
называется \odmkey{алфавитом}{Алфавит}. Элементы алфавита называются
\odmkey{буквами}{Буква} или \odmkey{символами}{Символ}.
Обычно отдельные символы обозначаются
начальными строчными буквами латинского алфавита.
Последовательность букв называется
\odmkey{словом}{Слово} (в данном алфавите).
Заметим, что в программировании слово часто называют
\odmkey{строкой}{Строка}, а в теории формальных языков~--- \odmkey{цепочкой}{Цепочка}.
Буквы в слове упорядочены (слева направо),
и могут быть перенумерованы.
Одна буква может входить
в слово несколько раз, каждый
отдельный экземпляр буквы в слове называется
\odmkey{вхождением}{Вхождение} буквы в слово.
Вхождения могут быть идентифицированы,
например, номером буквы в слове 
 относительно начала слова.
Обычно слова обозначают
строчными буквами греческого алфавита,
а номера записывают в квадратных скобках.
\begin{example}
Пусть $A=\{a,b\}$. Тогда 
$\alpha=aabb$~--- слово, 
а $\alpha[3]=b$~--- буква.
\end{example}

\begin{remark}
Вхождения можно идентифицировать не номером,
а смещением~--- это вопрос вкуса, который регулируется предварительным соглашением.
\end{remark}

Количество вхождений букв в слово называется \odmkey{длиной}{Длина!слова} слова:
$|\alpha| = |a_1\dots a_k| = k$.
\begin{example} $|$Ура-а-а!$| = 8$.
\end{example}

\odmkey{Пустое}{Слово!пустое} слово обозначается
$\varepsilon$,
$|\varepsilon|= 0$.

\end{frame}


\begin{frame}
\frametitle{1.1.1. Алфавит, слово и язык (2/9)}
Можно рассматривать не только вхождения
отдельных букв в слово, но и 
вхождения любых частей слова как отдельные слова.
Операция взятия части слова называется 
\odmkey{вырезкой}{Операция вырезки}.
Вырезку естественно также обозначить 
квадратными скобками. 
\begin{example}
самовар$[3..6]= $ мова.
\end{example}

Таким образом, определено отношение 
между словами <<быть частью>>.
Это отношение рефлексивно, транзитивно и антисимметрично,
а потому является нестрогим частичным порядком,
и обозначается здесь $\subset$.
\begin{example}
луг $\subset$ слуга.
\end{example}

\begin{remark} Очевидно, что 
$\alpha\subset\beta\Impl|\alpha|\le|\beta|$,
но $\Not(|\alpha|\le|\beta|\Impl\alpha\subset\beta)$!
\end{remark} 
\begin{example}
гул $\not\subset$ слуга.
\end{example}

Для слов определена операция 
\odmkey{конкатенации}{Операция!конкатенации}, 
или \odmkey{сцепления}{Операция!сцепления}.
Обычно операция конкатенации никак не обозначается, а сцепляемые слова 
просто записываются подряд. 

\begin{example}
Конкатенация слов <<пара>> и <<план>> образует слово <<параплан>>.
\end{example}

Сцепление можно применять многократно,
таким образом определяется \odmkey{степень}{Степень!слова} слова. Положим
$\alpha^1 \Def \alpha; \alpha^k\Def\alpha\alpha^{k-1}; k>1$. Можно 
определить нулевую степень слова:
$\alpha^0 \Def \varepsilon$.

\begin{example}
$a^4=aaaa; (ab)^3= ababab$.
\end{example}

\end{frame}


\begin{frame}
\frametitle{1.1.1. Алфавит, слово и язык (3/9)}

Одни и те же буквы могут входить в разные слова.
Это порождает отношения между словами, 
связанные с наличием или отсутствием общих частей. 

Если $\alpha=\alpha_1\alpha_2$, то $\alpha_1$ называется
\odmkey{началом}{Начало слова}, или \odmkey{префиксом}{Префикс},
слова $\alpha$, а $\alpha_2$~---
\odmkey{окончанием}{Окончание слова}, или
\odmkey{постфиксом}{Постфикс}, слова $\alpha$.
Если при этом $\alpha_1 \neq \varepsilon\And \alpha_2\neq\varepsilon$, то
$\alpha_1$ (соответственно, $\alpha_2$)
называется \odmkey{собственным началом}{Начало слова!собственное}
(соответственно, \odmkey{собственным окончанием}{Окончание
слова!собственное})
слова $\alpha$.

\begin{example}
Слова <<д>> и <<до>> являются собственными префиксами, а слова <<ом>> и <<м>> являются собственными постфиксами слова <<дом>>.
\end{example}

\begin{remark}
Каждое слово является своим постфиксом и префиксом (несобственным!).
\end{remark}

\medskip
Слова $\alpha$ и $\beta$ могут \odmkey{пересекаться}{Операция!пересечения слов}.
При пересечении либо первое слово является частью второго, 
то есть $\beta=\gamma\alpha\delta$,
либо постфикс первого слова является префиксом второго, 
то есть $\alpha=\delta\gamma, \beta=\gamma\zeta$.
Для пары слов $\alpha$ и $\beta$ тройка слов 
$\gamma$, $\delta$, $\zeta$ определяется однозначно,
если выбирать {\em наибольшее} пересечение

\begin{example}
Пусть $\alpha=abb$, $\beta=bbc$. Тогда можно
взять $\gamma\Assign b$, $\delta\Assign ab$, $\zeta\Assign bc$.
Однако наибольшее пересечение даёт единственная тройка 
$\gamma\Assign bb$, $\delta\Assign a$, $\zeta\Assign c$.
\end{example} 

Таким образом, на множестве слов определена
{\em операция} \odmkey{пересечения}{Пересечение!слов}. 
Эта операция не коммутативна! 
\end{frame}


\begin{frame}
\frametitle{1.1.1. Алфавит, слово и язык (4/9)}
Множество слов в алфавите
${A}$ обозначается~${A^*}$. 
По определению $\varepsilon \not\in A$, но $\varepsilon\in A^*$.
Если пустое слово исключается из рассмотрения, 
то множество слов в алфавите $A$ 
обозначается~$A^+$.

\smallskip
Произвольное множество $L$ слов в алфавите $A$ называется 
\odmkey{языком}{Язык (language)} в этом алфавите.

\begin{examples}
\begin{enumerate}
\item $A\subset A^*$ --- язык, состоящий из
однобуквенных слов;
\item $\{0,1\}^*$ --- язык двоичных последовательностей;
\item $A^*\subset A^*$ --- универсальный
язык всех слов.
\end{enumerate}
\end{examples}
 
К любым языкам применимы операции над множествами: объединение, пересечение,
разность и симметрическая разность.
Поскольку алфавит языка известен, то определена и
операция дополнения.

\begin{digress}
Обычные операции наивной теории множеств
играют сравнительно небольшую роль в 
теории языков, за исключением
операции объединения.
\end{digress}

\smallskip
Пусть есть два языка над одним алфавитом:
$L_1 \subset A^*, L_2\subset A^*$.  
Тогда операция сцепления слов индуцирует операцию 
сцепления языков: $L_1 L_2 
\Def \Set{\alpha\beta}{\alpha\in L_1 \And
\beta\in L_2}.$
\end{frame}

\begin{frame}
\frametitle{1.1.1. Алфавит, слово и язык (5/9)}
Сцепление языков можно применять многократно,
таким образом определяется \odmkey{степень}{Степень!слова} языка. Пусть $L\subset A^*$. Положим
$L^1 \Def L; L^k\Def L L^{k-1}; k>1$. Можно 
определить нулевую степень языка:
$L^0 \Def \{\varepsilon\}$.

Но сам алфавит $A$~--- это язык,
и значит определены степени этого языка $A^k$
для любого натурального числа $k$. 
Из этого наблюдения и определений
немедленно следует очевидная, 
но важная теорема.

\begin{theorem}
$A^k=\Set{\alpha\in A^*}{|\alpha|=k}$.
\end{theorem}

\begin{remark}
Другими словами, $A^k$~--- это множество слов длины $k$ в алфавите $A$.
\end{remark}

\begin{corollary}
$|A^k|=|A|^k =n^k.$
\end{corollary} 

Рассмотрим $\cal{A} = \{ A^0, A^1, A^2, \dots\}$~--- семейство классов эквивалентности
по отношению равенства длин слов. Тогда
$$A^* = \bigcup_{k=0}^\infty A^k, \qquad
A^+ = \bigcup_{k=1}^\infty A^k,$$
что вполне согласуется с предшествующими определениями.
\end{frame}

\begin{frame}
\frametitle{1.1.1. Алфавит, слово и язык (6/9)}
Множество $A^*$ бесконечно, но счётно, то есть 
его можно занумеровать. Рассмотрим некоторые нумерации, которые особенно часто используются
в других построениях.

Поскольку алфавит $A=\{\Tuple{a}{1}{n}\}$ конечен,
его буквы можно перенумеровать (то есть 
{\em линейно упорядочить}) и представить алфавит словом 
$\alpha\!\beta_A\Def a_1\dots a_{n}$,
то есть конкатенацией всех букв в порядке
возрастания номеров.


\begin{digress}
Слово <<алфавит>> происходит от прочтения
первых букв греческого алфавита:
альфа, бета, ... Аналогичное происхождение
имеет русское слово <<азбука>>: аз, буки, веди, ... \\Обозначение $\alpha\!\beta_A$ допустимо читать <<алфавит А>>.
\end{digress}

\begin{example}
Современная кириллица содержит 33 буквы,
и для этих букв узаконен следующий порядок:
$\alpha\!\beta_{\cal{K}}\Assign$ абвгдеёжзийклмнопрстуфхцчшщъыьэюя. 
Индекс $\cal{K}$ символизирует кириллицу.
\end{example}

Зафиксируем биекцию
$A=\{\Tuple{a}{1}{n}\} \sim 1..n$ следующим образом.

\vspace{-16pt}
\begin{columns}[t]
\begin{column}{10cm}
\begin{algorithmic}
\STATE \textbf{func $n_A(a : A) : 1..n$}
\FOR {\textbf{$i$ from $1$ to $n$}}
\STATE{\textbf{if $\alpha\beta_A[i]=a$ then return 
$i$ end if}}
\ENDFOR
\end{algorithmic}
\end{column}    
\begin{column}{5cm}
\begin{algorithmic}
\STATE \textbf{func $a_A(i : 1..n) : A$}
\STATE \textbf{return $\alpha\!\beta_A[i]$}
\end{algorithmic}
\end{column}
\end{columns}    


\end{frame}

\begin{frame}
\frametitle{1.1.1. Алфавит, слово и язык (7/9)}

Нумерация букв индуцирует нумерацию слов.
Пусть $\alpha \in A^k$. Положим 
$$n_A(\alpha)\Assign n_A(\alpha[k])
+ n \cdot n_A(\alpha[k-1]) +\dots
+n^{k-1}\cdot n_A(\alpha[1]) = 
\sum_{i=1}^{k} n_A(\alpha[i]) n^{k-i}. $$
%Тогда номером слова оказывается число,
%записанное в $n$-ричной системе счисления,
%причём младшие цифры записываются слева,
%а старшие --- справа ({\em little-endian}).
%

\vspace{-6pt}
\begin{example}
В этой нумерации для слов в кириллице имеем:
 
\begin{tabular}{l|l|l}
\hline
$n_{\cal{K}}(\text{а})=1$ & $n_{\cal{K}}(\text{аа})=1\cdot 33+1=~~34$ & $n_{\cal{K}}(\text{ааa})=1\cdot 33^2 +1\cdot 33 +1=~1123$  \\
$n_{\cal{K}}(\text{б})=2$ & $n_{\cal{K}}(\text{аб})=1\cdot 33+2=~~35$ & $n_{\cal{K}}(\text{ааб})=1\cdot 33^2+1\cdot 33+2=~1124$  \\
$\dots$ & $\dots$ & $\dots$ \\
$n_{\cal{K}}(\text{я})=33$ & $n_{\cal{K}}(\text{яя})=33\cdot 33+33=1122$ & $n_{\cal{K}}(\text{яяя})=33\cdot 33^2+ 33\cdot 33+33=37059$ \\
\hline
\end{tabular}
\end{example}

\medskip
\begin{theorem}
$\A{A}{A^*\sim \Bbb{N}_0}$.
\end{theorem}

\begin{proof}
Представленная нумерация реализует требуемую 
биекцию. Номер слова
легко вычисляется известным алгоритмом позиционной системы счисления, представленным
на следующем слайде слева.
Обратный алгоритм определения слова по номеру чуть сложнее, поскольку особым образом требуется
обрабатывать числа, кратные основанию $n$. Алгоритм представлен на следующем слайде справа.
Для завершения доказательства остаётся положить $n_A(\varepsilon)\Assign 0$.  
\end{proof}
 
\end{frame}

\begin{frame}
\frametitle{1.1.1. Алфавит, слово и язык (8/9)}

\vspace{-12pt}
\begin{columns}[t]
\begin{column}{5cm}
\begin{algorithmic}
\STATE{\textbf{func $n_A(\alpha : A^*) : 
\Bbb{N}_0$}}
\STATE{$i\Assign 0$}
\FOR {\textbf{$j$ from $1$ to $|\alpha|$ }}
\STATE{$i\Assign n\cdot i +n_A(\alpha[j])$}
\ENDFOR
\STATE{\textbf{return $i$}}
\end{algorithmic}
\end{column}    
\begin{column}{10cm}
\begin{algorithmic}
\STATE \textbf{func $a_A(i : \Bbb{N}_0) : A^*$}
\STATE{$\alpha\Assign \varepsilon$}
\WHILE {$i>0$}
\STATE\textbf{if $i\MOD n =0 $ then $a\Assign a_A(n); i\Assign i-n$}
\STATE\textbf{ else $a\Assign a_A(i\MOD n)$ end if}
\STATE{$\alpha\Assign a \alpha$; $i\Assign i\DIV n$}
\ENDWHILE
\STATE{\textbf{return $\alpha$}}
\end{algorithmic}
\end{column}
\end{columns}    

\medskip
\begin{remark}
Предложенная нумерация обладает рядом приятных свойств. \\Во-первых, нумерация монотонна по длине слова, то есть более длинные слова имеют номера больше, чем более короткие. 
\\Во-вторых, нумерация локально взаимно-однозначна:\\
$A\sim 1..n$, $A^2\sim (n+1)..(n^2+n)$ и вообще
$\A{k\in \Bbb{N}}{A^k \sim \frac{n^{k}-1}{n-1}.. n\frac{n^k-1}{n-1}}
$.\\
В-третьих, нумерация определяет естественный
линейный порядок на множестве слов:\\
$\alpha<\beta \Assign n_A(\alpha)<n_A(\beta)$,
что позволяет перечислить все слова в порядке возрастания.  
\end{remark}


\end{frame}

\begin{frame}
\frametitle{1.1.1. Алфавит, слово и язык (9/9)}
Часто на множестве слов используют 
\odmkey{лексикографический порядок}{Порядок!лексикографический} (или \odmkey{алфавитный}{Порядок!алфавитный}):\\
$\alpha<\beta\Def \beta=\alpha\gamma\And\gamma\neq\varepsilon 
\Or \E{j\le|\alpha|}{\alpha[j]<
\beta[j]\And\A{i \in 1..(j-1)}{\alpha[i]=
\beta[i]}}.$ 
Лексикогра\-фи\-ческий порядок линеен, так же как и введённый естественный порядок.  

\begin{remark}
Лексикографический порядок не позволяет построить генератор, который перечисляет
{\em все} слова в возрастающем порядке. 
\end{remark}

\begin{example}
В естественном порядке 
$ 1 < 2 <11 <12< 21 $, 
а в лексикографическом порядке 
$1<11<12<2<21$. 
\end{example}

\begin{digress}
Нумерация $n_A()$ очень похожа на 
обычную позиционную
систему счисления с основанием $n$. 
Отличие в том, что в позиционных системах
цифры берутся из диапазона $0..(n-1)$,
а в предложенной системе цифры берутся из
диапазона $1..n$. Наличие цифры $0$ с её 
особыми арифметическими свойствами влечёт два важных следствия. С одной стороны,
алгоритмы 
всех арифметических операций в обычной 
системе счисления
оказываются проще, чем алгоритмы 
в системе без нуля,
поэтому в компьютерах при вычислениях всегда применяются позиционные системы с нулём.
С другой стороны, из-за нуля
теряется однозначность соответствия
между словами, составленными из цифр, 
и натуральными числами. Проблему
создают \odmkey{ведущие нули}{Нуль!ведущий}~---
различные строки $0001$ и $1$ в позиционной системе соответствуют одному и тому же
числу <<один>>.
\end{digress}

\end{frame}

\subsubsection{1.1.2. Формальные грамматики}\label{1.1.2.}
\begin{frame}
\frametitle{1.1.2. Формальные грамматики (1/4)}

Если язык содержит только конечное множество
слов, то их можно задать просто перечислением.
Если же язык бесконечен, что чаще встречается,
то необходимо использовать специальные средства
для построения конечного описания бесконечного множества.
Возникновение и развитие самого известного из таких средств 
связывают с Ноамом Хомским.
  
\smallskip 
\odmkey{Формальная порождающая грамматика}{Грамматика!формальная}
\odmkey{}{Грамматика!порождающая} (далее просто грамматика)
$G=\langle \mathrm{M}, \Sigma, \Pi, S\rangle$~--- это совокупность конечного
алфавита $\mathrm{M}$ \odmkey{нетерминальных}{Символ!нетерминальный}
символов (то есть вспомогательных символов, не входящих в алфавит, над которым определяется язык), конечного алфавита $\Sigma$
\odmkey{терминальных}{Символ!терминальный} символов (то есть  символов алфавита, над которым определяется язык, 
$\mathrm{M}\cap \Sigma = \emptyset$), конечного множества \odmkey{правил вывода}{Правило!вывода}
(правила вывода называют также \odmkey{правилами подстановки}{Правило!подстановки}
или \odmkey{продукциями}{Продукция})
$\Pi=\Set{\alpha\rightarrow\beta}{\alpha,\beta\in (\mathrm{M}\cup \Sigma)^*}$ и
\odmkey{основного символа}{Символ!основной} (или \odmkey{аксиомы
грамматики}{Аксиома!грамматики}) $S \in \mathrm{M}$.


\begin{example}
Совокупность множеств \\
$\mathrm{M}\Assign \{N\}$,
$\Sigma\Assign\{a,b\}$,
$\Pi\Assign\{N\to NN, N\to \varepsilon, N\to a,
N\to b\}$ \\ образует грамматику  
$G_0=\langle \mathrm{M},\Sigma, \Pi, N\rangle$.
\end{example}

\end{frame}

\begin{frame}
\frametitle{1.1.2. Формальные грамматики (2/4)}

\begin{remark}
По традиции нетерминальные символы (или просто
\odmkey{нетерминалы}{Нетерминал}, 
иногда их также называют \odmkey{переменными}{Переменная}) обозначают прописными латинскими
буквами, терминальные символы (или просто
\odmkey{терминалы}{Терминал}) обозначают строчными латинскими
буквами, а цепочки терминальных и нетерминальных символов~---
строчными греческими буквами.
Было бы неудобно {\em множества} символов, то есть алфавиты, 
обозначать символами этих же алфавитов.
Поэтому здесь приняты обозначения алфавитов и множества правил прописными
греческими буквами, по первым буквам греческих слов
$\mathrm{M}\varepsilon\tau\alpha\beta\lambda\eta\tau\eta$ (Переменная),
$\Sigma\tau\alpha\theta o\rho o\varsigma$ (Постоянная),
$\Pi\alpha\rho\alpha\gamma\omega\gamma\iota\kappa o\varsigma
\ \nu o\mu o\varsigma$ (Производящее правило).
\end{remark}

\medskip
Поскольку все множества конечны, 
запись грамматики также конечна,
однако конечная грамматика может порождать
бесконечный язык на основе 
формального \odmkey{вывода}{Вывод}.

\medskip   
Говорят, что из цепочки $\gamma$ \odmkey{непосредственно
выводима}{Выводимость!непосредственная} цепочка $\delta$ в
грамматике $G$ (обозначение $\gamma\Rightarrow_G \delta$), если
$\gamma=\gamma_1\alpha\gamma_2$, $\delta=\gamma_1\beta\gamma_2$ и
$\alpha\rightarrow\beta\in \Pi$. Говорят, что из цепочки $\gamma$
\odmkey{выводима}{Выводимость} цепочка $\delta$ в грамматике $G$,
если существует \odmkey{вывод}{Вывод}, т.е. последовательность цепочек, в которой каждая следующая выводима из предыдущей по одному из правил вывода: 
$$\E{\Tuple{\sigma}{0}{n}}
{\A{i\in 0..(n-1)}{\sigma_{i}\Rightarrow_G\sigma_{i+1} \And \sigma_0 = \gamma
\And \sigma_n = \delta}}.$$ 
\end{frame}

\begin{frame}
\frametitle{1.1.2. Формальные грамматики (3/4)}

Число $n$ (количество цепочек в
последовательности, не считая первой) называется
\odmkey{длиной}{Длина!вывода} вывода. Выводимость $\delta$ из
$\gamma$ обозначают 
$\gamma\Rightarrow^n_G\delta$, если нужно
указать длину вывода или $\gamma\Rightarrow^*_G\delta$, если длина
вывода не указывается. Если грамматика подразумевается, то индекс
$G$ опускают.
Отметим, что $\Rightarrow_G$ является
отношением на множестве цепочек.
В таком случае отношение
 $\Rightarrow_{G}^{*}$ является транзитивным замыканием отношения $\Rightarrow_{G}$ на множестве цепочек.

\smallskip
\begin{example}
Последовательность формул $NN, NNN,  NNNN, aNNN, abNN, abbN, abba $ является выводом в грамматике $G_0$.
Другими словами, $NN\Rightarrow^6_{G_0} abba$.
\end{example}

\smallskip
\begin{remark}
Таким образом, грамматика~--- 
это {\em формальная теория}, или же
{\em исчисление}, в том строгом значении,
которое имеют эти термины в математической логике.
Теоремами этой теории являются терминальные цепочки, выводимые из аксиомы грамматики.   
\end{remark}

\smallskip
\begin{example}
В грамматике $G_0$ предыдущего примера
все слова из множества $\{a,b\}^*$, включая пустое слово, выводимы
из начального символа грамматики.
Другими словами, грамматика $G_0$~---
это такая малоинтересная теория,
в которой все формулы суть теоремы.   
\end{example}
\end{frame}

\begin{frame}
\frametitle{1.1.2. Формальные грамматики (4/4)}
Вывод называется \odmkey{бесповторным}{Вывод!бесповторный}, если
все цепочки в нем различны.

\begin{lemma}
Если существует вывод слова $\alpha$ в грамматике $G$, то существует бесповторный
вывод слова $\alpha$ в грамматике $G$.
\end{lemma}
\begin{proof}
Пусть
$\sigma_0=S,\sigma_1,\dots,\sigma_i,\dots,\sigma_j,\dots\sigma_n=\alpha$~---
вывод цепочки $\alpha$ в $G$, причём $\sigma_i = \sigma_j$. Очевидно, что
$\sigma_0=S,\sigma_1,\dots,\sigma_i,\sigma_{j+1},\dots\sigma_n=\alpha$
также будет выводом цепочки $\alpha$ в $G$.
\end{proof}

\smallskip
\odmkey{Языком}{Язык!порождаемый грамматикой}, порождаемым
грамматикой $G$ (обозначение $\frak{L}(G)$) называется множество всех
цепочек в терминальном алфавите, выводимых из начального символа грамматики:
$$
\frak{L}(G)\Def \Set{\alpha\in \Sigma^*}{S\Rightarrow_G^*\alpha}.
$$
\begin{remark}
В традициях терминологии логических исчислений говорят, что слова языка $\frak{L}(G)$ \odmkey{допускаются}{Слово!допустимое} грамматикой $G$.
\end{remark}

Один и тот же язык может быть порождён
разными грамматиками.
Говорят, что две грамматики $G$ и $G'$
\odmkey{эквивалентны}{Эквивалентность!грамматик} (обозначение
$G\sim G'$), если $\frak{L}(G)=\frak{L}(G')$.

\begin{example}
Пусть $\mathrm{M}\Assign \{N\}$,
$\Sigma\Assign\{a,b\}$,
${\Pi}'\Assign\{N\to \varepsilon, N\to aN,
N\to bN\}$, тогда грамматика  
${G_0}'=\langle \mathrm{M}, \Sigma, {\Pi}', N\rangle$ эквивалентна $G_0$,
то есть $\frak{L}(G_0)= \frak{L}({G_0}')$.
\end{example}
\end{frame}

\subsubsection{1.1.3. Перечислимость языков}
\label{1.1.3.}
\begin{frame}
\frametitle{1.1.3. Перечислимость языков (1/3)}
Язык, порождаемый всякой формальной грамматикой, 
является перечислимым множеством в смысле теории алгоритмов.
Действительно, следующий алгоритм перечисляет
множество слов языка $\frak{L}(G)$, порождаемого
грамматикой $G=\langle \mathrm{M}, \Sigma, \Pi, S\rangle$.

\medskip
\TODO{ Первая часть эссе 1.1.3. 
Алгоритм на псевдокоде, с обоснованием. \\
Пример: двоичные последовательности, представляющие натуральные числа} 
\end{frame}

\begin{frame}
\frametitle{1.1.3. Перечислимость языков (2/3)}
\medskip
Однако не все языки, как множества слов,
являются перечислимыми.
Существуют неперечислимые языки, просто потому,
что множество алгоритмов счётно, а множество языков несчётно.

\medskip
Если никаких ограничений на грамматики не накладывать, то
формальная порождающая грамматика оказывается предельно общим
механизмом описания языков, что видно из следующей теоремы.
\begin{theorem}
Для любого перечислимого множества слов $M$ существует такая
формальная порождающая грамматика $G$, что $M=\frak{L}(G)$.
\end{theorem}
\begin{proof}

\medskip
\TODO{ Вторая часть эссе 1.1.3. Есть доказательство на основе машины Тьюринга, но в грамматике 14 правил, очень тяжеловесно. 
Сократили до 11, файл прислан.
Есть надежда на сокращение за счёт применения нормальных
алгорифмов Маркова.
}
\end{proof}
\end{frame}

\begin{frame}
\frametitle{1.1.3. Перечислимость языков (3/3)}
\medskip
Столь высокая степень общности ведёт к тому, что многие
практически важные вопросы для грамматик общего вида оказываются
очень трудоёмкими или даже алгоритмически неразрешимыми. Поэтому
на практике применяются более узкие классы грамматик. 

\medskip
Общепринятая в настоящее время
\odmkey{классификация грамматик}{Классификация!грамматик} была
предложена Хомским\footnote{Ноам Хомский, р. 1928}.
\end{frame}


% Текст слайда подготовил Ляпустин Егор 5030102/20201 
% Дата 10.01.26
\begin{frame}
    \frametitle{Перечислимость языков (1/7)}
    Язык, порождаемый всякой формальной грамматикой, является перечислимым множеством в смысле теории алгоритмов. Действительно, следующий алгоритм перечисляет множество слов языка $\mathcal{L}(G)$, порождаемого грамматикой $G = \langle N, T, R, S \rangle$.
    \begin{columns}[T]
        \begin{column}{6cm}
            \begin{algorithmic}
                    \REQUIRE Грамматика $G = \langle N, T, R, S \rangle$
                    \ENSURE Множество всех терминальных строк, выводимых из $S$
                    \STATE $Q \gets$ пустая очередь
                    \STATE $L \gets$ пустое множество
                    \STATE $\text{PUSH}(Q, S)$
                \end{algorithmic}
        \end{column}
        \begin{column}{9cm}
          \begin{algorithmic}
            \WHILE{$Q \neq \emptyset$}
            \STATE $\alpha \gets \text{POP}(Q)$ \COMMENT{извлечь первый элемент}
           \IF{$\forall c \in \alpha: c \in T$}
                    \STATE $L \gets L \cup \{\alpha\}$
                    \ELSE
                    \STATE $A_1, A_2, \dots, A_k \text{ -- все нетерминалы в } \alpha$
                    \FOR{$i \gets 1$ \TO $k$}
                    \STATE $A \gets A_i$
                    \FOR{$\forall (A \to \beta) \in R$}
%                    \STATE $\gamma \gets \text{строка, полученная заменой одного вхождения $A$ на $\beta$ в $\alpha$}$

                    \STATE $\text{PUSH}(Q, \gamma)$
                    \ENDFOR
                    \ENDFOR
                    \ENDIF
                    \ENDWHILE

                    \RETURN $L$
                \end{algorithmic}
       \end{column}
    \end{columns}
\end{frame}

% Текст слайда подготовил Ляпустин Егор 5030102/20201 
% Дата 10.01.26
\begin{frame}
    \frametitle{Перечислимость языков (2/7)}

    \begin{ground}
        Используется поиск в ширину для перебора всех возможных выводов из начального символа $S$. Это гарантирует, что все слова, которые могут быть получены с помощью правил грамматики $R$, будут перечислены.

        Поскольку алгоритм при каждом шаге применяет только правила из $R$ к нетерминальным символам, каждая полученная строка соответствует допустимому выводу в грамматике $G$.
    \end{ground}
    \begin{example}

        Двоичные последовательности, представляющие десятичные числа. Количество символов $1$ в цепочке есть представляемое этой цепочкой натуральное число.

        $N = \{S\}$

        $T = \{1\}$

        $R = \{S \rightarrow 1, S \rightarrow 1S\}$
    \end{example}

\end{frame}

% Текст слайда подготовил Ляпустин Егор 5030102/20201 
% Дата 10.01.26
\begin{frame}
    \frametitle{Перечислимость языков (3/7)}

    \begin{example} (Продолжение)
        \begin{columns}[T]
            \begin{column}{0.48\textwidth}

                \textbf{Первая итерация:}
                \begin{itemize}

                    \item $\alpha = S$
                    \item 1S и 1 добавляется в $Q$
                \end{itemize}
            \end{column}
            \begin{column}{0.48\textwidth}

                \textbf{Вторая итерация:}

                \begin{itemize}
                    \item $\alpha = \text{1}S$
                    \item $\alpha$ не состоит только из терминалов
                    \item 11S, 11 добавляется в $Q$

                    \item $\alpha = \text{1}$
                    \item $\alpha$ состоит только из терминалов, добавляем в $L$.
                \end{itemize}
            \end{column}
        \end{columns}
    \end{example}
    Однако не все языки, как множества слов, являются перечислимыми. Существуют неперечислимые языки, просто потому, что множество алгоритмов счётно, а множество языков несчётно. Если никаких ограничений на грамматики не накладывать, то формальная порождающая грамматика оказывается предельно общим механизмом описания языков, что видно из следующей теоремы.
\end{frame}

% Текст слайда подготовил Ляпустин Егор 5030102/20201 
% Дата 10.01.26
\begin{frame}
    \frametitle{Перечислимость языков (4/7)}

    \begin{theorem}
        Для любого перечислимого множества слов $M$ существует такая формальная порождающая грамматика $G$, что $M = \mathcal{L}(G)$
    \end{theorem}
    \begin{proof}
        Рассмотрим правила грамматики, с помощью которой будем доказывать теорему.

        Пусть $T$ -- конечный алфавит терминалов, $\mathcal{A}$ -- нормальный алгорифм Маркова (НАМ) перечисляющий $M \subseteq T^*$, $ \Gamma \supseteq T $ -- aлфавит НАМ (включает терминалы и служебные символы), $ \omega_0 \in \Gamma^* $ -- начальное состояние НАМ ($\omega_0 = \varepsilon$), $ P = \{ r_1, r_2, \dots, r_n \} $ -- правила НАМ, где каждое $ r_i $ имеет вид $ u_i \to v_i $ ($ u_i, v_i \in \Gamma^* $)


        Пусть $ G = \langle N, T, R, S \rangle$:
        \begin{itemize}
            \item Терминалы: $ T $
            \item Нетерминалы: $ N = \{ S, A, B, C \} \cup \Gamma_{\text{nt}}$, где $ \Gamma_{\text{nt}} = \{ \underline{a} \mid a \in \Gamma \setminus T \} $ -- нетерминальные символы, соответствующие служебным символам НАМ.
            \item Аксиома грамматики: $ S $.
        \end{itemize}

    \end{proof}
\end{frame}
% Текст слайда подготовил Ляпустин Егор 5030102/20201 
% Дата 10.01.26
\begin{frame}
    \frametitle{Перечислимость языков (5/7)}

    \begin{proof} (Продолжение)

        Правила $ R $:
        \begin{enumerate}
            \small

            \item Инициализация: $S \to A \omega_0 B$

                  Здесь $ \omega_0 $ записывается как цепочка символов из $ \Gamma $
            \item Преобразование для моделирования правил НАМ: $ A \to A C$

                  Символ $ C $ будет ходить по слову и применять правила НАМ.
            \item Моделирование правил НАМ:

                  $ \forall  (u_i \to v_i)  \in P $, где $ u_i, v_i \in \Gamma^* $:
                  $$
                      C u_i \to v_i C
                  $$
                  Это позволяет $ C $ заменять вхождение $ u_i $ на $ v_i $ при движении слева направо.
            \item Проход $ C $ через терминалы:

                  $\forall a \in T $:
                  \[
                      C a \to a C
                  \]
            \item Удаление $ C $ при достижении конца слова: $C B \to B$

            \item Завершение вывода: $B \to \varepsilon, \quad A \to \varepsilon$
        \end{enumerate}

    \end{proof}
\end{frame}

% Текст слайда подготовил Ляпустин Егор 5030102/20201 
% Дата 10.01.26
\begin{frame}
    \frametitle{Перечислимость языков (6/7)}

    \begin{proof} (Продолжение)
        Вывод в $ G $ имитирует работу НАМ:
        Начинаем с применения правила (1) $ S \Rightarrow A \omega_0 B $. Правило (2) $ A \to A C $ может применяться многократно, генерируя несколько копий $ C $.

        Каждая копия $ C $ движется вправо (правило (4) $ C X \to Y C $), применяя одно из правил НАМ при встрече с $ u_i $, либо проскальзывая через терминалы.

        Когда $ C $ достигает $ B $, он исчезает: (5) $ C B \to B $

        Если после всех замен получено слово $ \omega \in T^* $, то применется правило (6) $ B \to \varepsilon $ и $ A \to \varepsilon $, оставляя только $ \omega $
        \begin{itemize}
            \item Если $ \omega \in M $:

                  Существует последовательность применений правил НАМ, переводящая $ \omega_0 $ в $ \omega $. В грамматике можно промоделировать ту же последовательность, выводя $ \omega $.
            \item Если $ \omega \in \mathcal{L}(G) $:

                  Тогда вывод в $ G $ соответствует некоторой корректной работе НАМ, значит, $ \omega $ порождается НАМ, т.е. $ \omega \in M $.

        \end{itemize}




        Таким образом, $ \mathcal{L}(G) = M $.

    \end{proof}


\end{frame}

% Текст слайда подготовил Ляпустин Егор 5030102/20201 
% Дата 10.01.26
\begin{frame}
    \frametitle{Перечислимость языков (7/7)}

    Столь высокая степень общности ведёт к тому, что многие практически важные вопросы для грамматик общего вида оказываются очень трудоёмкими или даже алгоритмически неразрешимыми. Поэтому на практике применяются более узкие классы грамматик.
    
    \bigskip
    \TODO{Идейно это именно то, что нужно.    
    Особенно нормальный алгориФм Маркова.
    Все обозначения неправильные, программы написаны ужасно,
    пример слишком слабый.
    В переработку.}
    

    \bigskip
    \TODO{Замечания ФАН 21 января. \\    
    По алгоритмам: присваивание $\Assign$, а стрелочки как раз push и pop. Пустое множество $\emptyset$.\\
    По обозначениям: кто такой $\Gamma$, обозначения для грамматик и языков другие. \\ 
    По примерам: сделана <<единиричная>> система, а требуется добавить двоичную.
    Вот она: $S\to 1\mid S (0\mid 1)$. Протокол выполнения алгоритма нужен,
    но это таблица = что на каждом шаге находится в каждой переменной. \\  
    По структуре: новые слайды должны быть вшиты в мои старые слайды и заканчиваться
    переходной фразой про классификацию. Плотность слайдов такая же, как везде.}

    
\end{frame}


\subsubsection{1.1.4. Классификация грамматик и языков}
\label{1.1.4.}


\begin{frame}
\frametitle{1.1.4. Классификация грамматик и языков (1/5)}

\begin{enumerate}
\item Грамматика \odmkey{типа 0}{Грамматика!типа 0}
(или грамматика \odmkey{с фразовой структурой}{Грамматика!с
фразовой структурой})~--- это грамматика общего вида без
ограничений на правила вывода. Класс таких грамматик обозначим
$\frak{T}_0$.
\item Грамматика \odmkey{типа 1}{Грамматика!типа 1}
(или грамматика \odmkey{непосредственных
составляющих}{Грамматика!непосредственных составляющих} или
\odmkey{контекст\-но-зависимая}{Грамматика!контекстно-зависимая})
~--- это грамматика, все правила которой имеют вид $\alpha
A\beta\rightarrow\alpha\gamma\beta$, где $\alpha,\beta, \gamma\in (\mathrm{M}\cup
\Sigma)^*$, $|\gamma|\geq 1$, $A\in \mathrm{M}$. Такое правило описывает замену
нетерминального символа $A$ на цепочку $\gamma$ в определенном
\odmkey{контексте}{Контекст!правила} ($\alpha,\beta$), который не
меняется при применении правила. Класс таких грамматик обозначим
$\frak{T}_1$.
\item Грамматика \odmkey{типа 2}{Грамматика!типа 2}
(или
\odmkey{контекстно-свободная}{Грамматика!контекстно-свободная}
грамматика)~---  это грамматика, все правила которой имеют вид
$A\rightarrow\alpha$, где $\alpha\in (\mathrm{M}\cup \Sigma)^*$, $A\in \mathrm{M}$. Такое
правило описывает замену нетерминального символа $A$ на цепочку
$\alpha$ вне зависимости от контекста. Класс таких грамматик
обозначим $\frak{T}_2$.
\item Грамматика \odmkey{типа 3}{Грамматика!типа 3}
(или \odmkey{регулярная}{Грамматика!регулярная} грамматика) ~---
это грамматика, все правила которой имеют вид либо
$A\rightarrow\alpha$, либо $A\rightarrow\alpha B$ (или
$A\rightarrow B\alpha$), где $\alpha\in \Sigma^*$,\\ $A,B\in \mathrm{M}$.
В правой части правила регулярной грамматики может быть не более
одного нетерминала.
Класс таких грамматик обозначим
$\frak{T}_3$.
\end{enumerate}
\end{frame}

\begin{frame}
\frametitle{1.1.4. Классификация грамматик и языков (2/4)}
Ясно, что $\A{i\in 1..3}{G\in \frak{T}_i\Impl G\in \frak{T}_{i-1}}$, поскольку всякое правило
грамматики класса $\frak{T}_i$ 
является частным случаем 
правила грамматики класса $\frak{T}_{i-1}$.

\begin{remark}
Исключением являются правила с пустой правой частью,
но их можно устранить, см. \DMref{п.~1.3.2}{1.3.2.}.
\end{remark}
\medskip
Один и тот же язык может быть описан разными грамматиками, в том
числе грамматиками разного типа. Принято относить язык к
{\em наибольшему} типу из типов описывающих его грамматик. То есть если
для языка $L$ есть две описывающие его грамматики $G'$ и $G''$,
$L=\frak{L}(G')\And L=\frak{L}(G'')\And G'\in \frak{T}_i\And G''\in \frak{T}_j\And i<j$, то
язык $L$ относится к типу $\frak{T}_j$.
\begin{example}
Рассмотрим грамматику
$G=\langle\{S,N,D\},\{0,1,2,3,4,5,6,7,8,9,-,+\},
\linebreak
\{S\rightarrow
N, S\rightarrow
+N, S\rightarrow
-N, N\rightarrow D, N\rightarrow ND, D\rightarrow 0, D\rightarrow 1,\ldots,D\rightarrow 9 \}, S\rangle$, описывающую множество десятичных целых чисел со знаком.
Эта грамматика относится к типу 2 (правило $N\rightarrow ND$
недопустимо в грамматике типа 3). Однако для того же самого языка
можно предложить грамматику
$G'=\langle\{S,N\},\{0,1,2,3,4,5,6,7,8,9,-,+\},
\linebreak
\{S\rightarrow N, 
S\rightarrow +N, 
S\rightarrow -N, 
N\rightarrow 0,
N\rightarrow 1,
\ldots,
N\rightarrow 9,
N\rightarrow 0N,
N\rightarrow 1N,
\ldots,\linebreak
N\rightarrow 9N
\}, S \rangle$,
которая является грамматикой типа 3. Таким образом, язык целых
десятичных чисел со знаком является языком типа 3.
\end{example}

\end{frame}

\begin{frame}
\frametitle{1.1.4. Классификация грамматик и языков (3/4)}

Вопрос об определении типа конкретной грамматики решается очень
просто~--- рассмотрением вида правил этой грамматики. К сожалению,
вопрос об определении типа языка оказывается алгоритмически
неразрешимым.
\begin{theorem} 
Следующая проблема является алгоритмически неразрешимой: дано
$L=\frak{L}(G)$, $G\in \frak{T}_i$, определить, существует ли $G'$,
такая что $M=\frak{L}(G') \And
G'\in \frak{T}_j\And i<j $.
\end{theorem}
\begin{proof} \TODO{Теорема Райса-Хиггинсона. 
Доказательство обязательно в обозначениях раздела 4.5 из Дискретной математики для программистов. Целесообразно добавить атом про проблему останова во Введении. 
Студенты, которые слушали курс Пастора по теории алгоритмов, могут помочь.}
\end{proof}
\end{frame}

\begin{frame}
\frametitle{1.1.4. Классификация грамматик и языков (4/4)}

\begin{digress}
Чем выше тип грамматики, 
порождающий язык, тем лучше ---
язык оказывается проще,
его легче анализировать,
описывать и применять.
Поэтому авторы новых
формальных языков стараются
использовать такие грамматические конструкции,
которые выражаются правилами
регулярных грамматик,
в крайнем случае контекстно-свободных грамматик.
Однако иногда приходится иметь дело
с унаследованными языками, 
не только естественными,
но и искусственными, в частности,
с ранними языками программирования.
В таких случаях ради совместимости
приходится жертвовать удобством.  
\end{digress}

\begin{example}
В языке Algol-60 была 
впервые официально использована теория 
формальных грамматик для описания стандарта языка.
При этом принятая  грамматика не позволяла 
определить, к которому $\textbf{if}$ 
относится $\textbf{else}$ 
в условном операторе $\textbf{if}
B_1 \textbf{then if} B_2 \textbf{then}
S_1 \textbf{else} S_2$. 
Язык оказался существенно 
синтаксически неоднозначным.
\end{example}

\end{frame}

% Текст слайда подготовил Ляпустин Егор 5030102/20201 
% Дата 10.01.26
\begin{frame}
    \frametitle{Классификация грамматик и языков (4/5)}
    Вопрос об определении типа конкретной грамматики решается очень просто — рассмотрением вида правил этой грамматики. К сожалению, вопрос об определении типа языка оказывается алгоритмически неразрешимым.
    \begin{theorem}
        Следующая проблема является алгоритмически неразрешимой: дано $M = \mathcal{L}(G)$, $G \in T_i$, определить, существует ли $G'$, такая что $M = \mathcal{L}(G') \text{ $\&$ } G' \in T_j \text{ $\&$ } i < j$
    \end{theorem}
    \begin{proof} От противного. Предположим, что существует алгоритм $\mathcal{A}$, который для заданного языка $M = \mathcal{L}(G)$ и грамматики $G \in T_i$ определяет, существует ли грамматика $G' \in T_j$, $i < j$, такая что $\mathcal{L}(G') = M$. Поскольку M рекурсивно перечислим, существует грамматика $G \in T_0$, такая что $\mathcal{L}(G) = M$.

        Пусть T – произвольная машина Тьюринга, $\omega$ – вход для этой машины.
        Построим язык М следующим образом:
        \[
            M =
            \begin{cases}
                \{\omega\},  & \text{если } T \text{ останавливается на входе } \omega; \\
                \varnothing, & \text{иначе.}
            \end{cases}
        \]
        Заметим, что язык М рекурсивно перечислимый, так как мы можем симулировать работу Т на $\omega$.
    \end{proof}
\end{frame}

% Текст слайда подготовил Ляпустин Егор 5030102/20201 
% Дата 10.01.26
\begin{frame}
    \frametitle{Классификация грамматик и языков (5/5)}
    \begin{proof} (Продолжение)
        Применим алгоритм из предположения $A$ к языку $M$ и грамматике $G \in T_0$, которая порождает $M$. Алгоритм $A$ должен определить, существует ли контекстно-свободная грамматика $G' \in T_2$ (поскольку $T_2$ имеет более низкий тип по сравнению с $T_0$), такая что $\mathcal{L}(G') = M$.

        Если алгоритм $A$ определит, что такая грамматика $G'$ существует, это будет означать, что $M$ является контекстно-свободным языком. Однако, поскольку $M$ построен на основе поведения машины Тьюринга $T$ на входе $\omega$, возможность определить, является ли $M$ контекстно-свободным, эквивалентна решению проблемы остановки для $T$ на $\omega$. Так как проблема остановки алгоритмически неразрешима, это приводит к противоречию с предположением о существовании алгоритма $A$.
    \end{proof}
    
    \bigskip
    \TODO{Неплохо. Исправить обозначения и вклеить в существующие слайды.
    Очень важно иметь во введении неразрещимость проблемы остановки. Также пока не было сказано, что такое симуляция работы машины Тьюринга.}
\end{frame}

\subsubsection{1.1.5. Неукорачивающие грамматики}
\label{1.1.5.}
\begin{frame}
\frametitle{1.1.5. Неукорачивающие грамматики (1/5)}

Грамматика $G=\langle \mathrm{M}, \Sigma, \Pi, S\rangle$ называется
\odmkey{неукорачивающей}{Грамматика!неукорачивающая}, если
длина левой части любого правила не превосходит
длины правой части:
$\A{r=\alpha\rightarrow\beta\in \Pi}{|\alpha|\leq|\beta|}$. При
любом выводе в такой грамматике длины цепочек могут разве что
возрастать. Поэтому, в частности, пустое слово $\varepsilon$ не
выводимо в такой грамматике.
\begin{example}
Контекстно-зависимые грамматики являются неукорачивающими.
\end{example}

Язык $L\subset \Sigma^*$ называется
\odmkey{разрешимым}{Язык!разрешимый}
(по аналогии с понятием {\em разрешимое множество}), если существует
алгоритм, который для произвольной цепочки $\alpha\in \Sigma^*$ за
конечное число шагов даёт ответ на вопрос, входит ли эта цепочка в
язык, т.~е. $\alpha\in L$ или $\alpha\not\in L$. Если при этом число шагов алгоритма определяется по длине цепочки
$\alpha$ и может быть указано заранее, то язык называется
\odmkey{легко разрешимым}{Язык!легко разрешимый}.
Если же алгоритм даёт ответ, когда \\$\alpha \in L$, и может не давать ответа (алгоритм не заканчивает работу),
когда $\alpha\not\in L$, то язык называется
\odmkey{полуразрешимым}{Язык!полуразрешимый}.
Алгоритм, разрешающий язык, называют  \odmkey{распознающим}{Язык!распознаваемый}. 

\begin{theorem}
Язык $\frak{L}(G)$, порождаемый неукорачивающей грамматикой
 $G=\langle \mathrm{M}, \Sigma, \Pi, S\rangle$, \\легко разрешимый.
\end{theorem}
\begin{proof}
Если $\alpha\in \frak{L}(G)$, то $S\Rightarrow^*\alpha$, существует бесповторный вывод $\Tuple{\sigma}{0}{k}$,
где $\sigma_0 = S$ и
$\sigma_k=\alpha$, причём 
 $\A{i\in 0..k}{|\sigma_i|\leq|\alpha|}$. Обозначим $n\Assign|\alpha|$.
%Если же \\ $\alpha\not\in \frak{L}(G)$, то $\Not
%S\Rightarrow^*\alpha$.
\end{proof}
\end{frame}

\begin{frame}
\frametitle{1.1.5. Неукорачивающие грамматики (2/5)}

\begin{proof} (Окончание)
Число $r(n)$  различных цепочек длиной не более
$n$:

\vspace{-12pt}
$$
r(n) = \sum_{i=1}^{n} \left|\mathrm{M}\cup \Sigma\right|^i
=\sum_{i=1}^n (|\mathrm{M}|+|\Sigma|)^i=
(|\mathrm{M}|+|\Sigma|)\frac{(|\mathrm{M}|+|\Sigma|)^n -1}{|\mathrm{M}|+|\Sigma| -1}.
$$

\vspace{-4pt}
Далее, длина бесповторной последовательности, составленной из
таких цепочек, также ограничена (например, величиной $r(n)$). Общее
количество $R(n)$ всех бесповторных последовательностей цепочек длиной не
более $n$ также ограничено:

\vspace{-16pt}
 $$R(n)=\sum_{i=0}^{r(n)} A(r(n),i) =
 \sum_{i=0}^{r(n)} \frac{r(n)!}{(r(n)-i)!}=
  r(n)!\sum_{i=0}^{r(n)}\frac{1}{i!}
  \leq r(n)!\sum_{i=0}^\infty\frac{1}{i!}
  =r(n)!\cdot e.$$ 

\vspace{-4pt}
Тем
более ограничено количество таких бесповторных
последовательностей, которые начинаются с $S$ и заканчиваются
$\alpha$.
Отсюда немедленно вытекает схема алгоритма распознавания. Нужно
перебрать все бесповторные последовательности цепочек длиной не
более $n$, которые начинаются с $S$ и заканчиваются $\alpha$, и для
каждой последовательности проверить, не является ли она допустимым
выводом $\alpha$ в $G$. 
%(то есть проверить, что каждый следующий
%элемент последовательности непосредственно выводим из
%предыдущего). 
Если такой вывод найдётся, то можно прекратить
дальнейший перебор и сделать вывод, что $\alpha\in \frak{L}(G)$. В
противном случае $\alpha\not\in \frak{L}(G)$.
\end{proof}
\end{frame}

\begin{frame}
\frametitle{1.1.5. Неукорачивающие грамматики (3/5)}
\begin{remark}
Число шагов перебора,
о котором идёт речь в доказательстве,
очень велико.
Это не случайно.
Можно показать, что в общем случае не существует эффективного
алгоритма распознавания языка, порождённого неукорачивающей
грамматикой. Другими словами, хотя число шагов работы алгоритма
распознавания можно предсказать заранее, но это число
может быть столь велико, что алгоритм практически не применим.
\end{remark}


\begin{digress}
Грамматически можно
описать любое перечислимое множество, в том числе и неразрешимое.
Поскольку грамматики типа $\frak{T}_1$ описывают распознаваемые, то есть
разрешимые языки, ясно что описание нераспознаваемых языков
является уделом грамматик типа $\frak{T}_0$. Хотя грамматики типа $\frak{T}_1$
удобны для описания языков (контекст очень полезен)
и порождают распознаваемые языки,
но алгоритмы распознавания имеют слишком
большую трудоёмкость. Поэтому грамматики
типа $\frak{T}_0$ и $\frak{T}_1$ {\em никогда} не применяются для описания
искусственных языков программирования.
В информационных технологиях
применяются грамматики типа $\frak{T}_2$ и $\frak{T}_3$.
Лингвисты иногда применяют грамматики
типа $\frak{T}_0$ и $\frak{T}_1$ для описания естественных языков,
но это только потому, что они не имеют права
изменять язык в угоду удобству описания,
а программисты могут и должны это делать.
Современная тенденция в языкотворчестве
ведет к исключению из формальных языков конструкций,
которые невозможно описать грамматиками типа $\frak{T}_2$ и $\frak{T}_3$.
\end{digress}

\end{frame}

\begin{frame}
\frametitle{1.1.5. Неукорачивающие грамматики (4/5)}

\begin{remark}
Любую грамматику можно доопределить так, 
чтобы все символы, входящие в исходные правила, были нетерминалами.
Для этого достаточно
для каждого терминального символа
$a$ ввести новый нетерминальный символ
$A$ и правило $A\to a$. После этого
нужно все вхождения терминальных
символов
в старые правила заменить 
новыми нетерминальными символами.
\end{remark}



\begin{theorem}
Для всякой неукорачивающей грамматики существует эквивалентная ей \\контекстно-зависимая грамматика.
\end{theorem}

\begin{proof}
Построим правила для контекстно-зависимой
грамматики $G_2$ на основе правил
неукорачивающей грамматики $G_1$.
Каждое правило 
$$X_1 X_2 \ldots X_n \to Y_1 Y_2 \ldots Y_{n+m} \in \Pi,\quad n > 0, m \geq 0$$
заменим на ряд из $2n$ следующих правил:
$X_1 X_2 \dots X_n \rightarrow Z_1 X_2 \dots X_n$,
\\
$Z_1 X_2 \dots X_n \rightarrow Z_1 Z_2 \dots X_n$,
$\ \dots$,
$\ Z_1 Z_2 \dots X_n \rightarrow Z_1 Z_2 \dots Z_n$,
и \\
$Z_1 Z_2 \dots Z_n \rightarrow Y_1 Z_2 \dots Z_n$, $\ Y_1 Z_2 \dots Z_n \rightarrow Y_1 Y_2 \dots Z_n$,
$\ \dots$,\\
$Y_1 Y_2 \dots Z_n \rightarrow Y_1 Y_2 \dots Y_n \dots Y_{n+m}$,
где $Z_1, Z_2, \dots, Z_n$~--- новые нетерминалы, 
для каждого исходного правила они свои.
\end{proof}
\end{frame}

\begin{frame}
\frametitle{1.1.5. Неукорачивающие грамматики (5/5)}
\begin{proof} (продолжение)
Правила с новыми нетерминалами могут применяться только в одной указанной последовательности. 
Другими словами, если будет применено правило 
$X_1 X_2 \dots X_n \rightarrow Z_1 X_2 \dots X_n$, то все последующие правила тоже будут применены, потому что иначе нетерминалы $Z_1, Z_2, ... Z_n$ будут присутствовать в выведенном слове, чего быть не должно.
Построенная грамматика эквивалентна исходной, потому что в результате применения новых правил 
цепочка $X_1 X_2 ... X_n$ перейдет в $Y_1 Y_2 ... Y_{n+m}$. 
Новая грамматика является контекстно-зависимой грамматикой по построению.
\end{proof}

\medskip
Из доказанной теоремы и предшествующих наблюдений
следует важный вывод:
классы языков, порождаемых неукорачивающими грамматиками и грамматиками непосредственных составляющих, совпадают. 

\begin{remark}
Неукорачивающие грамматики,
очевидно, не могут порождать
языки, содержашие пустое слово $\varepsilon$.
В большинстве случаев это обстоятельство
не является существенным.
Если, однако, требуется,
чтобы пустое слово входило
в язык, то достаточно
в неукорачивающей 
грамматике разрешить иметь
одно исключительное правило:
$S\to\varepsilon$. 
\end{remark}

\end{frame}

\subsubsection{1.1.6. Контекстно-свободные грамматики}\label{1.1.6.}
\begin{frame}
\frametitle{1.1.6. Контекстно-свободные грамматики (1/5)}
\begin{remark}
Для краткости оборот <<контекстно-свободная грамматика>> далее заменяется аббревиатурой КСГ. 
\end{remark}
Пусть задана КСГ $G=\langle \mathrm{M}, \Sigma, \Pi, S\rangle$
и вывод $S\Rightarrow_G^*\alpha$
в этой грамматике. Поскольку все правила
имеют вид $A\rightarrow \alpha$, то в выводе
$S=\sigma_0,\sigma_1,\dots,\sigma_n=\alpha$
для каждой пары $\sigma_i,\sigma_{i+1}$
соседних цепочек
следующая получается из предыдущей заменой
{\em одного} нетерминального символа.
Вообще говоря, возможно заменить
любой нетерминальный символ,
поэтому последовательность цепочек в выводе 
отнюдь не является однозначно определённой.
Вывод называется
\odmkey{левосторонним}{Вывод!левосторонний}
(\odmkey{правосторонним}{Вывод!правосторонний}),
если на каждом шаге заменяется самое левое
(соответственно, самое правое)
вхождение нетерминального символа в цепочку.

Вывод в КСГ
удобно представить в виде
\odmkey{дерева разбора}{Дерево!разбора}
(или \odmkey{дерева вывода}{Дерево!вывода})~---
упорядоченного ордерева,
корнем которого является аксиома грамматики,
нелистовыми узлами являются нетерминалы,
а листьями~--- терминалы. Если в выводе было применено
правило  $A\rightarrow \alpha$, то в дереве разбора проводятся
дуги из узла $A$ ко всем узлам, соответствующим буквам
цепочки $\alpha$, причем эти узлы располагаются
в том же порядке, что и буквы в цепочке $\alpha$.

\end{frame}


\begin{frame}
\frametitle{1.1.6. Контекстно-свободные грамматики (2/5)}

\begin{columns}[t]
\begin{column}{10.1cm}
Дерево вывода компактно
представляет множество эквивалентных выводов,
в которых применены одни и те же правила, но в разном порядке.
\begin{example}
Рассмотрим КСГ $G_1$ с правилами
$E\rightarrow E+T$,
$E\rightarrow T$,
$T\rightarrow T*F$,
$T\rightarrow F$,
$F\rightarrow a$~--- упрощенную грамматику арифметических выражений.
Дерево разбора справа в компактной форме представляет 10 (десять!) возможных выводов
цепочки $a+a$.
\end{example}

Ясно, что по дереву вывода можно однозначно
восстановить левосторонний (равно как и правосторонний)
вывод в данной грамматике.

\end{column}
\begin{column}{5cm}
\vspace{-20pt}
\begin{figure}[h]
\begin{center}
\includegraphics[height=50mm]{1_1_4_1}
\end{center}
\end{figure}
\end{column}
\end{columns}

\vspace{4pt}
Ещё одной наглядной формой 
представления выводов в КСГ является 
\odmkey{дерево порождения}{Дерево!порождения}.
Дерево порождения строится следующим образом.
Корнем дерева является узел, помеченный аксиомой грамматики 
$S$.
Если в грамматике есть правило $A\to\alpha$,
и в дереве построен узел, 
помеченный цепочкой $\beta A \gamma$,
то добавляется узел, помеченный цепочкой
$\beta\alpha\gamma$, и проводится дуга 
$(\beta A \gamma,\beta\alpha\gamma)$,
помеченная правилом $A\to\alpha$.

\end{frame}


\begin{frame}
\frametitle{1.1.6.Контекстно-свободные грамматики (3/5)}

\begin{example}
На рисунке представлено дерево порождения,
в котором все десять выводов предыдущего примера
показаны в явной форме. 
\end{example}
\begin{figure}[h]
\begin{center}
\includegraphics[width=150mm]{1_1_4_2}
\end{center}
\end{figure}

\end{frame}


\begin{frame}
\frametitle{1.1.6.Контекстно-свободные грамматики (4/5)}
Один и тот же язык может порождаться
разными (но эквивалентными) контекстно-свобод\-ными грамматиками.


\begin{example}
Рассмотрим ещё раз КСГ $G_1$ с правилами
$E\rightarrow E+T$,
$E\rightarrow T$,
$T\rightarrow T*F$,
$T\rightarrow F$,
$F\rightarrow a$ и КСГ
$G_2$ с правилами
$E\rightarrow E+E$,
$E\rightarrow E*E$,
$E\rightarrow a$.
Грамматика $G_2$  задает тот же
упрощенный язык арифметических выражений,
что и грамматика $G_1$.
Чтобы убедиться в этом, достаточно заметить,
что любой вывод в грамматике $G_1$ 
легко перестраивается в вывод 
в грамматике $G_2$ той же терминальной 
цепочки.
Действительно, вместо правила
$E\rightarrow E+T$ надо применить правило
 $E\rightarrow E+E$,
вместо правила $T\rightarrow T*F$
 надо применить правило
$E\rightarrow E*E$,
вместо правила $F\rightarrow a$ надо применить правило
 $E\rightarrow a$,
а применения правил 
$E\rightarrow T$ и 
$T\rightarrow F$ надо просто исключить из вывода.
\end{example}

\smallskip

Представленные грамматики $G_1$ и $G_2$
задают один и тот же язык,
но грамматика $G_2$ более компактна, то есть содержит меньше правил и меньше нетерминалов,
чем грамматика $G_1$. 
Возникает вопрос: на основании каких свойств
следует предпочесть одну КСГ другой? 
Ответ на поставленный вопрос зависит
от целей построения грамматического описания языка.

\smallskip
Например, для {\em публикации}
языка предпочтительна
более компактная грамматика $G_2$.

\end{frame}

\begin{frame}
\frametitle{1.1.6.Контекстно-свободные грамматики (5/5)}
Однако с точки зрения 
использования языка для {\em программирования }
 грамматика $G_2$ обладает двумя существенными
недостатками.

\smallskip
Во-первых,
КСГ $G_2$ \odmkey{неоднозначна}{Грамматика!неоднозначная}~---
на рисунке представлены два различных дерева разбора
для выражения $a+a*a$ в этой грамматике.

\begin{figure}[h]
\begin{center}
\includegraphics[width=110mm]{1_1_5.jpg}
\end{center}
\end{figure}


Во-вторых, грамматика $G_2$ не учитывает семантику языка
(приоритет арифметических операций).
Если использовать дерево разбора
для вычислений
(что и делается в большинстве реальных компиляторов),
то дерево на рисунке справа
даст неверный результат.

\end{frame}





\subsubsection{1.1.7. Линейные, регулярные и автоматные грамматики}\label{1.1.7.}
\begin{frame}
\frametitle{1.1.7. Линейные, регулярные и автоматные грамматики (1/5)}

Пусть задана порождающая грамматика $G=\langle \mathrm{M}, \Sigma, \Pi, S\rangle$.
Правило грамматики называется
\odmkey{начальным}{Правило!начальное}, если 
в левой части стоит аксиома грамматики $S$.
Правило грамматики называется
\odmkey{заключительным}{Правило!заключительное}, если оно имеет
вид $A\rightarrow\alpha$, где $A\in \mathrm{M}$, $\alpha\in \Sigma^*$. 
Правило
называется \odmkey{линейным}{Правило!линейное}
(соответственно,
\odmkey{леволинейным}{Правило!леволинейное} или
 \odmkey{праволинейным}{Правило!праволинейное}),
если оно имеет вид $A\rightarrow \alpha B\beta$ (соответственно, $A\rightarrow B\beta$ или
$A\rightarrow \alpha B$), где $A, B \in \mathrm{M}$ и 
$\alpha, \beta\in \Sigma^*$.

\smallskip
Порождающая грамматика, все правила которой заключительные или
линейные (соответственно, леволинейные или праволинейные) называется
\odmkey{линейной}{Грамматика!линейная} (соответственно, 
\odmkey{леволинейной}{Грамматика!леволинейная} или
\odmkey{праволинейной}{Грамматика!праволинейная}). 
Очевидно,
что праволинейные и леволинейные грамматики
являются линейными, а линейные грамматики
являются контекстно-свободными.

\smallskip
Язык, порождаемый линейной (соответственно,
леволинейной или праволинейной) грамматикой
называется    
\odmkey{линейным}{Язык!линейный} 
(соответственно, 
\odmkey{леволинейным}{Язык!леволинейный} или
\odmkey{праволинейным}{Язык!праволинейный}). 
\begin{example}
Язык $\Set{a^n b^n}{n\in \Bbb{N}}$ 
является линейным языком, потому что порождается 
линейной грамматикой
$\langle \{S, A\}, \{a,b\},
\{S\rightarrow aAb, A\rightarrow 
\varepsilon, A\rightarrow aAb\},S\rangle$.
\end{example}

\begin{remark}
Любое дерево вывода в линейной грамматике
имеет один нетерминальный <<ствол>>,
по бокам от которого висят терминальные 
<<листья>>, причём в леволинейных деревьях листья только справа, а в праволинейных деревьях листья только слева.
\end{remark}

\end{frame}


\begin{frame}
\frametitle{1.1.7. Линейные, регулярные и автоматные грамматики (2/5)}
\begin{theorem}
Для любой праволинейной грамматики 
существует эквивалентная леволинейная грамматика
и наоборот.
\end{theorem}
\begin{proof}
Пусть задана праволинейная грамматика. 
Построим для неё эквивалентную 
леволинейную грамматику: 
для каждого начального правила 
$S\rightarrow \alpha A$ введём
заключительное правило 
$A\rightarrow \alpha$,
для каждого праволинейного правила 
$A\rightarrow \alpha B$ введём
леволинейное правило 
$B\rightarrow A\alpha$,
для каждого заключительного правила 
$A\rightarrow \alpha$ введём
начальное правило 
$S\rightarrow A\alpha $.
Теперь любой вывод 
в исходной грамматике 
преобразуется в вывод в построенной грамматике:
достаточно применять новые правила начиная с конца исходного вывода.
\end{proof}

Таким образом праволинейные и леволинейные грамматики порождают один и тот же класс языков.
Поэтому определение в \DMref{п.~1.1.3}{1.1.3.}, 
гласящее, что леволинейные и
праволинейные грамматики называются
\odmkey{регулярными}{Грамматика!регулярная},
оказывается вполне корректным. 
Для определённости далее рассматриваются леволинейные грамматики. 
Регулярная грамматика
называется \odmkey{автоматной}{Грамматика!автоматная}, если все
терминальные цепочки $\alpha$ в правилах состоят ровно из одного
символа. Другими словами, леволинейная автоматная грамматика
содержит правила вида $A\rightarrow a$ и $A\rightarrow B a$, где
$A, B \in \mathrm{M}$, $a\in \Sigma$.

\end{frame}


\begin{frame}
\frametitle{1.1.7. Линейные, регулярные и автоматные грамматики (3/5)}

Терминальные цепочки в правилах регулярной грамматики могут быть
пустыми, и значит, язык, порождаемый регулярной грамматикой $G$,
может содержать пустое слово $\varepsilon$.  В то же время язык,
порождаемый автоматной грамматикой, очевидно, не содержит пустого
слова $\varepsilon$. 
Таким образом, формально эти классы языков различны, но на самом деле они почти эквивалентны.

\smallskip
Чтобы установить эту связь,
необходимо обозначить некоторые специальные свойства
грамматик, которые понадобятся также и в дальнейшем.    

Нетерминальный символ называется
\odmkey{непорождающим}{Символ!непорождающий},
если из него невозможно вывести терминальную цепочку.
Нетерминальный или терминальный  символ называется
\odmkey{недостижимым}{Символ!недостижимый},
если его невозможно вывести из аксиомы.
Если правило имеет вид $A\rightarrow B$, то оно называется
\odmkey{цепным}{правило!цепное}. 
Если в правиле $A\rightarrow
\alpha$ цепочка $\alpha$ пустая,
$\alpha=\varepsilon$, то правило принимает вид 
$A\rightarrow$  и называется
\odmkey{$\varepsilon$-правилом}{Правило!$\varepsilon$}.
Если в одном из начальных правил $S\rightarrow \alpha$
аксиома входит в правую часть,
то правило называется 
\odmkey{$S$-правилом}{Правило!$S$}.

Недостижимые и непорождающие символы, цепные правила,
$\varepsilon$-правила и $S$-правила грамматик
создают различные помехи в рассуждениях и применении алгоритмов, но при этом все такие особенности грамматик являются {\em устранимыми} (см. \DMref{п.~1.3.2}{1.3.2.}).

\end{frame}


\begin{frame}
\frametitle{1.1.7. Линейные, регулярные и автоматные грамматики (4/5)}
Грамматики, в которых нет непорождающих 
и недостижимых символов, а также нет 
цепных правил, 
$\varepsilon$-правил и
$S$-правил называются \odmkey{приведёнными}{Грамматика!приведённая}.

\smallskip
На практике языки с пустыми словами, как правило, не используются,
а грамматики составляются квалифицированными авторами и
не содержат непорождающих 
и недостижимых символов, а также неподходящих правил.
Поэтому мы докажем утверждение 
об эквивалентности регулярных и автоматных грамматик
в существенно более
слабой, но практически более полезной форме, а именно в форме
алгоритма преобразования {\em приведённой}  грамматики к нужному виду.

\begin{theorem}
Для любой приведённой регулярной грамматики существует эквивалентная автоматная
грамматика.
\end{theorem}

\begin{proof}
Пусть задана леволинейная регулярная грамматика $G=\langle \mathrm{M}, \Sigma,
\Pi, S\rangle$, где каждое правило $r\in \Pi$ имеет вид $A\rightarrow
\alpha$ или $A\rightarrow B\alpha$, причем $A, B \in \mathrm{M}$,\\
$\alpha\in \Sigma^*$ и $|\alpha|\geq 1$. Построим автоматную грамматику
$G'=\langle \mathrm{M}', \Sigma, \Pi', S\rangle$  в соответствии с алгоритмом на следующем слайде.
 Очевидно, что всякому выводу в грамматике $G$ однозначно
 соответствует вывод в грамматике $G'$. Таким образом, грамматики
 эквивалентны.
\end{proof}
\end{frame}

\begin{frame}
\frametitle{1.1.7. Линейные, регулярные и автоматные грамматики (5/5)}


\begin{algorithmic}
\REQUIRE Регулярная леволинейная грамматика
$G=\langle \mathrm{M}, \Sigma, \Pi, S\rangle$
%, не содержащая цепных правил и $\varepsilon$-правил.
\ENSURE
Автоматная леволинейная грамматика
$G'=\langle \mathrm{M}', \Sigma, \Pi',S\rangle$.
\STATE{$\mathrm{M}'\Assign \mathrm{M}; \Pi'\Assign\emptyset$}
 \COMMENT { $r= A \rightarrow \alpha$ или $r= A\rightarrow B\alpha$ }
   \FOR {\textbf{$r\in \Pi$}}
  \STATE {\textbf{if $|\alpha|=1$ then 
  $\Pi'\Assign \Pi'\cup r;$ next for $r$ end if}}
    \COMMENT{ переносим без изменений }
      \STATE {\textbf{$n\Assign|\alpha|-1$}}
    \COMMENT{ количество новых нетерминалов }
    \STATE{\textbf{$\mathrm{M}'\Assign \mathrm{M}'\cup$ new $A_n$ }}
    \COMMENT{ добавляем нетерминал $A_n$ }
    \STATE {\textbf{$\Pi'\Assign \Pi'\cup(A\rightarrow A_n\alpha [n+1])$}}
    \COMMENT{ добавляем правило
     $A\rightarrow A_na_{n+1}$ }
    \FOR{\textbf{$i$ from $n-1$ to 1}}
       \STATE {\textbf{$\mathrm{M}'\Assign \mathrm{M}'\cup$ new $A_i$ }}
       \COMMENT{ добавляем нетерминал $A_i$ }
       \STATE{\textbf{$\Pi'\Assign \Pi'\cup (A_{i+1}\rightarrow A_i \alpha [i+1])$}}
       \COMMENT{ добавляем правило $A_{i+1}\rightarrow A_i a_{i+1}$ }
    \ENDFOR
    \STATE{\textbf{if $r=A\rightarrow \alpha$ then $\Pi'\Assign \Pi'\cup(A_1\rightarrow \alpha[1])$}}
      \COMMENT{ заключительное правило }
     \STATE {\textbf{else $\Pi'\Assign \Pi'\cup(A_1\rightarrow B\alpha[1])$ end if}}
      \COMMENT{ незаключительное правило }
 \ENDFOR
\end{algorithmic}

\end{frame}

\subsubsection{1.1.8. Языки записи грамматик}\label{1.1.8.}
\begin{frame}
\frametitle{1.1.8. Языки записи грамматик (1/9)}
Запись правил грамматики в общей форме довольно многословна, не всегда легко читается и
существенно использует рекурсию. 

\begin{example}
Набор правил
$N\to D$, $N\to ND$, $D\to 0$,
$D\to 1$, $D\to 2$, $D\to 3$, 
$D\to 4$, $D\to 5$, 
$D\to 6$, $D\to 7$, $D\to 8$, $D\to 9$
по существу означает, что целое без знака~--- это
последовательность из одной или более цифр, что не всегда очевидно
без дополнительных пояснений.
\end{example}

\smallskip
Были предложены многочисленные усовершенствования записи грамматик, в особенности контекстно-свободных и 
автоматных грамматик,
которые сводятся к введению дополнительных
\odmkey{метасимволов}{Метасимвол}. Дополнительные метасимволы не
меняют выразительной силы  грамматик, но делают запись грамматики более
удобной для восприятия человеком. Обычно рассматриваются следующие
метасимволы.

\begin{enumerate}
\item Вертикальная черта $\Ms{|}$~--- разделение альтернатив. Правило
$$
A\rightarrow \alpha \Ms{|}\beta
$$
эквивалентно паре правил
$$
A\rightarrow \alpha  \qquad A\rightarrow\beta
$$
\end{enumerate}
\end{frame}

\begin{frame}
\frametitle{1.1.8. Языки записи грамматик (2/9)}
\begin{enumerate}

\item[2.] Круглые скобки $\Ms{(\ )}$~--- введение неявного нетерминала.
Круглые скобки могут содержать несколько альтернатив, разделенных
метасимволом $\Ms{|}$. Правило
$$
A\rightarrow \alpha\Ms{(}\beta\Ms{)}\gamma
$$
эквивалентно паре правил
$$
A\rightarrow \alpha B\gamma \qquad B\rightarrow\beta
$$
где $B$~--- это новый нетерминал.
\item[3.] Квадратные скобки $\Ms{[\ ]}$~--- необязательная конструкция.
Цепочка, заключенная в квадратные скобки, может присутствовать или
отсутствовать. Правило
$$
A\rightarrow \alpha\Ms{[}\beta\Ms{]}\gamma
$$
эквивалентно паре правил
$$
A\rightarrow \alpha\gamma \qquad A\rightarrow\alpha\beta\gamma.
$$
\end{enumerate}
\end{frame}

\begin{frame}
\frametitle{1.1.8. Языки записи грамматик (3/9)}
\begin{enumerate}
\item[4.] Фигурные скобки $\Ms{\{\ \}}$~--- повторение конструкции.
Цепочка, заключенная в фигурные скобки, может повторяться
несколько раз. Правило
$$
A\rightarrow \alpha\Ms{\{}\beta\Ms{\}}\gamma
$$
эквивалентно тройке правил
$$
A\rightarrow \alpha B\gamma \qquad B\rightarrow\beta  \qquad B\rightarrow B \beta
$$
где $B$~--- это новый нетерминал.
\item[5.] Диез $\Ms{\#}$ в фигурных скобках~--- итерирование одной конструкции через другую конструкцию.
Этот метасимвол является обобщением предыдущего:
\\
$\Ms{\{}\alpha\Ms{\}}=\Ms{\{}\alpha\Ms{\#}\varepsilon\Ms{\}}$. Правило
$$
A\rightarrow \alpha\Ms{\{}\beta\Ms{\#}\gamma\Ms{\}}\delta
$$
эквивалентно тройке правил
$$
A\rightarrow \alpha B\delta \qquad B\rightarrow\beta 
\qquad B\rightarrow B
\gamma\beta
$$
где $B$~--- это новый нетерминал.
\end{enumerate}
\end{frame}
\begin{frame}
\frametitle{1.1.8.Языки записи грамматик (4/9)}
\begin{enumerate}
\item[6.] Точка с запятой $\Ms{;}$~--- разделитель правил в множестве правил.
Строчка
$$
A\rightarrow \alpha \Ms{;} B\rightarrow\beta
$$
эквивалентна паре правил
$$
A\rightarrow \alpha \qquad B\rightarrow\beta 
$$
\item[7.] Точка $\Ms{.}$~--- признак конца множества правил.
\item[8.] Двоеточие $\Ms{:}$~--- замена метасимвола
$\rightarrow$, которого нет на стандартной клавиатуре.
\end{enumerate}

\smallskip
\begin{remark}
Обычно считают, что метасимволы 
$\Ms{[\ ]}$, $\Ms{\{\ \}}$ и $\Ms{\#}$ также вводят
неявный нетерминал, а также что символ разделения альтернатив имеет приоритет перед всеми скобками, так что писать дополнительные круглые скобки
не обязательно.
\end{remark}

\smallskip
\begin{example}
В приведённых обозначениях грамматику целых чисел с возможным знаком можно
описать следующим образом:
$$N \to \Ms{[} + \Ms{|} - \Ms{]}
\Ms{(}  1 \Ms{|} 2 \Ms{|} 3 \Ms{|}
 4 \Ms{|} 5 \Ms{|} 6 \Ms{|} 7 \Ms{|} 8 
 \Ms{|} 9 \Ms{)}
\Ms{\{} 0 \Ms{|} 1 \Ms{|} 2 \Ms{|} 3 \Ms{|}
 4 \Ms{|} 5 \Ms{|} 6 \Ms{|} 7 \Ms{|} 8 
 \Ms{|} 9 \Ms{\}} ~\Ms{|}~ 0 ~\Ms{.}$$
Здесь не порождаются ведущие нули, 
а ноль знака не имеет. 
В такой записи грамматики отсутствует явная рекурсия и не используются
вспомогательные нетерминалы.
\end{example}
\end{frame}

\begin{frame}
\frametitle{1.1.8. Языки записи грамматик (5/7)}

\vspace{-2pt}
Иногда метасимволы $\Ms{\{\ \}}$ определяют так, чтобы конструкция
повторялась 0 или более раз, т.е. могла отсутствовать. В этом
случае правило
$$
A\rightarrow \alpha\Ms{\{}\beta\Ms{\}}\gamma
$$
считается эквивалентным трём правилам
$$
A\rightarrow \alpha B\gamma 
\qquad B\rightarrow\varepsilon 
\qquad B\rightarrow B \beta
$$
где $B$~--- это новый нетерминал, а $\varepsilon$~--- пустая
цепочка.

\begin{remark} Символы 
$\Ms{:\ |\ (\ )\ [\ ]\ \{\ \}\ \#\ ;\ .}$
присутствуют на клавиатуре и ими 
удобно пользоваться при описании 
грамматики. Однако эти же символы 
часто используются как терминальные 
символы языка, 
и тогда встаёт вопрос,
как отличить терминальный символ
от метасимвола? Здесь этот вопрос
решён {\em визуально}, 
за счёт использования цвета.
Если система программирования
не допускает различения символов по цвету,
то приходится придумывать 
дополнительные средства, например, заключать
терминальные символы в кавычки $"\ "$.
\end{remark}

Нетрудно убедиться, что все эти расширения (а также некоторые
другие, применяемые на практике) не выводят за пределы
контекстно-свободных грамматик типа $\frak{T}_2$.

\end{frame}

\begin{frame}
\frametitle{1.1.8. Языки записи грамматик (6/9)}

Таким образом, сосуществуют и используются 
{\em несколько языков} для записи порождающих
формальных грамматик. Такая множественность 
языков и диалектов является скорее правилом,
нежели исключением. Какой вариант использовать~---
дело вкуса. 

\smallskip
Очевидно, что использование метасимволов
сокращает запись грамматики делает её более наглядной
(при том условии, однако, что читатель грамматики
точно знает значения всех метасимволов).
С другой стороны, использование метасимволов
существенно усложняет, например, доказательство
эквивалентности грамматик, поскольку
приходится рассматривать в доказательстве 
каждый метасимвол отдельно. 
 
\begin{example} 
На следующем слайде для сравнения приведены четыре варианта
записи грамматики двух языков:
строгого языка  записи грамматик с использованием $\rightarrow$ и
расширенного языка с использованием метасимволов.
Для наглядности в этих записях использованы
одинаковые нетерминальные символы:
$G$~--- собственно грамматика,
$N$~--- алфавит нетерминальных символов целевого языка,
$T$~--- алфавит терминальных символов целевого языка,
$R$~--- правило грамматики,
$S$~--- цепочка символов.
Все терминальные символы заключены в кавычки $"\ "$.
Правила строгого языка считаются разделенными
некоторым непечатаемым символом (конец строки, пробел, ...),
который здесь обозначен $\sqcup$.
\end{example}

 


\end{frame}

\begin{frame}
\frametitle{1.1.8. Языки записи грамматик (7/9)}
\begin{columns}[t]
\begin{column}{7.0cm}
Грамматика строгого языка, записанная на строгом языке

$G \rightarrow R$

$G \rightarrow R\sqcup G$

$R \rightarrow N\ "\!\rightarrow\!"\ S$

$S \rightarrow N $

$S \rightarrow T $

$S \rightarrow N S$

$S \rightarrow T S$

\end{column} 
\begin{column}{8.0cm}
Грамматика расширенного языка, записанная на строгом языке

\begin{columns}[t]
\begin{column}{3.7cm}
$G \rightarrow R\ \ "."$

$G \rightarrow R\ ";"\ G$

$R  \rightarrow N\ "\!:\!"\ S$

$S \rightarrow N $

$S \rightarrow T $

$S \rightarrow N S$

$S \rightarrow T S$
\end{column}
\begin{column}{4.3cm}
$S \rightarrow S\  "\!\mid\!"\  S $

$S \rightarrow "\{"\ S\  "\#"\  S\  "\}"$

$S \rightarrow "("\ S\ ")" $

$S \rightarrow "["\ S\ "]" $

$S \rightarrow "\{"\ S\ "\}" $
\end{column}
\end{columns}
\end{column} 
\end{columns}

\medskip
\begin{columns}[t]
\begin{column}{7.0cm}
Грамматика строгого языка, записанная на расширенном языке

$G\ :\ \{\ R\ \#\ \sqcup\ \}\ ; $

$R\ :\ N\ "\!\rightarrow\!"\ S\ ;$

$S\ :\ \{\ N \mid T \}\ . $
\end{column} 
\begin{column}{8.0cm}
Грамматика расширенного языка, записанная на расширенном языке

$G\ :\ \{\ R\ \#\ ";" \}\ "."\ ; $

$R\ :\ N\ \ "\!:\!"\ S\ ;$

$S\ :\ \{\ N \mid T \}\ \ \mid \ \ 
"(" S\ \#\ "\!\mid\!"\ ")" \ \ \mid \ \ \hfill\linebreak 
"["\ S\ \#\ "\!\mid\!"\ "]" \ \ \mid \ \ 
"\{" S \ [\ "\#"\ S\ ]\ "\}" \ 
. $
\end{column} 
\end{columns}
\end{frame}

\begin{frame}
\frametitle{1.1.8. Языки записи грамматик (8/9)}
Почти одновременно с появлением теории формальных
грамматик для
практических целей был предложен целый ряд других средств.
Исторически самым заметным явилось описание языка программирования
Алгол-60 с помощью \odmkey{нормальных форм
Бэкуса-Наура}{Нормальная форма!Бэкуса-Наура}
 (сокращенно БНФ).
Описание языка c помощью нормальных форм 
Бэкуса\footnote{Джон Бэкус (1924--2007)}--Наура\footnote{Питер Наур (1928--2016)} 
состоит из набора 
\odmkey{основных символов}{Символ!основной} языка, 
набора названий конструкций (или
\odmkey{металингвистических
переменных}{Переменная!металингвистическая}) и набора
\odmkey{металингвистических формул}{Формула!металингвистическая}.
Каждая металингвистическая формула состоит
из левой части, в которой
стоит металингвистическая переменная, \odmkey{металингвистической
связки}{Связка!металингвистическая} $::=$ (читается <<это есть>>)
и правой части, в которой указывается один или несколько вариантов
построения конструкции в левой части. Каждый вариант~--- это
последовательность металингвистических переменных и основных
символов. Варианты разделяются символом $|$ (читается <<или>>).
Металингвистические переменные заключаются в угловые скобки, чтобы
не путать их с основными символами. Считается, что метасимволы $<,
>$, $|$ и $::=$ не входят в число основных символов.

\end{frame}

\begin{frame}
\frametitle{1.1.8. Языки записи грамматик (9/9)}


\begin{example}
Язык целых десятичных чисел со знаком в обозначениях Бэкуса-Наура:\\
<Целое со знаком> $::=$ <Целое без знака> | $+$ <Целое без знака> |
$-$ <Целое без знака> \\
<Целое без знака> $::=$ 0 | <Цифра не ноль> <Последовательность цифр>\\
<Последовательность цифр> $::=$ | <Последовательность цифр> <Цифра>\\
<Цифра> ::= 0 | <Цифра не ноль> \\
<Цифра не ноль> ::=  1 | 2 | 3 | 4 | 5 | 6 | 7 | 8 | 9
\end{example}


Совершенно ясно, что нормальные формы Бэкуса-Наура~--- это один из
способов записать контекстно-свободную грамматику, 
то есть грамматику типа $\frak{T}_2$.  Действительно, основные
символы~--- это алфавит терминальных символов $\Sigma$,
металингвистические переменные~--- это алфавит нетерминальных
символов $\mathrm{M}$, 
металингвистические формулы~--- это множество правил
вывода $\Pi$, в которых вместо символа $\to$
использована связка $::=$, а аксиома грамматики обычно ясна из контекста.

\begin{remark}
Форма Бэкуса--Наура
была предназначена, по замыслу авторов,
прежде всего
для точного и в то же время
{\em удобочитаемого}
описания синтаксиса
языков программирования,
поэтому используются развёрнутые названия конструкций.
\end{remark}
\end{frame}

\subsubsection{1.1.9. Альтернативные средства задания синтаксиса}\label{1.1.9.}
\begin{frame}
\frametitle{1.1.9. Альтернативные средства задания синтаксиса (1/9)}

\begin{digress}
Нормальные формы Бэкуса--Наура 
оказали большое влияние 
на развитие техники определения 
искусственных языков, в частности, языков программирования.
В язык БНФ были добавлены различные усовершенствования,
подобные рассмотренным выше метасимволам,
и появились различные диалекты расширенной БНФ (РБНФ).
До сих пор эта форма 
описания синтаксиса признаётся в качестве 
стандарта
и используется как авторами языков,
так и разработчиками трансляторов.
Если бы дело ограничивалось 
сравнительно немногочисленными языками
программирования общего назначения,
то металингвистических формул было бы достаточно.
Но для специализированных языков 
предметных областей формы Бэкуса--Наура
оказались недостаточно {\em удобны}.
Требуются более лаконичные, более 
наглядные, более современные средства,
поскольку языков предметных областей 
в сотни раз больше, время их жизни в десятки
раз меньше, а пользователями
описания языка являются специалисты предметной
области, не знающие и не нуждающиеся в 
знаниях тонкостей синтаксического анализа 
и компиляции.
\end{digress}

\smallskip
В этом параграфе рассматриваются три разных 
подхода к заданию синтаксиса: 
использование особых графов (синтаксические диаграммы Вирта), 
использование языков разметки (XML) 
и использования языков задания иерархических структур (YAML).  
\end{frame}

\begin{frame}
\frametitle{1.1.9. Альтернативные средства задания синтаксиса (2/9)}
\odmkey{Синтаксические диаграммы}{Диаграмма!синтаксическая}~---
это особый вид орграфов, предложеный Виртом\footnote{Никлаус Вирт (1934--2024)}
для записи синтасиса языка Паскаль. 
Каждая синтаксическая диаграмма~--- это сеть с одним
источником и одним стоком, которые называются <<вход>> и <<выход>>,
соответственно. По традиции вход и выход не изображают на
диаграмме как узлы, а рисуют приходящую <<из ниоткуда>> входную
стрелку и уходящую <<в никуда>> выходную стрелку. Узлы на
синтаксической диаграмме изображают двух типов: с прямыми углами
(они соответствуют нетерминалам) и со скругленными углами (они
соответствуют терминалам). Диаграммы имеют имя, оно обычно пишется
в левом верхнем углу диаграммы.

\begin{example}
На рисунке приведены три диаграммы, $A$, $B$ и $C$. Диаграмма
$A$~--- это \odmkey{тривиальная}{Диаграмма!тривиальная} диаграмма,
она имеет только вход, выход и соединяющую их дугу, которую мы
видим на диаграмме. Диаграммы $B$ и $C$~--- это
\odmkey{элементарные}{Диаграмма!элементарная} диаграммы, они имеют
один узел, помеченный буквой (терминалом или нетерминалом) и входную и выходную дуги.

\vspace{-4pt}
\begin{figure}[h]
\begin{center}
\includegraphics[width=120mm]{1_1_9_2.jpg}
\end{center}
\end{figure}
\end{example}

\end{frame}

\begin{frame}
\frametitle{1.1.9. Альтернативные средства задания синтаксиса (3/9)}
Над синтаксическими диаграммами определены операции
\odmkey{последовательного соединения}{Операция!последовательного соединения} 
(рисунок ниже, диаграмма $E$),
\odmkey{параллельного соединения}{Операция!параллельного соединения} (диаграмма $F$) и
\odmkey{итерации}{Операция!итерации} (диаграмма $G$), которые позволяют
строить более сложные диаграммы из более простых. На
рисунке $\alpha$ и $\beta$~--- это {\em любые} диаграммы, сколь угодно сложные, но имеющие один вход и один выход.
\begin{figure}[h]
\begin{center}
\includegraphics[width=140mm]{1_1_9_3.jpg}
\end{center}
\end{figure}

Язык в целом описывается {\em совокупностью} синтаксических диаграмм.
Каждая конструкция языка (то есть нетерминальный символ или
металингвистическая переменная) определяется на отдельной
диаграмме и записывается в качестве названия этой диаграммы. 
Один нетерминал должен быть указан как начальный (аксиома).
При этом
диаграмма определяет нетерминал в следующем смысле. Каждый путь на
синтаксической диаграмме, ведущий от входа к выходу~--- это
цепочка, выводимая из данного нетерминала.

\end{frame}

\begin{frame}
\frametitle{1.1.9. Альтернативные средства задания синтаксиса (4/9)}

\begin{example}
На рисунке представлена пара синтаксических диаграмм,
определяющих целые числа со знаком.
\begin{figure}[h]
\begin{center}
%\includegraphics[width=120mm]{odm35-03.eps}
\TODO{диаграммы}
\end{center}
\end{figure}
\end{example}

\begin{example}
На рисунке представлены синтаксические диаграммы
строгого языка описания грамматик из предыдущего параграфа.
\begin{figure}[h]
\begin{center}
\TODO{диаграммы}
\end{center}
\end{figure}
\end{example}


Говорят, что совокупность синтаксических диаграмм
\odmkey{правильно построена}{Диаграмма!синтаксическая!правильно
построенная}, если для любого нетерминала есть ровно одна
синтаксическая диаграмма и все синтаксические диаграммы построены
из тривиальных и элементарных с помощью операций последовательного
и параллельного соединения и итерации.

\end{frame}

	% Иллюстрации к слайду подготовил Буслама Анис
	% Дата 10.01.2026
\begin{frame}
	\frametitle{1.1.9. Альтернативные средства задания синтаксиса (4/9)}
	
	\begin{example}
		На рисунке представлена пара синтаксических диаграмм,
		определяющих целые числа со знаком.
		\begin{figure}[h]
			\begin{center}
				\includegraphics[width=95mm]{1.1.9_4_1.png}
			\end{center}
		\end{figure}
	\end{example}
	
	\begin{example}
		На рисунке представлены синтаксические диаграммы
		строгого языка описания грамматик из предыдущего параграфа.
		\begin{figure}[h]
			\begin{center}
				\includegraphics[width=95mm]{1.1.9_4_2.png}
			\end{center}
		\end{figure}
	\end{example}
	
	
	Говорят, что совокупность синтаксических диаграмм
	\odmkey{правильно построена}{Диаграмма!синтаксическая!правильно
		построенная}, если для любого нетерминала есть ровно одна
	синтаксическая диаграмма и все синтаксические диаграммы построены
	из тривиальных и элементарных с помощью операций последовательного
	и параллельного соединения и итерации.
	
\end{frame}

\begin{frame}
\frametitle{1.1.9. Альтернативные средства задания синтаксиса (5/9)}

\begin{theorem}
Для любой правильно построенной совокупности синтаксических
диаграмм существует эквивалентная 
контекстно-свободная грамматика.
\end{theorem}

\begin{proof}
Рассмотрим синтаксическую диаграмму для некоторого нетерминала
$A$. Построим совокупность контекстно-свободных правил, так что
каждому пути в диаграмме будет соответствовать вывод из $A$ по
данным правилам и обратно. Для этого введем необходимое количество
новых нетерминалов, соответствующих структурным элементам
диаграммы, построенной по правилам, приведённым на
предыдущих слайдах, и добавим правила в
соответствии со следующей таблицей.

\begin{tabular}{|l|l|l|}
\hline
Диаграмма & Рисунок & Правила \\
\hline
Тривиальная &  $A$ & $A\rightarrow\varepsilon$ \\
Элементарная с терминалом $a$ & $B$ & $B\rightarrow a$ \\
Элементарная с нетерминалом $D$& $C$ & $C\rightarrow D$ \\
Последовательное соединение $\alpha$ и $\beta$&  $E$ & $E\rightarrow \alpha\beta$ \\
Параллельное соединение $\alpha$ и $\beta$& $F$ & $F\rightarrow \alpha\mid\beta$\\
Итерация $\alpha$ через $\beta$& $G$Е & $G\rightarrow \alpha\mid E \beta\alpha$\\
\hline
\end{tabular}

Очевидно, что правила в правом столбце контекстно-свободные.
\end{proof}

\end{frame}

\begin{frame}
\frametitle{1.1.9. Альтернативные средства задания синтаксиса (6/9)}
\begin{proof} (Окончание)
Далее, по всякому пути на диаграмме можно построить вывод цепочки,
прочитываемой вдоль пути. Вывод строится в обратном порядке:
сначала мы имеем путь и заменяем в нем прохождение пустых ребер на
нетерминалы по $\varepsilon$-правилам (см. первую строку таблицы),
затем заменяем элементарные диаграммы на нетерминалы, затем пары
нетерминалов в соответствии с операциями и так далее, пока не
получим один нетерминал~--- название диаграммы.

Обратно, по всякому выводу в построенном множестве правил, который
заканчивается цепочкой, состоящей из терминалов и нетерминалов,
присутствующих на диаграмме в виде узлов, можно указать путь,
соответствующий данному выводу.

Построение правил (введение дополнительных нетерминалов)
производится для каждой диаграммы независимо (все новые
нетерминалы различны). Поэтому выводу в грамматике в целом
соответствует путь в иерархии диаграмм и обратно.
\end{proof}

\begin{digress}
Синтаксические диаграммы полезно сопоставить с источниками
(см.~\DMref{п.~1.2.3}{1.2.3.}). 
И в том и в другом случае для описания
языка используется граф специального вида, а
язык определяется как множество некоторых путей в таком
графе. Но в источниках <<лингвистическая>> нагрузка приходится на
ребра, а в синтаксических диаграммах~--- на вершины.
\end{digress}

\end{frame}

\begin{frame}
\frametitle{1.1.9. Альтернативные средства задания синтаксиса (7/9)}
Широко используемым на практике средством задания синтаксиса 
являются \odmkey{языки разметки}{Язык!разметки}.

\TODO{Кратко идея SGML, далее HTML, XML и XML Scheme 
с обязательным сопоставлением с КСГ.
Ещё можно (но не очень нужно) замечание про Mark Down}
\end{frame}

\begin{frame}
\frametitle{1.1.9. Альтернативные средства задания синтаксиса (8/9)}
Сосуществуют и активно используются возникшие в разное время 
и по различным причинам несложные, но полезные средства 
описания иерархических текстовых структур.

\TODO{Далее можно YAML или JSON или что угодно, но обязательно
с переводом с птичьего языка на язык КСГ. }
 
\end{frame}

\begin{frame}
\frametitle{1.1.9. Альтернативные средства задания синтаксиса (9/9)}

\TODO{Опционально: слова про метод раскрутки (bootstapping), поднять себя за шнурки ботинок, с демонстрацией выразительной силы любого из рассмотренных средств. 
Нечто похожее на слайд 1.1.8.(7/9)}

\end{frame}

%Добавил пробел в названии слайда Штейнберг Георгий 5030102/20202 
% Дата 10.01.26
\begin{frame}
\frametitle{1.1.8. Языки записи грамматик (4/9)}
\begin{enumerate}
\item[6.] Точка с запятой $\Ms{;}$~--- разделитель правил в множестве правил.
Строчка
$$
A\rightarrow \alpha \Ms{;} B\rightarrow\beta
$$
эквивалентна паре правил
$$
A\rightarrow \alpha \qquad B\rightarrow\beta 
$$
\item[7.] Точка $\Ms{.}$~--- признак конца множества правил.
\item[8.] Двоеточие $\Ms{:}$~--- замена метасимвола
$\rightarrow$, которого нет на стандартной клавиатуре.
\end{enumerate}

\smallskip
\begin{remark}
Обычно считают, что метасимволы 
$\Ms{[\ ]}$, $\Ms{\{\ \}}$ и $\Ms{\#}$ также вводят
неявный нетерминал, а также что символ разделения альтернатив имеет приоритет перед всеми скобками, так что писать дополнительные круглые скобки
не обязательно.
\end{remark}

\smallskip
\begin{example}
В приведённых обозначениях грамматику целых чисел с возможным знаком можно
описать следующим образом:
$$N \to \Ms{[} + \Ms{|} - \Ms{]}
\Ms{(}  1 \Ms{|} 2 \Ms{|} 3 \Ms{|}
 4 \Ms{|} 5 \Ms{|} 6 \Ms{|} 7 \Ms{|} 8 
 \Ms{|} 9 \Ms{)}
\Ms{\{} 0 \Ms{|} 1 \Ms{|} 2 \Ms{|} 3 \Ms{|}
 4 \Ms{|} 5 \Ms{|} 6 \Ms{|} 7 \Ms{|} 8 
 \Ms{|} 9 \Ms{\}} ~\Ms{|}~ 0 ~\Ms{.}$$
Здесь не порождаются ведущие нули, 
а ноль знака не имеет. 
В такой записи грамматики отсутствует явная рекурсия и не используются
вспомогательные нетерминалы.
\end{example}
\end{frame}



% Исправил нумерацию слайда (было 5/7) Штейнберг Георгий 5030102/20202 
% Дата 10.01.26
\begin{frame}
\frametitle{1.1.8. Языки записи грамматик (5/9)}

\vspace{-2pt}
Иногда метасимволы $\Ms{\{\ \}}$ определяют так, чтобы конструкция
повторялась 0 или более раз, т.е. могла отсутствовать. В этом
случае правило
$$
A\rightarrow \alpha\Ms{\{}\beta\Ms{\}}\gamma
$$
считается эквивалентным трём правилам
$$
A\rightarrow \alpha B\gamma 
\qquad B\rightarrow\varepsilon 
\qquad B\rightarrow B \beta
$$
где $B$~--- это новый нетерминал, а $\varepsilon$~--- пустая
цепочка.

\begin{remark} Символы 
$\Ms{:\ |\ (\ )\ [\ ]\ \{\ \}\ \#\ ;\ .}$
присутствуют на клавиатуре и ими 
удобно пользоваться при описании 
грамматики. Однако эти же символы 
часто используются как терминальные 
символы языка, 
и тогда встаёт вопрос,
как отличить терминальный символ
от метасимвола? Здесь этот вопрос
решён {\em визуально}, 
за счёт использования цвета.
Если система программирования
не допускает различения символов по цвету,
то приходится придумывать 
дополнительные средства, например, заключать
терминальные символы в кавычки $"\ "$.
\end{remark}

Нетрудно убедиться, что все эти расширения (а также некоторые
другие, применяемые на практике) не выводят за пределы
контекстно-свободных грамматик типа $\frak{T}_2$.

\end{frame}


% Текст слайда подготовил Штейнберг Георгий 5030102/20202 
% Дата 10.01.26
\begin{frame}
\frametitle{1.1.9. Альтернативные средства задания синтаксиса (7/9)}
Широко используемым на практике средством задания синтаксиса 
являются \odmkey{языки разметки}{Язык!разметки}~--- это языки для описания структуры и представления текста.
\begin{itemize}

 \item\odmkey{SGML}{Язык!разметки!SGML}~--- это стандарт, определяющий правила построения языков разметки. Он использует \odmkey{DTD}{Язык!разметки!DTD} (Document Type Definition), чтобы задать множество тегов (нетерминалов) и правила их вложенности. Так как SGML допускает множество сокращений (например, незакрытые теги), написание парсеров для него затруднительно.
 
 \item\odmkey{HTML}{Язык!разметки!HTML}~--- это язык разметки, изначально написанный на SGML, его множество нетерминалов и правила фиксированы стандартом.
 
 \item\odmkey{XML}{Язык!разметки!XML}~--- это доработанный SGML со строгим синтаксисом, используемый для хранения и передачи структурированных данных. Благодаря строгому синтаксису парсеры для XML проще и работают быстрее. Для описания структуры XML-документов используется \odmkey{XML Schema}{Язык!разметки!XML Schema}, который позволяет задать грамматику.

\end{itemize}
 \begin{remark}
 SGML и XML являются инструментами для определения грамматик, DTD и XML Scheme~--- грамматиками, а HTML~--- языком, порожденным грамматикой.
 \end{remark}
\end{frame}




% Текст слайда подготовил Штейнберг Георгий 5030102/20202 
% Дата 10.01.26
\begin{frame}[fragile]
\frametitle{1.1.9. Альтернативные средства задания синтаксиса (8/9)}
    \smallskip
\begin{example}
Сопоставление с КСГ $G=\langle \mathrm{M}, \Sigma, \Pi, S\rangle$:
\begin{enumerate}
    \item Схема документа (DTD, XML Schema)~--- Грамматика $G$.
    \item Множество тегов (Elements)~--- Алфавит нетерминалов $\mathrm{M}$.
    \item Текстовое содержимое (PCDATA)~--- Алфавит терминалов $\Sigma$.
    \item Правила вложенности тегов~--- Правила вывода $\Pi$.
    \item Корневой элемент (Root)~--- Аксиома грамматики $S$.
\end{enumerate}
 Например, язык целых десятичных чисел со знаком XML-документ с внутренней схемой (Internal DTD):
\scriptsize % Уменьшаем шрифт, чтобы код влез на слайд
\begin{verbatim}
<?xml version="1.0" encoding="utf-8"?>
<!DOCTYPE signedInteger [
  <!ELEMENT signedInteger (sign?, unsignedInteger)>
  <!ELEMENT sign (#PCDATA)>
  <!ELEMENT unsignedInteger (zero | (nonzero, digit*))>
  <!ELEMENT zero (#PCDATA)> <!ELEMENT nonzero (#PCDATA)> 
  <!ELEMENT digit (#PCDATA)> ]>
\end{verbatim}
\end{example}

\TODO{Замечания ФАН 21 января. Вклеить в структуру, доработать примеры.}
\end{frame}





\begin{frame}
\frametitle{1.1.А. Онтология формальных языков, грамматик и автоматов}
\end{frame}

\begin{frame}
\frametitle{1.1.Б. Темы для самостоятельной работы студентов}

\begin{enumerate}
\item  \SW{Грамматические описания нетекстовых языков} \\
В некоторых случаях формализм порождающих
грамматик удаётся применить, когда знаками языка являются не символы в тексте, а другие объекты. Например, фигуры и линии на диаграмме графа, 
жесты пользователей смартфонов, ноты в партитуре и тому подобное. 
Требуется подготовить небольшой, но представительный набор убедительных примеров.

\item \SW{История развития и перспективы теории формальных грамматик} \\ 
Требуется подготовить познавательный научно-популярный доклад с фотографиями ученых,  датами, событиями, примерами применения и прогнозами на будущее.

\item \SW{Классификация формальных языков на основе ...} \\ 
Классификация языков на основе формальных грамматик уже развита Хомским и другими учёными.
Предлагается попробовать построить классификацию на основе использования 
иных формализмов. Например, провести классификацию языков,
которые могут быть заданы с помощью графов и орграфов различных видов.

\end{enumerate}

\end{frame}


\subsection{1.2. Регулярные выражения}
\begin{frame}
\frametitle{1.2. Регулярные выражения (1/3)}

Порождающие грамматики~--- не единственное средство определения
языков. Языки являются множествами, поэтому над ними определены
все операции над множествами 
(объединение, пересечение, разность и 
симметрическая разность). Кроме того, для
элементов языков~--- слов~--- определена операция конкатенации,
что позволяет определить для языков и другие операции. 
Таким образом,
существует возможность определить язык
{\em алгебраически}, как выражение,
построенное из более простых элементов.

Основополагающий вклад в развитие теории регулярных выражений 
внес Клини.\footnote{Стивен Коул Клини (1909--1994)}

\medskip
Удивительный факт: далёкие как от исчислений (формальные грамматики),
так и от алгебраических структур (регулярные выражения) формализмы,
в частности, ориентированные {\em графы}, при использовании для определения 
языков приводят к одним и тем же результатам. Обсуждение этого наблюдения
начинается в этом разделе, а продолжается и заканчивается в разделе
\DMref{2.2}{2.2}.     

\end{frame}

\begin{frame}
\frametitle{1.2. Регулярные выражения (2/3)}

\begin{tabular}{p{10mm}|p{40mm}|p{95mm}}
  № & Название параграфа
  & Ключевые термины и обозначения \\
  \hline
\DMref{1.2.1.}{1.2.1.} &  Регулярные операции
  & { регулярные операции, 
  объединение, альтернация, сложение, умножение, конкатенация, сцепление,
  итерация, замыкание Клини}
   \\
\DMref{1.2.2.}{1.2.2.} & Алгебра регулярных выражений 
  & 
{ элементарный язык,
  регулярный язык, регулярное выражение, 
  алгебра регулярных множеств $\cal{R}$, алгебра регулярных выражений $\frak{R}$,
  эквивалентность регулярных выражений}
  \\
\DMref{1.2.3.}{1.2.3.}  & Языковые графы &
{граф языка, источник, полный источник, эквивалентные источники,
  двухполюсник, начальные и заключительные узлы, значение пути, размер источника,
  замкнутое множество узлов } 
  \\ 
\DMref{1.2.4.}{1.2.4.}  & Детерминированные и недетерминированные источники &
{ теорема детерминизации } 
\end{tabular}

\end{frame}


\begin{frame}
\frametitle{1.2. Регулярные выражения (3/3)}

\begin{tabular}{p{10mm}|p{45mm}|p{90mm}}
  № & Название параграфа
  & Ключевые термины и обозначения \\
  \hline
\DMref{1.2.5.}{1.2.5.} & Расширение языка регулярных выражений &
{операции повторения, итерация Цейтина, символы шаблона, символьные классы}\\   
\DMref{1.2.6.}{1.2.6.} & Программирование в регулярных выражениях &
{POSIX, PERL, SNOBOL4, REFAL }\\   
\DMref{1.2.7.}{1.2.7.} & Уравнения с регулярными коэффициентами &
{ }\\   
\DMref{1.2.8.}{1.2.8.} & Характеристика класса регулярных языков &
{язык Дика  }   
\end{tabular}

\end{frame}


\subsubsection{1.2.1. Регулярные операции}
\label{1.2.1.}
\begin{frame}
\frametitle{1.2.1. Регулярные операции (1/3)}

Пусть $L_1, L_2\subset A^*$. Рассмотрим следующие три операции над языками $L_1, L_2$.

\odmkey{Объединение}{Объединение!языков}
(\odmkey{альтернация}{Операция!альтернации},
\odmkey{сложение}{Сложение!языков}), обозначается знаками $\cup, ~|~, +$:
$$
L_1\cup L_2 = L_1 ~|~ L_2 = L_1 + L_2\Def \Set{\alpha \in A^*}{\alpha \in L_1 \Or \alpha\in L_2}.$$

\begin{remark}
В определении указано,
что языки определены 
над одним алфавитом, \\$L_1, L_2\subset A^*$,
но это сделано только ради удобства 
обозначений и не ограничивает общности.
Можно считать, что $L_1\subset A_1^*$, 
$L_2\subset A_2^*$, и тогда $L_1\cup L_2\subset (A_1\cup A_2)^*$.
\end{remark}

\odmkey{Умножение}{Умножение!языков}
(\odmkey{конкатенация}{Операция!конкатенации},
\odmkey{сцепление}{Сцепление!языков}), обозначается знаком $\cdot$ или
опускается:
$$
L_1 L_2= L_1\cdot L_2\Def \Set{\alpha\in A^*}
{\alpha=\alpha_1\alpha_2 \And \alpha_1\in L_1 \And \alpha_2\in
L_2}.
$$

\begin{remark}
Можно сцеплять язык сам с собой:
$L, LL, LLL, \dots$,
что естественно обозначить с
помощью показателя степени,
то есть $L^1\Assign L, L^2\Assign LL,
L^3\Assign LLL$ и так далее.
При этом по определению 
считают, что $L^0\Assign \{\varepsilon\}$.  
\end{remark}

\odmkey{Итерация}{Итерация!языков} (\odmkey{замыкание Клини}{Замыкание!языков}), обозначается знаком $*$:

\vspace{-12pt}
$$
L_1^*\Def \varepsilon\cup L_1\cup L_1 L_1\cup\dots =
\bigcup_{i=0}^\infty L_1^i
.$$

\end{frame}

\begin{frame}
\frametitle{1.2.1. Регулярные операции (2/3)}

Введённые операции называются
\odmkey{регулярными}{Операция!регулярная}. Эти операции имеют
особое значение для теории языков и автоматов.

\medskip
Поскольку результатом регулярной 
операции над языками является язык,
из этих операций с помощью суперпозиции 
можно строить
сколь угодно сложные  выражения,
которые называются \odmkey{регулярными}{Выражение!регулярное}. 

\medskip
Для регулярных операций обычно устанавливают следующий приоритет:
$*$, $\cdot$, $|$, что позволяет опускать лишние скобки при записи
выражений.

\begin{remark}
Операция альтернации является
привычной операцией объединения из
наивной теории множеств,
но операции конкатенации и итерации
не являются 
и не выражаются операциями наивной теории множеств. 
\end{remark}

\begin{examples}
Пусть язык 
$L_1 =\{a,b,c,\dots,x,y,z\}$~--- это латинский алфавит, то есть множество однобуквенных слов,
а язык 
$L_2=\{1,2,3,\dots\}$~--- это натуральные
числа в десятичной записи.
Тогда $L_1 L_2$~--- это множество 
идентификаторов, которые начинаются
с буквы и продолжаются цифрами,
$L_1^*$~--- это множество всех 
слов, составленных из букв латинского алфавита,
включая пустое слово $\varepsilon$,
а $L_2^* = L_2$.
\end{examples}

\end{frame}


\begin{frame}
\frametitle{1.2.1. Регулярные операции (3/3)}
\begin{theorem}
Пусть $L_1, L_2, L_3\subset A^*$. Тогда регулярные операции
обладают следующими свойствами.
\begin{enumerate}
  \item Альтернация коммутативна:\\[-2pt]
  $L_1~|~ L_2 = L_2~|~ L_1$;\\[-1pt]
  \item Альтернация ассоциативна:\\[-2pt]
  $L_1~|~ (L_2~|~ L_3) = (L_1~|~ L_2)~|~L_3$;\\[-1pt]
  \item Конкатенация дистрибутивна относительно альтернации слева:\\[-2pt]
  $L_1(L_2~|~ L_3) = L_1 L_2~|~ L_1 L_3$;\\[-1pt]
  \item Конкатенация дистрибутивна относительно альтернации справа:\\[-2pt]
$(L_1~|~ L_2)L_3 = L_1 L_3~|~ L_2 L_3$;\\[-1pt]
  \item Конкатенация ассоциативна:\\[-2pt]
  $(L_1 L_2)L_3 = L_1(L_2 L_3)$;\\[-1pt]
    \item Итерация идемпотентна:\\[-2pt]
  $(L_1^*)^* = L_1^*$;\\[-1pt]
      \item Альтернация идемпотентна:\\[-2pt]
  $L_1~|~L_1 = L_1$.\\[-1pt]
\end{enumerate}
\end{theorem}
\begin{proof}
Очевидно из определений.
\end{proof}
\end{frame}

\subsubsection{1.2.2. Алгебра регулярных выражений}\label{1.2.2.}
\begin{frame}
\frametitle{1.2.2. Алгебра регулярных выражений (1/8)}

Пусть задан конечный алфавит $A=\{\Tuple{a}{1}{k}\}$. Язык $L_{a_i}\Assign\{a_i\}$,
состоящий из одного однобуквенного слова, называется
\odmkey{элементарным}{Язык!элементарный} и обозначается $a_i$.
Пустое слово $\varepsilon$ тоже считается элементарным языком.  

\begin{digress}
Используемые обозначения многозначны.
Поясним это утверждение примером.\\
Пусть $A=\{a,b,c,\dots,x,y,z\}$. 
Тогда символ $a$ именует одновременно три различных объекта:
\begin{enumerate}
\item $a$~--- это буква латинского алфавита, $a\in A$;\\[-2pt]
\item $a$~--- это однобуквенное слово, $a\in A^*$;\\[-2pt]
\item $a$~--- это элементарный язык, $a\subset A^*$.\\[-2pt]
\end{enumerate}
Аналогично символ $\varepsilon$ именует одновременно три различных объекта:
\begin{enumerate}
\item $\varepsilon$~--- это символ {\em греческого}
алфавита, $\varepsilon\not\in A$;\\[-2pt]
\item $\varepsilon$~--- это обозначение пустого слова, $\varepsilon\in A^*$;\\[-2pt]
\item $\varepsilon$~--- это элементарный язык, $\varepsilon\subset A^*$.\\[-2pt]
\end{enumerate}
Такие обозначения удобны, 
потому что лаконичны, и не создают 
путаницы, потому что 
тип объекта всегда восстанавливается из контекста. В частности,
аргументы регулярных операций~--- это 
всегда языки, но не слова и не буквы.
\end{digress} 
\end{frame}

\begin{frame}
\frametitle{1.2.2. Алгебра регулярных выражений (2/8)}


Язык (то есть множество слов в алфавите $A$) называется \odmkey{регулярным}{Язык!регулярный}, если он
может быть построен из элементарных языков (букв алфавита $A$ и пустого слова 
$\varepsilon$)
с помощью суперпозиции регулярных операций. Семейство регулярных языков является подмножеством
множества всех языков над алфавитом $A$.
Множество регулярных языков обозначим $\cal{R}$:
$\cal{R} \subset 2^{A^*}.$
Алгебра $\langle \cal{R}; ~|~, ~\cdot~, ~{}^*\rangle$ называется \odmkey{алгеброй
регулярных множеств}{алгебра!регулярных!множеств}. Эта алгебра
конечнопорождённая; ее образующими являются
элементарные языки.

\smallskip
Каждый регулярный язык может быть описан формулой, которая называется \odmkey{регулярным
выражением}{Выражение!регулярное}. Операндами регулярных выражений являются обозначения 
элементарных языков (то есть буквы алфавита $A$), а операциями~--- регулярные операции 
$|~\cdot~{}^*$. 
Множество регулярных выражений обозначим $\frak{R}$.
Алгебра $\langle \frak{R}; ~|~, ~\cdot~, ~*\rangle$ называется \odmkey{алгеброй
регулярных выражений}{алгебра!регулярных!выражений}.  Эта алгебра
также конечнопорождённая; ее образующими являются
буквы алфавита $A$ и обозначение $\varepsilon$.

\smallskip
Каждое регулярное выражение порождает единственный регулярный язык.
Один и тот же язык может быть порождён бесконечным числом
различных регулярных выражений.
\begin{example}
Различные регулярные выражения 
$\varepsilon,\ \varepsilon\varepsilon,\ \varepsilon~|~\varepsilon,\ \varepsilon^*,\ \dots$
порождают один и тот же регулярный язык $\{\varepsilon\}$. 
\end{example}
  

\end{frame}

\begin{frame}
\frametitle{1.2.2. Алгебра регулярных выражений (3/8)}
Соответствие регулярных выражений и регулярных языков является 
{\em гомоморфизмом}, который уважает все регулярные операции.
Говорят, что регулярные выражения
\odmkey{эквивалентны}{Эквивалентность!регулярных выражений}, если они задают
один и тот же язык. 
Эквивалентность регулярных 
выражений обозначают, как обычно, знаком $=$.
Таким образом, по теореме о гомоморфизме алгебра классов эквивалентности
регулярных выражений {\em изоморфна} алгебре регулярных множеств.
На этом обстоятельстве основана общепринятая в математике практика:
манипулировать конечными дескрипторами (в данном случае регулярными выражениями)
бесконечных объектов (в данном случае, регулярных языков).

\begin{examples}
\begin{enumerate}
  \item Обозначение $A^*$~--- это правильное регулярное выражение:
  $
A^*=(a_1~|~\dots~|~ a_k)^*.
$
  \item Любое слово регулярно: оно построено из букв
  конкатенацией.
  \item Любой конечный язык регулярен:
$
L=\{\alpha_1,\dots, \alpha_k\},\quad
  L=\alpha_1~|~\dots~|~\alpha_k.
$
  \item  Язык $(aba)^*$  регулярен.
\end{enumerate}
\end{examples}



В предыдущем параграфе
установлены семь важных эквивалентностей,
с помощью которых можно установить многие другие
полезные свойства операций алгебры регулярных
выражений  


\end{frame}

\begin{frame}
\frametitle{1.2.2. Алгебра регулярных выражений (4/8)}
Пустое множество $\emptyset\subset A^*$ 
является языком, этот язык считается регулярным по определению, $\emptyset\in R, \emptyset\ne\varepsilon$. 
Напомним, что $A^*= \varepsilon +A + AA+ \dots =
\bigcup_{i=0}^\infty A^i$ и введём дополнительное обозначение:
$A^+ \Def A + AA+ \dots = \bigcup_{i=1}^\infty A^i$.
  
Пользуясь свойствами регулярных операций и определениями, 
нетрудно установить, что в алгебре 
регулярных выражений
выполняются следующие соотношения:

\vspace{-16pt}
\begin{columns}[t]
\begin{column}{7cm}
\begin{enumerate}
\item[8.] $A+\emptyset = A$;
\item[9.] $A\emptyset = \emptyset A=\emptyset$;
\item[10.] $A^* +A =A^*$;
\item[11.] $\emptyset^*=\varepsilon$;
\item[12.] $\varepsilon^*=\varepsilon$;
\item[13.] $A^* + \varepsilon=A^*$;
\end{enumerate}
\end{column}
\begin{column}{7cm}
\begin{enumerate}
\item[14.] $A\varepsilon = \varepsilon A = A$;
\item[15.] $A^+=AA^*$;
\item[16.] $A^+ + A =A^+$;
\item[17.] $\emptyset^+ = \emptyset$;
\item[18.] $\varepsilon^+=\varepsilon$;
\item[19.] $A^+ + \varepsilon=A^*$.
\end{enumerate}
\end{column}
\end{columns} 

\vspace{2pt}
\begin{remark}
Алгебра регулярных выражений является моноидом относительно конкатенации и языка $\varepsilon$,
является коммутативным моноидом
относительно альтернации и пустого языка
$\emptyset$, и является некоммутативным
полукольцом с единицей относительно 
указанных операций и элементов.  
\end{remark}
\end{frame}

\begin{frame}
\frametitle{1.2.2. Алгебра регулярных выражений (5/8)}

Пользуясь указанными простыми свойствами, в алгебре регулярных выражений
можно доказывать различные равенства.

\begin{example}
 Доказать, что
$(b+aa^*b)+(b+aa^*b)(a+ba^*b)^*(a+ba^*b) = a^*b(a+ba^*b)^*$.
Доказательство:
$(b+aa^*b)+(b+aa^*b)(a+ba^*b)^*(a+ba^*b) =$\\
$=(b+aa^*b)(\varepsilon+(a+ba^*b)^*(a+ba^*b)) =
(b+aa^*b)(\varepsilon+(a+ba^*b)^+) = $\\
$=(b+aa^*b)(a+ba^*b)^* = 
(\varepsilon+aa^*)b(a+ba^*b)^* = (\varepsilon+a^+)b(a+ba^*b)^* = a^*b(a+ba^*b)^*$.
\end{example}

Имеют место некоторые менее очевидные свойства.

\begin{lemma}
$\A{L\subset A^*}
   {\A{\alpha\in L^*, \alpha\neq\varepsilon}
      {\E{\Tuple{\alpha}{1}{k}\in L}
         {\alpha=\alpha_1\dots\alpha_k}
      }
   }$.
\end{lemma}
Другими словами, если слово принадлежит итерации языка, то это слово всегда можно
разложить в конкатенацию слов, принадлежащих языку и обратно.  

\begin{theorem}
$(A+B)^*=(A^*+B^*)^* =(A^*B^*)^*$.
\end{theorem}
\begin{proof}
Обозначим $L_1\Assign (A+B)^*, L_2\Assign 
(A^*+B^*)^*, L_3=(A^*B^*)^*$.
\proofsection{$L_1\subset L_2$}
$A+B\subset A^* + B^* \Impl (A+B)^*\subset 
(A^* + B^*)^*$.
\proofsection{$L_2\subset L_1$}
Если $\alpha\in (A^*+B^*)^*$, то по лемме 
$\E{\Tuple{\alpha}{1}{k}\in A^*+B^*}{\alpha=
\alpha_1\dots\alpha_k}$.
Но по той же лемме
$\A{\alpha_i}{\E{\Tuple{\alpha}{i_1}{i_j}\in A+
B}{\alpha_i=\alpha_{i_1}\dots\alpha_{i_j}}}$.
Значит $\alpha\in (A+B)^*$.
\end{proof}
\end{frame}

\begin{frame}
\frametitle{1.2.2. Алгебра регулярных выражений (6/8)}
\begin{proof} (Окончание)
\proofsection{$L_1\subset L_3$}
Если $\alpha\in (A+B)^*$, то по лемме 
$\E{\Tuple{\alpha}{1}{k}\in A+B}{\alpha=
\alpha_1\dots\alpha_k}$.
Но если $\alpha_i\in A+B$, 
то $\alpha_i\in A^*B^*$ 
и значит $\alpha\in (A^*B^*)^*$.
\proofsection{$L_3\subset L_1$}
Если $\alpha\in (A^*B^*)^*$, то по лемме 
$\E{\Tuple{\alpha}{1}{k}\in A^*B^*}{\alpha=
\alpha_1\dots\alpha_k}$.
Далее, если $\alpha_i \in A^*B^*$,
то $\E{\beta_i\in A^{k_1},
\gamma_i \in B^{k_2}}{\alpha_i = \beta_i \gamma_i}$.
Но слова $\beta_i$ и $\gamma_i$
также можно разложить по лемме:
$\E{\Tuple{\beta}{i_1}{i_{k_1}}\!\in\!A}{\beta_i=\beta_{i_1}\dots\beta_{i_{k_1}}}$,
$\E{\Tuple{\gamma}{i_1}{i_{k_2}}\!\in\!B}{\gamma_i=\gamma_{i_1}\dots\gamma_{i_{k_2}}}$,
 и значит слово $\alpha_i$
 разложено в конкатенацию слов,
 каждое из которых принадлежит $A+B$,
 откуда $\alpha\in (A+B)^*$. 
\end{proof}

\begin{corollary} {\NB{[1]}}
$(AB)^*A = A(BA)^*$.
\end{corollary}

\begin{proof} Раскрасим буквы в доказываемом равенстве, чтобы было видно, какому языку принадлежит какое слово:
$(A\textcolor{blue}{B})^*\textcolor{red}{A}=A(BA)^*$.
Если $\alpha\in(A\textcolor{blue}{B})^*\textcolor{red}{A}$,
то  $\alpha=\alpha_1\textcolor{red}{\alpha_2}$, 
причём $\alpha_1\in(A\textcolor{blue}{B})^*\And \alpha_2\in \textcolor{red}{A}$. 
Далее, если $\alpha_1\in(A\textcolor{blue}{B})^*$, 
то по лемме\\ 
$\E{\Tuple{\alpha_{1}}{1}{k}\in A\textcolor{blue}{B}}{\alpha_{1}=\alpha_{1_1}\ldots\alpha_{1_k}}$, 
причём
 $\A{i}{\alpha_{1_i}=\gamma_i\textcolor{blue}{\beta_i}\And
\gamma_i\in A \And \textcolor{blue}{\beta_i}\in \textcolor{blue}{B}}$.\\ 
Таким образом, $\alpha=\gamma_1\textcolor{blue}{\beta_1}\ldots\gamma_k\textcolor{blue}{\beta_k}\textcolor{red}{\alpha_2}\And\forall\gamma,\textcolor{red}{\alpha}\in A,\forall\textcolor{blue}{\beta} \in \textcolor{blue}{B}$. По ассоциативности конкатенации: $\alpha=\gamma_1(\textcolor{blue}{\beta_1}\gamma_2\ldots\textcolor{blue}{\beta_{k-1}}\gamma_k\textcolor{blue}{\beta_k}\textcolor{red}{\alpha_2})$. Значит $\alpha\in A(BA)^*$. Мы доказали, что \\$(AB)^*A\subset A(BA)^*$. Вложение $A(BA)^*\subset (AB)^*A$ доказывается аналогично.
\end{proof}
\end{frame}

\begin{frame}
\frametitle{1.2.2. Алгебра регулярных выражений (7/8)}

\begin{corollary} {\NB{[2]}}
$(A^* B)^* =(A+B)^* B + \varepsilon$.
\end{corollary}
\begin{proof}
 Пусть $\alpha=\varepsilon$. 
 Тогда $\alpha\in (A^*B)^* \And \alpha\in (A+B)^*B+\varepsilon$ и можно сразу предположить, что
 $\alpha\neq\varepsilon$.
\proofsection{$(A^*\textcolor{blue}{B})^*\subset (A+B)^*B+\varepsilon$} Если $\alpha\in(A^*\textcolor{blue}{B})^*$, то по лемме \\
$\E{\alpha_1,\ldots,\alpha_k\in A^*\textcolor{blue}{B}}
{\alpha=\alpha_1\ldots\alpha_k}$,
причём $\alpha_i=\gamma_i\textcolor{blue}{\beta_i}
\And\gamma_i\in A^{j_i}\And\textcolor{blue}{\beta_i}\in \textcolor{blue}{B}$. По лемме 
$\E{\gamma_{i_1},\ldots,\gamma_{i_{j_i}}\in A}
{\gamma_i=\gamma_{i_1}\ldots\gamma_{i_{j_i}}}$. 
В итоге $\alpha=\gamma_{1_1}\ldots\gamma_{1_{j_1}}\textcolor{blue}{\beta_1}\ldots\gamma_{k_1}\ldots\gamma_{k_{j_k}}\textcolor{blue}{\beta_k}$, где\\ 
$\gamma_s\in A, \textcolor{blue}{\beta_r} \in \textcolor{blue}{B}$. Расставим скобки и получим $\alpha=(\gamma_{1_1}\ldots\gamma_{1_{j_1}}\textcolor{blue}{\beta_1}\ldots\gamma_{k_1}\ldots\gamma_{k_{j_k}})\textcolor{blue}{\beta_k}$, 
\\$\alpha\in(A+B)^*B$ или, расширяя множество, $\alpha\in(A+B)^*B+\varepsilon$.
 \proofsection{$(\textcolor{red}{A}+\textcolor{blue}{B})^*B+\varepsilon\subset(A^*B)^*$}. Пусть $\alpha\in(\textcolor{red}{A}+\textcolor{blue}{B})^*B+\varepsilon$ и $\alpha\neq\varepsilon$, тогда $\alpha\in(\textcolor{red}{A}+\textcolor{blue}{B})^*B$. По определению конкатенации $\alpha=\textcolor{magenta}{\gamma}\beta$, где $\gamma\in (\textcolor{red}{A}+\textcolor{blue}{B})^j,\beta\in B$. 
 Слово $\textcolor{magenta}{\gamma}$ раскладываем по лемме: 
 $\E{\textcolor{magenta}{\gamma_1},\ldots,\textcolor{magenta}{\gamma_j}\in (\textcolor{red}{A}+\textcolor{blue}{B})}{\textcolor{magenta}{\gamma}=\textcolor{magenta}{\gamma_1}\ldots\textcolor{magenta}{\gamma_j}}$, $\alpha=\textcolor{magenta}{\gamma_1}\ldots\textcolor{magenta}{\gamma_j}\beta$. Получили $\alpha\in(A^*B)^*$.
\end{proof}

\medskip 
\begin{corollary} {\color{blue}\bf [3]}
$(AB^*)^*=A(A+B)^*+\varepsilon$.
\end{corollary}
\begin{proof}
Упражнение.
\end{proof}

\medskip
Доказанные эквивалентности позволяют упрощать регулярные выражения. 
\end{frame}

\begin{frame}
\frametitle{1.2.2. Алгебра регулярных выражений (8/8)}

\begin{example}
Доказать, что $(ABA(BA)^*B+AB)(A(BA)^*B)^*=A(BA)^*B$. \\
 Для начала проведем несколько преобразований\\ $(ABA(BA)^*B+AB)(A(BA)^*B)^*\overset{14}{=}(ABA(BA)^*B+A\varepsilon B)(A(BA)^*B)^*\overset{3,4}{=}$\\ $\overset{3,4}{=}(A(BA(BA)^*+\varepsilon) B)(A(BA)^*B)^*\overset{15}{=}(A((BA)^++\varepsilon) B)(A(BA)^*B)^*\overset{19}{=}$ \\
 $\overset{19}{=}(A(BA)^* B)(A(BA)^*B)^*\overset{15}{=}(A(BA)^*B)^+$\\
 Теперь осталось доказать, что $(A(BA)^*B)^+=A(BA)^*B$\\
 \proofsection{$A(BA)^*B\subset (A(BA)^*B)^+$} По определению введенной  операции $L^+$.\\
 \proofsection{$A(\textcolor{red}{B}\textcolor{green}{A})^*\textcolor{blue}{B})^+\subset A(BA)^*B$} Заметим, 
 что если пустое слово $\varepsilon$ не содержится ни в языке $A$, ни в языке $B$, то оно не содержится и в выражениях.
  Поэтому можно доказывать не $(A(\textcolor{red}{B}\textcolor{green}{A})^*\textcolor{blue}{B})^+\subset A(BA)^*B$, 
 а $(A(\textcolor{red}{B}\textcolor{green}{A})^*\textcolor{blue}{B})^*\subset A(BA)^*B+\varepsilon$.
 Пусть $\alpha\in(A(\textcolor{red}{B}\textcolor{green}{A})^*\textcolor{blue}{B})^*$. 
По лемме 
$\E{\alpha_1,\ldots,\alpha_k\!\in\!A(\textcolor{red}{B}\textcolor{green}{A})^*\textcolor{blue}{B}}
{\alpha\!=\!\alpha_1\ldots\alpha_k}$, причём
$\alpha_i\!=\!\gamma_i\theta_i\textcolor{blue}{\beta_i}$, 
$\gamma_i\!\in\! A$, 
$\theta_i\!\in\!(\textcolor{red}{B}\textcolor{green}{A})^{l_i}$,
$\textcolor{blue}{\beta_i}\!\in\!\textcolor{blue}{B}$. 
Далее $\E{\theta_{i_1},\ldots,\theta_{i_{l_i}}\!\in\!(\textcolor{red}{B}\textcolor{green}{A})}{\theta_i\!=\!\theta_{i_1}\ldots\theta_{i_{l_i}}}$. 
Наконец, $\theta_{i_j}\!=\!\textcolor{red}{\lambda_{ij}}\textcolor{green}{\mu_{ij}}$, 
где $\textcolor{red}{\lambda_{ij}}\!\in\!\textcolor{red}{B}$, 
$\textcolor{green}{\mu_{ij}}\!\in\!\textcolor{green}{A}$.
Имеем $\alpha\!=\!\gamma_1\textcolor{red}{\lambda_{11}}\textcolor{green}{\mu_{11}}\ldots\textcolor{red}{\lambda_{1l_1}}\textcolor{green}{\mu_{1l_1}}\textcolor{blue}{\beta_1}\gamma_2\ldots\textcolor{blue}{\beta_{k-1}}\gamma_k\textcolor{red}{\lambda_{k1}}\textcolor{green}{\mu_{k1}}\ldots\textcolor{red}{\lambda_{kl_k}}\textcolor{green}{\mu_{kl_k}}\textcolor{blue}{\beta_k}$, где $\gamma,\textcolor{green}{\mu}\!\in\! A$, 
$\textcolor{red}{\lambda},\textcolor{blue}{\beta}\in B$. Получили $\alpha\in A(BA)^*B$.
 \end{example}
 Особенность этого примера: мы построили такой язык $L$, что $L^+=L$. 
\end{frame}

\subsubsection{1.2.3. Языковые графы }
\label{1.2.3.}
\begin{frame}
\frametitle{1.2.3. Языковые графы (1/5)}
Для описания языка можно использовать граф. 
\odmkey{Граф языка над
алфавитом $A$}{Граф!языка} (или \odmkey{источник}{Источник})~---
это ориентированный (мульти)граф, на дугах которого написаны буквы (или пустое слово
$\varepsilon$) и в котором выделены
\odmkey{начальные}{Узел!начальный} и
\odmkey{заключительные}{Узел!заключительный} узлы.

\smallskip
Далее
используются следующие обозначения. Если $s$ и $t$~--- узлы
источника $H$, то $\langle s, t\rangle$ обозначает путь из $s$ в
$t$, а $\alpha= V(\langle s, t\rangle)$ означает слово, полученное
конкатенацией букв, написанных на дугах пути $\langle s,
t\rangle$. Такое слово $\alpha \in A^*$ будем называть
\odmkey{значением}{Значение!пути} пути. Очевидно, что $|V(\langle
s, t \rangle)|\leq |\langle s, t\rangle|$  и $|V(\langle s, t
\rangle)| = |\langle s, t \rangle|$ тогда и только тогда, когда
$\langle s, t \rangle$  не содержит дуг, нагруженных~$\varepsilon$.

\smallskip
Говорят, что источник {\em представляет} некоторый язык, а именно,
язык, состоящий из слов, которые являются значениями всех путей,
ведущих из начальных узлов в заключительные.
Очевидно, что всякий источник представляет некоторый язык (но не
всякий язык может быть представлен графом).

\begin{remark}
Множества начальных и заключительных узлов могут пересекаться и
даже совпадать. На диаграммах начальные узлы помечают входной стрелкой
<<из ниоткуда>>, а заключительные узлы выделяют жирной границей. Мультидугу обычно изображают одной линией и указывают
сразу множество букв через запятую.
\end{remark}
\end{frame}

\begin{frame}
\frametitle{1.2.3. Языковые графы (2/5)}

\begin{columns}[t]
\begin{column}[T]{8.5cm}
\begin{example}
Справа представлена диаграмма источника
десятичных записей натуральных чисел.
\end{example}
\end{column}
\begin{column}[T]{6.5cm}
\includegraphics[width=6.5cm]{1_2_3_1.jpg}
\end{column}
\end{columns}

Наряду с путями, которые ведут из начальных
узлов в заключительные узлы 
и представляют <<правильные>> слова языка,
в источнике могут быть пути,
которые не ведут из начальных узлов 
в заключительные, и тем самым представляют
<<неправильные>>  слова, то есть 
не принадлежащие языку.
Источник, в котором существуют пути для 
{\em всех} слов в алфавите, естественно назвать
\odmkey{полным}{Источник!полный}. 
Если два источника представляют один язык, то их
называют \odmkey{эквивалентными}{Источник!эквивалентный}.

\begin{theorem}
Для любого источника существует 
эквивалентный полный источник.
\end{theorem}

\begin{proof}
Если в  графе из любого узла исходят дуги, помеченные всеми буквами алфавита,
то источник, очевидно, полный.
Любой источник можно превратить в полный
следующим образом.
Нужно добавить ещё один узел <<поглотитель>> и
назначить его не заключительным. 
Для каждого из старых узлов,
в которых не все буквы алфавита использованы
на исходящих дугах,
нужно провести новые дуги в узел поглотитель и пометить их недостающими
буквами.
В самом этом узле поглотителе нужно добавить петли для всех букв.
Источник станет полным, а представляемый язык не изменится.      
\end{proof}

\end{frame}


\begin{frame}
\frametitle{1.2.3. Языковые графы (3/5)}

\odmkey{Размером}{Размер!источника}
(\odmkey{величиной}{Величина!источника}) источника $H$ называется
количество узлов в нем (обозначение $|H|$).

\medskip
Подмножество $X\subset H$ узлов источника называется
\odmkey{замкнутым}{Подмножество!замкнутое}, если
$$
x\in X\And x
\stackrel{\varepsilon}\longrightarrow y \Impl y\in X
$$
(то~есть
в замкнутое подмножество входят все узлы, в которые можно попасть
<<даром>>~--- по слову $\varepsilon$).

\begin{remark}
Если в источнике нет пустых дуг, то все подмножества узлов
замкнутые.
\end{remark}

\medskip
Языковой граф (источник) называется
\odmkey{двухполюсным}{Граф!двухполюсный} или
\odmkey{двухполюсником}{Граф!двухполюсник}, если он имеет один
начальный и один заключительный узел. Обозначим их $s$ и $t$
соответственно.

\begin{theorem}
Для любого источника существует эквивалентный ему двухполюсник.
\end{theorem}

\begin{columns}[t]
\begin{column}[T]{9cm}
\begin{proof}
Пусть $H$~--- заданный источник, $S$~--- множество начальных
узлов, $T$~--- множество заключительных узлов. Построим
эквивалентный двухполюсный источник так, как показано на
рисунке.
\end{proof}
\end{column}
\begin{column}[T]{6cm}

\vspace{-12pt}
\begin{figure}[h]
\begin{center}
\includegraphics[width=6cm]{1_2_3.jpg}
\end{center}
\end{figure}
\end{column}
\end{columns}

\end{frame}

\begin{frame}
\frametitle{1.2.3. Языковые графы (4/5)}
Введенное понятие двухполюсника исчерпывающим образом
характеризует множество языков, описываемых регулярными
выражениями.

\begin{theorem}
Для любого регулярного выражения $E$ существует двухполюсник $H$,
представляющий $E$.
\end{theorem}
\begin{proof}
Построим $H$ индукцией по структуре $E$.

Буквы: для каждой буквы  $a$, входящей в выражение $E$, построим
источник, как показано на рисунке~{\em а}.
Альтернация: построим параллельное соединение источников,  как
показано на рисунке~{\em б}.
Конкатенация: построим последовательное соединение источников,  как
показано на рисунке~{\em в}.
Итерация: построим зацикливание источников,  как показано на
рисунке~{\em г}.
\end{proof}

\begin{figure}[h]
\begin{center}
\includegraphics[width=15cm]{1_2_4_1.jpg}
\end{center}
\end{figure}

\end{frame}

\begin{frame}
\frametitle{1.2.3. Языковые графы (5/5)}
\begin{remark}
Пустые дуги нужны для наглядности построения и простоты
доказательства. При проведении реального построения графа по
регулярному выражению, как правило, можно обойтись склейкой узлов,
как показано на рисунке. Склеенные узлы и добавленные дуги выделены красным.
\end{remark}

\begin{figure}[h]
\begin{center}
\includegraphics[width=140mm]{1_2_4_2.jpg}
\end{center}
\end{figure}

\begin{example}
На рисунке показан двухполюсник,
представляющий десятичные числа со знаком.

\vspace{-8pt}
\begin{figure}[h]
\begin{center}
\includegraphics[width=140mm]{1_2_3_2.png}
\end{center}
\end{figure}

\end{example}


\end{frame}

\subsubsection{1.2.4. Детерминированные и недетерминированные источники}\label{1.2.4.}
\begin{frame}
\frametitle{1.2.4. Детерминированные и недетерминированные источники (1/5)}

Источник (граф языка) называется
\odmkey{детерминированным}{источник!детерминированный}, если он
содержит одну начальную вершину, не содержит пустых дуг и
все дуги, исходящие из одного узла, помечены разными буквами.
В противном случае,
то есть если нарушено хотя бы одно из
трёх условий, источник называется
\odmkey{недетерминированным}{источник!недетерминированный}.


\begin{example}
На рисунке слева показан детерминированный источник,
представляющий общепринятый язык записи 
вещественных чисел с фиксированной запятой,
а справа показан эквивалентный недетереминированный источник,
в котором нарушены все три условия детерминированности.


\begin{center}
\includegraphics[width=12cm]{1_2_5_1.jpg}
\end{center}


\end{example}

\end{frame}


\begin{frame}
\frametitle{1.2.4. Детерминированные и недетерминированные источники (2/5)}

\begin{theorem} {\NB{[Детерминизации]}}
Для любого  источника $H$ с $n$ узлами существует эквивалентный
ему детерминированный источник $H'$, имеющий не более $2^n$ узлов.
\end{theorem}
\begin{proof}
Пусть $\bar{H}$~--- множество узлов исходного источника $H$,
причем $S\subset\bar{H}$~--- множество начальных узлов,
а $T\subset\bar{H}$~--- множество заключительных узлов.
Построим источник $H'$. Рассмотрим все замкнутые подмножества
$\bar{H}$ (включая пустое). Обозначим их $\bar{H}'\Assign 2^{[\bar{H}]}$
и будем считать их узлами источника $H'$.  Таких узлов не более
$2^n$, $|H'|\leq 2^n$. Положим $S'\in \bar{H}'$~--- наименьшее замкнутое подмножество
$\bar{H}$, содержащее все начальные узлы $H$,
$S'\Assign[S]$. Все элементы $\bar{H}'$,
содержащие хотя бы один заключительный узел $\bar{H}$, объявим
заключительными и обозначим 
$T'\Assign\Set{q'\in\bar{H}'}{\E{q\in\bar{H}}{q\in q'\And q\in T}}$.

Построим теперь дуги источника $H'$. Рассмотрим $q_i'\in \bar{H}'$
($q_i'\subset \bar{H}$) и символ $a$.
Положим $q_j'\Assign \Set{q_j\in \bar{H}}{
\E{q_i\in q_i'}{\E{ \langle q_i, q_j\rangle}{ V(\langle q_i,
q_j\rangle)=a}}}$~--- множество узлов, в которые ведут пути из
$q_i'$ со значением $a$. Если $q_j'$ не замкнуто, то положим
$q_j'\Assign [q_j']$. Если нет путей со значением $a$, исходящих из
узлов множества $q_i'$,  то положим $q_j'\Assign \emptyset$, $\emptyset\in \bar{H}'$.
Множество $q_j'$ является узлом источника $H'$, $q_j'\in\bar{H}'$.
Проведем дугу $q_i'\stackrel{a}\longrightarrow q_j'$.
Таким
образом, каждому $q_i'$ и каждому $a$ соответствует ровно одна
дуга, так что источник $H'$ детерминированный по построению.
\end{proof}
\end{frame}


\begin{frame}
\frametitle{1.2.4. Детерминированные и недетерминированные источники (3/5)}
\begin{proof} (Продолжение)
Осталось показать, что $H\sim H'$. Для этого сначала покажем, что
$
\E{\langle S', q_j'\rangle}
{V(\langle S', q_j'\rangle) = \alpha}
\Equiv
\A{q_j\in q_j'}
{\E{q_i\in S}
{V(\langle q_i, q_j \rangle)=\alpha}}.
$
Индукция по длине $\alpha$. База: если
$\alpha=\varepsilon$, то по построению $q_j' = S'$, так~как в $H'$
нет пустых дуг. Если $\alpha=a$,
то, база индукции верна по построению $H'$.
Индукционный переход: пусть
доказываемое утверждение верно для $|\alpha|\leq k$. Рассмотрим
$\alpha a$.

Пусть $\E{\langle S', q_j'\rangle}
{V(\langle S', q_j'\rangle)= \alpha a}$.
Обозначим
$q_\ell'$~--- предпоследний узел на этом пути, так что
$V(\langle S', q_\ell'\rangle)= \alpha$ и
$q_\ell'\stackrel{a}\longrightarrow q_j'$.
По индукционному предположению в источнике $H$
$\A{q_\ell\in q_\ell'}{\E{q_i\in S}{V(\langle q_i, q_\ell \rangle)=\alpha}}$.
По построению источника $H'$, если
$V(\langle q_\ell', q_j'\rangle) =a$, то
$\A{q_j\in q_j'}{\E{q_\ell\in q_\ell'}{V(\langle q_\ell, q_j\rangle)=a}}$ и, таким
образом, в источнике $H$ имеем путь $\langle q_i,\dots, q_\ell,
q_j\rangle$, причём $V(\langle q_i, q_\ell\rangle)=\alpha$ и
$V(\langle q_\ell, q_j\rangle)=a$.
Пусть теперь \\$\A{q_j\in q_j'}
{\E{q_i\in S}
{V(\langle q_i, q_j \rangle)=\alpha a}}$. Рассмотрим
множество $q_\ell'$ узлов, в которые можно попасть из $q_i$ по путям
со значением $\alpha$. Аналогичным образом имеем путь
$\langle S',\dots, q_\ell', q_j'\rangle$, причём
$V(\langle S',q_\ell'\rangle)=\alpha$ и $V(\langle q_\ell', q_j'\rangle)=a$.
Из доказанного свойства и определения заключительных узлов следует, что в источнике $H'$ путь $\alpha$  из начального узла $S'$ в какой-либо
заключительный узел существует тогда и только тогда, когда в источнике $H$ существует
путь $\alpha$ из некоторого начального узла в какой-либо
заключительный узел.
\end{proof}
\end{frame}


\begin{frame}
\frametitle{1.2.4. Детерминированные и недетерминированные источники (4/5)}
\begin{remark}
Если интерпретировать источник как алгоритм
порождения слов языка, то теорема детерминизации показывает,
что в некоторых случаях по недетерминированному алгоритму
можно построить функционально эквивалентный детерминированный.
\end{remark}

\begin{example}
В доказательстве теоремы используется построение, при котором
количество узлов экспоненциально возрастает. Однако это не
означает, что детерминированный источник всегда имеет больше
узлов, чем недетерминированный. На рисунке приведён
соответствующий пример.

\vspace{-4pt}
\begin{figure}[h]
\begin{center}
\includegraphics[width=110mm]{1_2_5_4.jpg}
\end{center}
\end{figure}
\end{example}


\end{frame}


\begin{frame}
\frametitle{1.2.4. Детерминированные и недетерминированные источники (5/5)}

\begin{remark}
Детерминированный источник~--- это граф автомата без выходов
(см. \DMref{п.~2.2.1}{2.2.1.}).
Недетерминированный источник~--- это недетерминированный конечный автомат,
то~есть автомат, в котором нарушено условие
автоматной детерминированности (см. 
\DMref{п.~2.1.4}{2.1.4.}).
Если интерпретировать источник как алгоритм
порождения слов языка, то недетерминированный источник
на каждом
шаге  <<выбирает>> одно из ребер, по которому можно
перейти. Теорема детерминизации показывает,
что с точки зрения порождения языков, возможности 
этих двух классов автоматов совпадают,
хотя эти классы автоматов различны с точки зрения
программной реализации.
\end{remark}

\bigskip
\TODO{ Вопрос о программной реализации детерминированных и детерминированных 
источников обозначен, но не закончен, это открытый TODO. Нужно либо подготовить хороший атом с примером, либо перенести этот слайд в следующий семестр.}

\begin{remark}
Свойства источника, как графа, связаны со свойствами языка, как множества слов.
Например, источник конечного языка всегда можно представить в виде
леса корневых деревьев, то есть бесконтурного орграфа. 
Источник бесконечного языка, напротив, обязательно содержит хотя бы один контур.  
\end{remark}

\end{frame}

\subsubsection{1.2.5. Расширения языка регулярных выражений}
\label{1.2.5.}

\begin{frame}
\frametitle{1.2.5. Расширения языка регулярных выражений (1/4)}

Трёх введённых регулярных операций
(альтернация, конкатенация и замыкание Клини) {\em достаточно} 
для описания класса регулярных языков.
Обнако набор регулярных операций может быть {\em расширен} многими способами,
например, введением дополнительных операций, которые не меняют
выразительной мощности регулярных выражений, но делают их более
{\em удобными}.

\smallskip
Рассмотрим операцию дополнения:
$$\overline{L}\Def
\Set{\alpha\in A^*}{\alpha\not\in L}.$$

\begin{theorem}
Если $L$~--- регулярный язык, то
его дополнение $\overline{L}$ также регулярно.
\end{theorem}
\begin{proof}
Если в полном
источнике языка $L$ заключительные узлы
объявить незаключительными и наоборот, то полученный граф будет
источником языка $\overline{L}$.
\end{proof}

\smallskip
Имея регулярную операцию дополнения, легко
вывести, что операция пересечения также регулярна:
$$
L_1\cap L_2 =
\overline{\overline{L}_1~|~\overline{L}_2}.
$$

\begin{remark}
Алгебра $\langle \cal{R}; \cup, \cap, \overline{\phantom{L}}\rangle$ является булевой алгеброй. 
\end{remark}

\end{frame}


\begin{frame}
\frametitle{1.2.5. Расширения языка регулярных выражений (2/4)}

Помимо введения дополнительных операций над регулярными языками, 
можно вводить сокращённые обозначения для регулярных выражений.

\begin{example}
\odmkey{Символы шаблона}{Символ!шаблона}, которые используются в операционных
системах для задания маски имени файла, имеют следующий смысл:
знак <<?>>, который обозначает один произвольный символ, это сокращение
для регулярного выражения $a_1 | \dots | a_n$,
а знак <<*>>, соответственно, сокращение для
$(a_1 | \dots | a_n)^*$, где
$\Tuple{a}{1}{n}$~--- все символы, которые разрешается использовать
в именах файлов.
\end{example}

\begin{remark}
Важно понимать, что конкретное значение
символов шаблона зависит от контекста,
потому что в разных операционных системах и даже
в разных версиях одной системы наборы допустимых символов различны.
\end{remark}

\smallskip
Пусть далее в этом параграфе $A=\{\Tuple{a}{1}{n}\}$~--- некоторый фиксированный конечный алфавит, причём буквы 
алфавита считаются перенумерованными, а $P$ и $Q$~--- произвольные регулярные 
выражения над этим алфавитом. Между алгеброй регулярных языков и алгеброй 
регулярных выражений установлен изоморфизм, и это означает, что к регулярным выражениям можно применять регулярные операции, получая новые регулярные выражения.        

\end{frame}


\begin{frame}
\frametitle{1.2.5. Расширения языка регулярных выражений (3/4)}
В частности, в соответствии с \DMref{п.~1.2.1}{1.2.1.} и \DMref{п.~1.2.2}{1.2.2.}\\
$ 
PQ \Def\Set{\alpha\beta}{\alpha\in P \And \beta\in Q}, \
P^1\Def P, \ 
P^k\Def PP^{k-1},\  
P^*\Def\bigcup_{i=0}^{\infty}P^i,\ P^+\Def PP^* .$

Введем дополнительное обозначение для опциональной, необязательной конструкции, которое  выражается через базовые регулярные операции:
$$
P? \Def \varepsilon \mid P.
$$

\begin{example}
$(ab)?c = c|abc.
$
\end{example}

\begin{remark}
Операции $P^*$, $P^+$ и $P?$ можно рассматривать как
\odmkey{операции повторения}{Операция!повторения}:\\
$P^*$ --- нуль или более повторений,
$P^+$ --- одно или более повторений,
$P?$ --- нуль или одно повторение $P$.
\end{remark}

Удобно ввести \odmkey{итерацию с разделителем}{Итерация!с разделителем}
(обобщённую итерацию, итерацию Цейтина\footnote{Григорий Самуилович Цейтин (1936--2022)}).
Пусть $P$ и $Q$~--- регулярные выражения:
$$
P \operatorname{\#} Q \Def P(QP)^*.
$$

\end{frame}


\begin{frame}
\frametitle{1.2.5. Расширения языка регулярных выражений (4/4)}
%\begin{example}
%$A\#B = A~|~ABA~|~ABABA~|~ABABABA~|~\dots$
%\end{example}

\vspace{-12pt}
\begin{columns}
\begin{column}{11cm}
\begin{digress}
Выразительную силу и пользу итерации Цейтина легко оценить,
взглянув на источник выражения $P \operatorname{\#} Q$.  
\end{digress}
\end{column}
\begin{column}{4cm}
\begin{center}
\includegraphics[width=2.4cm]{1_2_5_44.jpg}
\end{center}
\end{column}
\end{columns}

\smallskip
Часто используют \odmkey{символьные классы}{Класс!символьный}, то есть
подмножества символов заданного алфавита $A$. 
Подмножество можно задать перечислением элементов или же диапазоном, 
поскольку алфавит линейно упорядочен. 

\begin{example}
$[abcde] = [a-e] = (a|b|c|d|e)$. 
\end{example}

\smallskip
Символьному классу 
можно дать имя и многократно использовать это имя в других регулярных выражениях 

\begin{example} Символьный класс <<цифра>>:
$\backslash d = [0-9] = [0123456789] =  (0|1|2|3|4|5|6|7|8|9)$. 
\end{example}

\smallskip
Используют в том числе и инвертированные символьные классы:
$
[\char`\^ S] = [~\overline{S}~]
$

\begin{remark}
Операция дополнения определена только
относительно фиксированного алфавита.
Следует быть осторожными при использовании 
инвертированных символьных классов, поскольку результат
зависит от ряда трудно учитываемых факторов:
версии операционной системы, 
действующей в данный момент схемы кодирования и т.д.
\end{remark}

\begin{example}
Если $A=\{0,1,2,3\}$, то ${[\char`\^{}01]} = {2|3}$.
\end{example}

\end{frame}


%\begin{frame}
%\frametitle{1.2.5. Расширения языка регулярных выражений (5/5)}
%Ещё два популярных варианта сокращения:
%
%\smallskip
%\begin{itemize}
%\item \odmkey{точка}{Точка (.)} \texttt{.}~--- В одних программах точка является сокращенным обозначением символьного класса, пораждающего элементарный язык, совпадающий с любым элементом некоторого базового языка, а в других– с любым элементом некоторого базового языка, кроме символа новой строки.
%\item \odmkey{сокращённые классы}{Класс!сокращённый} вида
%\texttt{\textbackslash d}, \texttt{\textbackslash w}, \texttt{\textbackslash s}
%(и отрицания \texttt{\textbackslash D}, \texttt{\textbackslash W}, \texttt{\textbackslash S})
%обычно понимаются как сокращения для некоторых символьных классов.
%\end{itemize}
%\end{frame}

\subsubsection{1.2.6. Программирование в регулярных выражениях}
\label{1.2.6.}

\begin{frame}
\frametitle{1.2.6. Программирование в регулярных выражениях (1/5)}
Регулярные выражения являются элегантным и эффективным средством
решения многих задач обработки текстов и встроены в большинство
унаследованных и современных систем программирования.

\bigskip 
\TODO{Посмотреть POSIX, PERL, SNOBOL4, REFAL. Понять основные идеи. \\
Эффективный алгоритм сопоставления с образцом = парсинг с присваиваниями. \\
Обусловленные (условные, с условиями) регулярные выражения. \\
Обязательно хороший пример. Я люблю такой 08.01.2026 => 2026, Jan, 8 \\ 
но можно взять любой, только понятный до конца и без ошибок в коде.\\
Если это пример вдобавок будет выдавать :) в ответ на дату 31.06.2026, то честь и слава. }

\end{frame}

\begin{frame}
\frametitle{1.2.6. Программирование в регулярных выражениях (1/5)}
\TODO{ Замечания ФАН 20 января. \\
В целом хорошо = 4 * 4 = 16. \\
Слить слайды 1 и 2.\\
Расшифровки всех аббревиатур. \\
Извлечение значений для переменных = присваивание переменным значений, это важнейшее место сказано коряво. \\
Backtracing = поиск с возвратами. Есть великолепные графики, показывающие сравнительную производительность юниксоидных регулярных выражений и питона, из которых хорошо видно, что питон = отстой. Это неплохо бы дать в отступлении. \\
В примерах надо объяснять (кратко) все буквы и имена методов}




Регулярные выражения являются элегантным и эффективным средством
решения многих задач обработки текстов и встроены в большинство
унаследованных и современных систем программирования.
\end{frame}


% Текст слайда подготовил Грушин Андрей 13.01.2026
\begin{frame}
\frametitle{1.2.6. Программирование в регулярных выражениях (2/5)}
Исторически и практически регулярные выражения развивались в разных традициях.

\smallskip
\begin{itemize}
\item \odmkey{POSIX}{POSIX} (BRE/ERE) --- стандарт для многих UNIX-утилит
(\texttt{grep}, \texttt{sed}, \texttt{awk}):
акцент на переносимости и предсказуемости.

\smallskip
\item \odmkey{PERL/PCRE}{PERL@PERL/PCRE} --- <<расширенные>> регулярные выражения: захваты, подстановки, дополнительные проверки контекста, обратные ссылки и др.

\smallskip
\item \odmkey{SNOBOL4}{SNOBOL4} --- сопоставление с образцом является центральной
идеей языка: образцы можно комбинировать, успешное сопоставление
естественно связывает переменные с найденными фрагментами.

\smallskip
\item \odmkey{REFAL}{REFAL} --- сопоставление и переписывание выражений по правилам: акцент не на поиске в строке, а на систематических преобразованиях структур.
\end{itemize}
\end{frame}


% Текст слайда подготовил Грушин Андрей 13.01.2026
\begin{frame}
\frametitle{1.2.6. Программирование в регулярных выражениях (3/5)}
В программировании важная задача при работе с регулярными выражениями ---
\odmkey{сопоставление с образцом}{Сопоставление с образцом}:
проверка соответствия строки или её фрагмента некоторому образцу,
и извлечение значений для переменных.

\smallskip
Существует два подхода:
\begin{itemize}
\item Автоматный подход (POSIX): регулярное выражение
преобразуется в автомат, сопоставление выполняется проходом по слову.

\item Backtracking (Perl): перебираются
альтернативы и разбиения, фиксируется первый успешный разбор.
\end{itemize}

\begin{remark}
Если выражение допускает несколько разборов одной строки, то результат
определяется правилами сопоставления в выбранной системе.
\end{remark}

\begin{example}
Выражение \texttt{(a|aa)}, строка \texttt{'aa'}.\\
PERL (backtracking): совпадение \texttt{'a'}.
POSIX (leftmost-longest): совпадение \texttt{'aa'}.
\end{example}
\end{frame}


% Текст слайда подготовил Грушин Андрей 13.01.2026
\begin{frame}
\frametitle{1.2.6. Программирование в регулярных выражениях (4/5)}
В некоторых системах регулярных выражений встречаются
\odmkey{условные регулярные выражения}{Условные регулярные выражения} ---
конструкции, позволяющие выбирать один из образцов в зависимости от условия:
\[
\texttt{(?(cond)\;yes\;|\;no)}.
\]

\begin{remark}
Условные конструкции не входят в POSIX-стандарт.
Они подчёркивают, что на практике регулярные выражения нередко дополняются
средствами управления ходом сопоставления.
\end{remark}

\begin{example}
Требуется распознавать время в одном из форматов: \texttt{HH} или \texttt{HH:MM},
при условии: минуты обязательны, если присутствует двоеточие:
\[
\texttt{(\textbackslash d\textbackslash d)(:)?(?(2)\textbackslash d\textbackslash d)}.
\]
Группа 2 соответствует символу \texttt{':'}.
Условная часть \texttt{(?(2)\textbackslash d\textbackslash d)} требует две цифры минут, если двоеточие было сопоставлено.

\smallskip
Строки \texttt{'09'} и \texttt{'09:30'} подходят, а \texttt{'09:'} и \texttt{'0930'} --- нет.
\end{example}
\end{frame}


% Текст слайда подготовил Грушин Андрей 14.01.2026
\begin{frame}
\frametitle{1.2.6. Программирование в регулярных выражениях (5/5)}
\begin{example}
Строка задаёт дату в формате
\texttt{'DD.MM.YYYY'}. Требуется преобразовать её к виду \texttt{'YYYY, Mon, D'},
при некорректной дате вывести \texttt{':)'}.
\begin{algorithmic}
\STATE \textbf{Input:} $s$

\STATE $R_{\mathrm{valid}} \Assign$ \texttt{\^{}(?:}
\ \texttt{31\textbackslash.(?:01|03|05|07|08|10|12)\textbackslash.\textbackslash d\{4\}}
\STATE \quad \texttt{|30\textbackslash.(?:01|03|04|05|06|07|08|09|10|11|12)\textbackslash.\textbackslash d\{4\}}
\STATE \quad \texttt{|29\textbackslash.(?:01|03|04|05|06|07|08|09|10|11|12)\textbackslash.\textbackslash d\{4\}}
\STATE \quad \texttt{|(?:0[1-9]|1\textbackslash d|2[0-8])\textbackslash.(?:0[1-9]|1[0-2])\textbackslash.\textbackslash d\{4\}}\texttt{)\$}

\STATE $R_{\mathrm{extract}} \Assign$
\texttt{\^{}0?(\textbackslash d+)\textbackslash.0?(\textbackslash d+)\textbackslash.(\textbackslash d\{4\})\$}

\STATE $Mon[1..12] \Assign
[\texttt{Jan},\texttt{Feb},\texttt{Mar},\texttt{Apr},\texttt{May},\texttt{Jun},\texttt{Jul},\texttt{Aug},\texttt{Sep},\texttt{Oct},\texttt{Nov},\texttt{Dec}]$

\IF{$match(R_{\mathrm{valid}}, s)$}
    \STATE $(d, m_{\mathrm{str}}, y)\Assign capture(R_{\mathrm{extract}}, s)$
    \STATE $m\Assign toInt(m_{\mathrm{str}})$
    \STATE \textbf{print} $y\texttt{, }Mon[m]\texttt{, }d$
\ELSE
    \STATE \textbf{print} \texttt{:)}
\ENDIF
\end{algorithmic}
\end{example}
\end{frame}

\subsubsection{1.2.7. Уравнения с регулярными коэффициентами}\label{1.2.7.}
\begin{frame}
\frametitle{1.2.7. Уравнения с регулярными коэффициентами (1/3)}
В алгебре регулярных выражений можно 
составлять и решать {\em уравнения}, пользуясь 
свойствами регулярных операций, установленными
в \DMref{п.~1.2.2}{1.2.2.}.
Коэффициентами и неизвестными таких
уравнений являются регулярные множества 
слов, поэтому их называют
\odmkey{уравнениями с  регулярными коэффициентами}{Уравнение!с регулярными коэффициентами}.
\begin{example}
Уравнение $X=aX+b$ имеет решение $X=a^*b$.\\
Действительно, $X=a^*b=(\varepsilon+aa^*)b
=\varepsilon b +aa^*b = b+aX=aX+b$.
\end{example}
\begin{digress}
Регулярные операции обладают многими 
удобными свойствами, но они не имеют
обратных операций, 
а потому <<школьные>> приёмы
решения алгебраических уравнений,
наподобие правила 
<<перенести неизвестные в левую часть>>
оказываются неприменимыми.
\end{digress}
В уравнении коэффициентами могут 
быть не только отдельные символы или конкретные слова,
как в предыдущем примере, 
но и любые регулярные выражения.
В этом случае решение ищется в форме
регулярного выражения над коэффициентами
уравнения.
\begin{example}
Уравнение $X=EX+B$ имеет решение $X=E^*B$
при любых регулярных коэффициентах $E$ и $B$.
\end{example} 

Однако даже для линейного уравнения решение
может не быть единственным.
\end{frame}

\begin{frame}
\frametitle{1.2.7. Уравнения с регулярными коэффициентами (2/3)}
\begin{example}
Если в уравнении $X=EX+B$
окажется, что $\varepsilon\in E$
(т.~е. $E^+=E^*$), то
выражение $X=E^*(B+C)$ 
является решением для произвольного $C$.  
Действительно, $EX +B= EE^*(B+C) +B =
E^+(B+C) +B= (E^+B+B)+E^+C =
(E^++\varepsilon)B + E^+C$,
и учитывая, что $E^+=E^++\varepsilon=E^*$  в данном случае, \\имеем 
$(E^++\varepsilon)B + E^+C = E^*B+E^*C=E^*(B+C)=X$.
\end{example}

\medskip
\TODO{
Связь с теоремой о наименьшей неподвижной точке функции (опционально)\\
См. Ахо -- Ульман Т1 стр.126 и далее.\\  
Системы уравнений с регулярными коэффициентами
\\
Развернутый пример :
система уравнений 1(0|1)*\\
то есть, двоичные записи натуральных чисел\\
Очень нужно доделать параграф. }

\end{frame}


% Текст слайда подготовил Зеленов Никита 11.01.2026

\subsubsection{1.2.7. Уравнения с регулярными коэффициентами}\label{1.2.7.}

\begin{frame}
\frametitle{1.2.7. Уравнения с регулярными коэффициентами (1/7)}
В алгебре регулярных выражений можно 
составлять и решать {\em уравнения}, пользуясь 
свойствами регулярных операций, установленными
в \DMref{п.~1.2.2}{1.2.2.}.
Коэффициентами и неизвестными таких
уравнений являются регулярные множества 
слов, поэтому их называют
\odmkey{уравнениями с  регулярными коэффициентами}{Уравнение!с регулярными коэффициентами}.

\begin{example}
Уравнение $X=aX+b$ имеет решение $X=a^*b$.\\
Действительно, $X=a^*b=(\varepsilon+aa^*)b
=\varepsilon b +aa^*b = b+aX=aX+b$.
\end{example}

\begin{digress}
Регулярные операции обладают многими 
удобными свойствами, но они не имеют
обратных операций, 
а потому <<школьные>> приёмы
решения алгебраических уравнений,
наподобие правила 
<<перенести неизвестные в левую часть>>
оказываются неприменимыми.
\end{digress}

В уравнении коэффициентами могут 
быть не только отдельные символы или конкретные слова,
как в предыдущем примере, 
но и любые регулярные выражения.
В этом случае решение ищется в форме
регулярного выражения над коэффициентами
уравнения.

\begin{example}
Уравнение $X=EX+B$ имеет решение $X=E^*B$
при любых регулярных коэффициентах $E$ и $B$.
\end{example} 

Однако даже для линейного уравнения решение
может не быть единственным.
\end{frame}

\begin{frame}
\frametitle{1.2.7. Уравнения с регулярными коэффициентами (2/7)}
\begin{example}
Если в уравнении $X=EX+B$
окажется, что $\varepsilon\in E$
(т.~е. $E^+=E^*$), то
выражение $X=E^*(B+C)$ 
является решением для произвольного $C$.  
Действительно, $EX +B= EE^*(B+C) +B =
E^+(B+C) +B= (E^+B+B)+E^+C =
(E^++\varepsilon)B + E^+C$,
и учитывая, что $E^+=E^++\varepsilon=E^*$  в данном случае, \\имеем 
$(E^++\varepsilon)B + E^+C = E^*B+E^*C=E^*(B+C)=X$.
\end{example}

\begin{remark}
Причина неоднозначности: при $\varepsilon\in E$ уравнение $X=EX+B$
имеет целое семейство неподвижных точек (решений), отличающихся добавкой $E^*C$.
Ниже будет выделено \emph{наименьшее} (по включению) решение.
\end{remark}

\begin{theorem}[Лемма Ардена]
Рассмотрим уравнение $X=EX+B$ над регулярными множествами.
\begin{enumerate}
\item Множество $E^*B$ является решением.
\item $E^*B$ --- \emph{наименьшее} решение по включению.
\item Если $\varepsilon\notin E$, то решение единственно и равно $E^*B$.
\end{enumerate}
\end{theorem}
\end{frame}

\begin{frame}
\frametitle{1.2.7. Уравнения с регулярными коэффициентами (3/7)}


\begin{proof}
(1) уже проверено в примере: $E(E^*B)+B=E^+B+B=(E^++\varepsilon)B=E^*B$.

(2) Пусть $X$ --- любое решение, т.~е. $X=EX+B$.
Тогда $X\supseteq B$.
Подставляя рекуррентно, получаем
$X=EX+B\supseteq E B + B$,
далее $X\supseteq E(EB+B)+B=E^2B+EB+B$, и вообще
$X\supseteq B+EB+\cdots+E^nB$ для всех $n\ge 0$.
Переходя к объединению по $n$, имеем $X\supseteq \bigcup_{n\ge 0}E^nB=E^*B$.

(3) Пусть $\varepsilon\notin E$ и $X$ --- решение.
Покажем $X\subseteq E^*B$.
Возьмём слово $w\in X$ минимальной длины среди слов из $X\setminus E^*B$ (если такие есть).
Из равенства $X=EX+B$ следует: $w\in B$ или $w=uv$, где $u\in E$, $v\in X$.
Если $w\in B$, то $w\in E^0B\subseteq E^*B$ --- противоречие.
Если $w=uv$ с $u\in E$, то поскольку $\varepsilon\notin E$, имеем $|u|\ge 1$, значит $|v|<|w|$.
По минимальности $w$ получаем $v\in E^*B$, следовательно $w=uv\in EE^*B=E^+B\subseteq E^*B$ --- противоречие.
Итак, $X\subseteq E^*B$, вместе с (2) получаем $X=E^*B$.
\end{proof}

\begin{digress}
Связь с теоремой о наименьшей неподвижной точке (ср. Ахо--Ульман, т.~1, стр.~126 и далее):
уравнение $X=EX+B$ можно понимать как задачу найти неподвижную точку функции
$F(X)=EX+B$ на решётке регулярных языков.
\end{digress}
\end{frame}

\begin{frame}
\frametitle{1.2.7. Уравнения с регулярными коэффициентами (4/7)}


\begin{theorem}[Наименьшая неподвижная точка для $F(X)=EX+B$]
Пусть $\mathcal{P}(\Sigma^*)$ упорядочено по включению.
Функция $F(X)=EX+B$ монотонна, и её наименьшая неподвижная точка равна
\[
\mu F=\bigcup_{n\ge 0}F^n(\varnothing)=\bigcup_{n\ge 0}E^nB=E^*B.
\]
\end{theorem}

\begin{proof}
Монотонность: если $X\subseteq Y$, то $EX\subseteq EY$, значит $F(X)\subseteq F(Y)$.

Обозначим $X_0=\varnothing$, $X_{n+1}=F(X_n)$.
Тогда $X_1=B$, $X_2=EB+B$, $X_3=E^2B+EB+B$ и по индукции
$X_n=\bigcup_{i=0}^{n-1}E^iB$.
Следовательно, $\bigcup_{n\ge 0}X_n=\bigcup_{i\ge 0}E^iB=E^*B$.

Проверим неподвижность:
$F(E^*B)=E(E^*B)+B=E^+B+B=E^*B$.
Пусть $X$ --- любая неподвижная точка $F$, т.~е. $X=F(X)$.
Тогда $X$ является решением уравнения $X=EX+B$, а значит по лемме Ардена
$E^*B\subseteq X$.
Итак, $E^*B$ --- наименьшая неподвижная точка.
\end{proof}
\end{frame}

\begin{frame}
\frametitle{1.2.7. Уравнения с регулярными коэффициентами (5/7)}
\odmkey{Система линейных уравнений с рег. коэффициентами}{Уравнение!система с регулярными коэффициентами}
--- это набор уравнений вида
\[
\begin{cases}
X_1 = E_{11}X_1 + E_{12}X_2 + \cdots + E_{1n}X_n + B_1,\\
\ \ \dots\\
X_n = E_{n1}X_1 + E_{n2}X_2 + \cdots + E_{nn}X_n + B_n,
\end{cases}
\]
где все $E_{ij}$ и $B_i$ --- регулярные выражения (языки), а $X_i$ --- неизвестные языки.


\begin{remark}
Удобно писать в векторной форме: $X=EX+B$, где
$X=(X_1,\dots,X_n)^T$, $B=(B_1,\dots,B_n)^T$,
а $E=(E_{ij})$ --- матрица регулярных коэффициентов.
\end{remark}

\begin{theorem}[Существование наименьшего решения системы]
Система $X=EX+B$ имеет наименьшее (по покомпонентному включению) решение,
равное наименьшей неподвижной точке монотонной функции $F(X)=EX+B$.
Это решение получается итерациями Клини:
$X^{(0)}=\varnothing$, $X^{(k+1)}=F(X^{(k)})$, и
$X^{(\infty)}=\bigcup_{k\ge 0}X^{(k)}$.
\end{theorem}

\begin{proof}
Функция $F$ монотонна покомпонентно (как и в одномерном случае).
На полной решётке $\mathcal{P}(\Sigma^*)^n$ существует наименьшая неподвижная точка,
равная объединению итераций от $\varnothing$.
Это объединение даёт вектор языков, который удовлетворяет $X=EX+B$ и минимален среди решений.
\end{proof}
\end{frame}

\begin{frame}
\frametitle{1.2.7. Уравнения с регулярными коэффициентами (6/7)}
\begin{example}[Развёрнутый пример: двоичные записи натуральных чисел]
Рассмотрим язык
\[
L=1(0\mid 1)^*,
\]
то есть все двоичные записи натуральных чисел без ведущих нулей (кроме $0$).
Построим систему уравнений по автомату:

\smallskip

начальное состояние $q_0$ --- прочитали ещё ничего;\\
состояние $q_1$ --- уже прочитали начальную $1$ и затем читаем любые $0/1$; \\
оба состояния принимающие: только $q_1$ (пустая строка после стартовой $1$ допустима).


Обозначим $X_0$ --- язык слов, ведущих из $q_0$ в принимающее состояние,
$X_1$ --- язык слов, ведущих из $q_1$ в принимающее состояние.
Тогда:
\[
\begin{cases}
X_0 = 1X_1,\\
X_1 = 0X_1 + 1X_1 + \varepsilon.
\end{cases}
\]
\end{example}
\end{frame}

\begin{frame}
\frametitle{1.2.7. Уравнения с регулярными коэффициентами (7/7)}

Во втором уравнении сгруппируем:
$X_1=(0+1)X_1+\varepsilon$.
Так как $\varepsilon\notin (0+1)$, по лемме Ардена имеем
\[
X_1=(0+1)^*\varepsilon=(0+1)^*=(0\mid 1)^*.
\]
Подставляя в первое уравнение, получаем
\[
X_0=1X_1=1(0\mid 1)^*.
\]


\begin{remark}
Если нужно включить число $0$, то язык записей натуральных чисел будет
$L_0=\{0\}+1(0\mid 1)^*$.
\end{remark}

\bigskip
\TODO{Неплохо. \\
Немного пример с автоматом недоделан. \\
Обознчения: греческие буквы, названия теорем, сноски ... \\
11 удовлетворительных атомов}

\end{frame}



% Текст слайда подготовил Грушин Андрей 05.01.2026
% поправил ФАН 09.01.2026
\subsubsection{1.2.8. Характеристика класса регулярных языков}\label{1.2.8.}
\begin{frame}
\frametitle{1.2.8. Характеристика класса регулярных языков (1/3) }
Следующее наблюдение помогает очертить класс регулярных языков. 

\begin{lemma} {\NB{[о разрастании]}}
Пусть $L$~--- регулярный язык над алфавитом $\Sigma$.
Тогда существует такое $n\in \Bbb{N}$, что для любого слова $\omega\in L$
длины $|\omega|\ge n$ найдутся слова $\alpha,\beta,\gamma\in\Sigma^*$, 
для которых верно:
$\quad
\alpha\beta\gamma=\omega,\quad \beta\ne\varepsilon,\quad |\alpha\beta|\le n,\quad
\A{i\ge 0}{\alpha \beta^i \gamma\in L} 
$.
\end{lemma}


\begin{proof}
Так как язык $L$ регулярен, существует полный детерминированный источник
$H$ языка $L$ (см.~\DMref{п.~1.2.3}{1.2.3.}, \DMref{п.~1.2.4}{1.2.4.}). Обозначим через $n$ число узлов
источника $H$.

Рассмотрим слово $\omega\in L$, $|\omega|\ge n$, и зафиксируем путь в $H$,
порождающий слово $\omega$.
Пусть $q_k$~--- узел, в котором оказывается источник после порождения первых $k$
символов слова $\omega$ (в частности, $q_0$~--- начальный узел).
Среди узлов $q_0,q_1,\dots,q_n$ имеется совпадение:
$q_k=q_m$ для некоторых $0\le k < m \le n$ (принцип Дирихле).
Положим
$\alpha$~--- префикс слова $\omega$ длины $k$,
$\beta$~--- подслово слова $\omega$, прочитанное на участке пути от $q_k$ до $q_m$,
и $\gamma$~--- оставшийся суффикс, так что $\omega=\alpha\beta\gamma$.
Тогда $|\beta|=m-k>0$ и $|\alpha\beta|= m \le n$.

Поскольку участок пути $q_k\To{\beta}q_m$ является циклом ($q_k=q_m$),
его можно пройти $i$ раз, после чего продолжить путь по $\gamma$.
Следовательно, для любого $i\ge 0$ существует путь со значением $\alpha\beta^i\gamma$
из начального узла в заключительный, то есть $\alpha\beta^i\gamma\in L$.
\end{proof}

\end{frame}

% Текст слайда подготовил Грушин Андрей 06.01.2026
% ФАН 10.01.2026
\begin{frame}
\frametitle{1.2.8. Характеристика класса регулярных языков (2/3) }

\begin{remark}
Лемма о разрастании применима только к {\em бесконечным} языкам,
в которых длина слов не ограничена сверху.
Все конечные языки регулярны.   
\end{remark}

\begin{example}
Язык $D\Assign\Set{a^nb^n}{n\ge 1}$ не является регулярным языком.
Действительно, рассмотрим слово $\omega=a^nb^n \in D$ при некотором 
достаточно большом $n$
и допустим, что оно разбито на три части $\omega=\alpha\beta\gamma$ так,
что $|\beta|>0$. Рассмотрим слово $\alpha\beta\beta\gamma$.
Если $\beta\in a^+$ или $\beta\in b^+$, то
$\alpha\beta\beta\gamma\not\in D$, потому что нарушается баланс $a$ и $b$.
Если же $\beta\in a^+b^+$, то
$\alpha\beta\beta\gamma\not\in D$, потому что возникает чередование $ba$.
\end{example}

\begin{digress}
Рассмотренный пример является частным случаем 
\odmkey{языка Дика}{Язык!Дика}\footnote{Вальтер Франц Антон фон Дик (1856--1934)},
то есть языка правильно построенных скобочных выражений в математике.
Правильно построенные последовательности открывающих и закрывающих скобок
являются контекстно-свободным языком, поскольку порождаются грамматикой
$S\to ()\mid(S)\mid SS$, но не являются регулярным языком, то есть не могут быть 
описаны и/или распознаны никаким регулярным выражением.
В этом проявляется ограниченность регулярных выражений: такие выражения
<<не умеют считать>>, точнее говоря, умеют считать только до числа $n$~--- количество
узлов в графе источнике языка.
\end{digress}
\end{frame}

% Текст слайда подготовил Грушин Андрей 06.01.2026
% ФАН 10.01.2026
\begin{frame}
\frametitle{1.2.8. Характеристика класса регулярных языков (3/3) }

\begin{theorem}
Язык регулярных выражений $\frak{R}$ не является регулярным языком.
\end{theorem}
\begin{proof}
От противного.
Пусть $\Sigma$  
содержит некоторый исходный алфавит 
плюс используемые метасимволы регулярных выражений $(~)~|~*$. По лемме о разрастании существует число $n$,
такое что если $|\omega|\ge n \And \omega\in \frak{R}$, то существует
такое слово $\beta$, $|\beta|>0$, что $\omega=\alpha\beta\gamma$
и $\A{i}{\alpha\beta^i\gamma\in\frak{R}}$, причём $|\alpha\beta|\le n$.
Рассмотрим слово
\[
\omega=\underbrace{( \dots ( \dots (}_{n~\text{символов}~(}\ \cdots\ 
\underbrace{)\dots )\dots)}_{n~\text{символов}~)}.
\]
Это корректное регулярное выражение, содержащее $n$ пар сложенных скобок.
Не имеет значения, что обозначено многоточием $\cdots$. 
Так как первые $n$ символов слова $\omega$ --- это открывающие скобки,
то слово $\beta$ состоит только из открывающих скобок, 
и не содержит закрывающих скобок. 
Но тогда при разрастании слова $\beta$ нарушается баланс скобок и значит
  $\alpha\beta^i\gamma\not\in\frak{R}$ для $i>1$.
  Противоречие, число $n$ не существует. 
\end{proof}

\begin{remark}
Если вложенность скобок ограничена константой, то язык математических выражений
со скобками (язык Дика) оказывается регулярным.
\end{remark}

\end{frame}



\begin{frame}
\frametitle{1.2.А. Онтология формальных языков, грамматик и автоматов}
\end{frame}

\begin{frame}
\frametitle{1.2.Б. Темы для самостоятельной работы студентов}

\begin{enumerate}
\item  \SW{Языки обработки символьной информации} \\
Существуют и используются языки программирования, 
в частности, REFAL, SNOBOL и многие другие, специально 
предназначенные для работы с формальными текстами и
опирающиеся, в частности, на регулярные выражения.
Кроме общих слов, необходим один красивый и убедительный пример 
элегантного решения нетривиальной задачи на обсуждаемом языке.   

\item \SW{Программные инструменты работы с регулярными выражениями} \\ 
В различных системах программирования разные инструменты показывают
большой разброс по производительности
Требуется подготовить обзорный научно-попу\-лярный доклад с графиками 
и рекомендациями.

\item \SW{Эволюция регулярных выражений} \\ 
Предлагаются различные нововведения, которые как расширяют класс 
регулярных выражений с целью охвата новых задач,
так и сужают класс регулярных выражений с целью повышения эффективности.
Требуется выбрать те или иные расширения/сужения
и провести сравнительный анализ и составление с тем, что в учебнике.

\end{enumerate}

\end{frame}



\subsection{1.3. Контекстно-свободные грамматики}\label{1.3.}
\begin{frame}
\frametitle{1.3. Контекстно-свободные грамматики}

Для определения синтаксиса реальных
языков программирования и 
языков предметной области 
чаще всего используются
контекстно-свободные грамматики, (напомним
также введённое обозначение КСГ, \DMref{п.~1.1.3}{1.1.3.}). 
Это обстоятельство объясняется тем,
что выразительной силы регулярных языков иногда оказывается
недостаточно для достижения 
практических целей, а для 
контекстно-зависимых грамматик 
неизвестно достаточно эффективных алгоритмов 
синтаксического анализа.
Таким образом,
КСГ образуют <<золотую середину>>:
эти грамматики позволяют, 
с одной стороны,
формально описать <<почти весь>> 
реальный язык программирования,
и, с другой стороны,
позволяют <<достаточно эффективно>>
реализовать все необходимые 
алгоритмы работы с реальным языком программирования.

\medskip  
Поэтому изучение свойств КСГ 
и алгоритмов, связанных с этим
классом грамматик, необходимо 
для всех последующих построений. 
В этом разделе изложена {\em теория}
 контекстно\-свободных грамматик, 
которая используется в прикладных 
построениях других разделов. 
\end{frame}

\subsubsection{1.3.1. Выводимость в контекстно-свободных грамматиках}\label{1.3.1.}
\begin{frame}
\frametitle{1.3.1. Выводимость в контекстно-свободных грамматиках (1/3)}

Выводимость в КСГ является частным случаем общего
понятия выводимости, рассматриваемого в \DMref{п.~1.1.2}{1.1.2.}. Однако, поскольку форма правил вывода (продукций)
в КСГ ограничена (левая часть правила~--- это всегда одиночный
нетерминал), выводимость в КСГ обладает 
рядом дополнительных полезных свойств и допускает наглядное 
представление в виде {\em упорядоченного корневого дерева},
которое называется  \odmkey{дерево разбора}{Дерево!разбора} или  \odmkey{дерево вывода}{Дерево!вывода} 
(см. \DMref{п.~1.1.5}{1.1.5.}).


\begin{remark}
Для грамматик более общего вида, контекстно-зависимых
грамматик и грамматик с фразовой структурой (\DMref{п.~1.1.3}{1.1.3.}), 
также можно предложить графоподобные представления
выводов в этих грамматиках.
Однако такие представления оказываются
принадлежами более сложным классам объектов
(в частности, гиперорграфам), 
диаграммы которых не столь наглядны, как привычные
диаграммы корневых деревьев.
В то же время, хотя представление вывода деревом разбора
безусловно возможно в случае линейных грамматик
(поскольку в таких грамматиках левая 
часть каждого правила~---  
это одиночный нетерминал, \DMref{п.~1.1.6}{1.1.6.}),
но представление вывода в виде корневого дерева
не добавляет наглядности, поскольку диаграмма 
такого дерева всегда имеет форму 
<<ствола>>, состоящего из одной линии, к которой
подвешены терминалы.  
\end{remark}

\end{frame}

\begin{frame}
\frametitle{1.3.1. Выводимость в контекстно-свободных грамматиках (2/3)}

Напомним, что в корневом ориентированном упорядоченном дереве
\odmkey{сечением}{Сечение!дерева}
называется множество узлов,
которое содержит ровно по одному излу из каждого пути из корня
в лист. 
В частности, сам корень, а также крона, являются сечениями.


\odmkey{Дерево разбора}{Дерево!разбора}, 
или \odmkey{дерево вывода}{Дерево вывода}
в КСГ $G=\langle N, T, R, S\rangle$~---
это ориентированное упорядоченное дерево,
все узлы которого помечены 
символами из $N\cup T \cup \{\varepsilon\}$
таким образом, что если в дереве
какой-то узел помечен символом $A$,
а потомки этого узла помечены символами,
образующими слово 
$\alpha$,
то в множестве правил грамматики $R$
существует правило $A\to\alpha$.


\begin{theorem}
Если $G=\langle N, T, R, S\rangle$~--- контекстно-свободная грамматика, то 
$S\Rightarrow^*\alpha$ тогда и только тогда,
когда существует дерево вывода с корнем $s$ и кроной $\alpha\in T^*$.
\end{theorem}

\begin{proof}
 \TODO{Эссе 1.3. часть 1. Доказательство. По выводу можно построить дерево, а по дереву можно построить вывод. 
}
\end{proof}

\begin{remark} 
Теорема утверждает существование дерева вывода,
но не утверждает, что такое дерево единственно
(см. \DMref{п.~1.3.6}{1.3.6.}).
\end{remark}
\end{frame}

\begin{frame}
\frametitle{1.3.1. Выводимость в контекстно-свободных грамматиках (3/3)}
\TODO{Эссе 1.3. часть 2. Отступление. Аналогия с логическими исчислениями. 
\\Секвенциальная форма = сечение = узел в дереве порождения, см. 1.1.5. 
}

\TODO{Эссе 1.3. часть 3. Отступление. Аналогия с орграфами.  Правило $A\to \alpha B \beta$ порождает дугу $A\to B$. Отношение достижимости соответствует выводимости. }

Таким образом, вопрос о выводимости в КСГ
сводится к вопросу о построении дерева разбора. 

\medskip 
Понятие выводимости в КСГ полезно обобщить, введя выводимость     
из одного символа и из цепочки символов, как терминальных, так и нетерминальных.
Действительно, пусть
задана грамматика $G=\langle \mathrm{M}, \Sigma, \Pi, S \rangle$ и
произвольная цепочка $\alpha \in (\mathrm{M}\cup\Sigma)^*$.
Положим $G' \Assign \langle \mathrm{M}+S', \Sigma, \Pi+(S'\to\alpha), S' \rangle$,
то есть добавим в грамматику новую аксиому $S'$
и новое правило $S' \to\alpha$. Легко видеть, что
$\A{\omega\in\Sigma^*}{S\Rightarrow^*_G \omega \Equiv S'\Rightarrow^*_{G'} \omega}$. 
Это наблюдение позволяет считать, что
$ \alpha\Rightarrow^*_G \omega \Def S'\Rightarrow^*_{G'} \omega$.
Другими словами, любой символ или 
цепочка символов $\alpha$ в грамматике $G$ также порождают некоторый язык, а именно множество выводимых слов, поэтому правомерно использовать обозначение
$\frak{L}(\alpha)$:
$$\frak{L}(\alpha)\Def\Set{\omega\in\Sigma^*}{\alpha\Rightarrow^*_G \omega}.$$
       
\end{frame}


% СЛАЙД 1: ДОКАЗАТЕЛЬСТВО ТЕОРЕМЫ
% Текст слайда подготовил Швачко Никита 5030102/20202
% Дата 11.01.2026
\begin{frame}
\frametitle{1.3.1. Выводимость в КСГ: доказательство (4/7)}

\begin{proof}
\proofsection{Необходимость: по выводу строим дерево.}
Пусть $S = \alpha_0 \Rightarrow \alpha_1 \Rightarrow \ldots \Rightarrow \alpha_n = \alpha$~--- вывод длины $n$.
Построим последовательность деревьев $D_0, D_1, \ldots, D_n$, где крона $D_i$ есть $\alpha_i$.

Положим $D_0$~--- дерево из одной вершины, помеченной $S$.

Пусть $\alpha_i = \beta A \gamma$ и применяется правило $A \to X_1 \ldots X_k$.
Тогда $D_{i+1}$ получается из $D_i$ добавлением к листу, помеченному $A$, 
потомков $X_1, \ldots, X_k$.
Крона $D_{i+1}$ будет $\alpha_{i+1} = \beta X_1 \ldots X_k \gamma$.

Дерево $D_n$~--- искомое: его корень~--- $S$, его крона~--- $\alpha$.

\proofsection{Достаточность: по дереву строим вывод.}
Рассмотрим последовательность сечений $C_0, C_1, \ldots, C_m$, где $C_0 = \{S\}$ (корень), 
а $C_m$~--- крона дерева.
Каждое $C_{i+1}$ получается заменой нетерминала $A \in C_i$ его потомками в дереве.
Последовательность крон сечений образует вывод $S \Rightarrow^* \alpha$.
\end{proof}

\end{frame}



% СЛАЙД 2: ПРИМЕР 1 — Дерево разбора арифметического выражения
% Текст слайда подготовил Швачко Никита 5030102/20202
% Дата 11.01.2026
\begin{frame}
\frametitle{1.3.1. Пример: дерево разбора выражения (5/7)}

\begin{example}
Рассмотрим КСГ арифметических выражений 
$G = \langle N, T, R, E \rangle$, где $N = \{E, T, F\}$, $T = \{a, +, *, (, )\}$, с правилами:
\begin{align*}
(1)\ E &\to E + T, & (2)\ E &\to T, \\
(3)\ T &\to T * F, & (4)\ T &\to F, \\
(5)\ F &\to (E), & (6)\ F &\to a.
\end{align*}

\odmkey{Левосторонний вывод}{Вывод!левосторонний} цепочки $a + a * a$:
\[
E \Rightarrow E + T \Rightarrow T + T \Rightarrow F + T \Rightarrow a + T \Rightarrow a + T * F \Rightarrow a + F * F \Rightarrow a + a * F \Rightarrow a + a * a.
\]

Левый разбор: $\pi = (1, 2, 4, 6, 3, 4, 6, 6)$.

Дерево разбора отражает приоритет: $a + (a * a)$, т.е. умножение выполняется раньше.
\end{example}

\end{frame}



% СЛАЙД 3: ОТСТУПЛЕНИЕ — Аналогия с логикой и орграфами
% Текст слайда подготовил Швачко Никита 5030102/20202
% Дата 11.01.2026
\begin{frame}
\frametitle{1.3.1. Аналогии: логика и орграфы (6/7)}

\begin{digress}[Аналогия с логическими исчислениями]
В секвенциальном исчислении доказательство 
представляется деревом, где каждый узел~--- секвенция (утверждение), 
а рёбра соответствуют правилам вывода.

Аналогия: \emph{секвенциальная форма соответствует \odmkey{сечению дерева разбора}{Сечение!дерева}} 
(см. \DMref{п.~1.1.5}{1.1.5.}).

В обоих случаях корень~--- исходное утверждение (аксиома $S$), 
а листья~--- элементарные факты (терминалы).
\end{digress}

\begin{digress}[Аналогия с ориентированными графами]
Построим граф достижимости нетерминалов: 
для каждого правила $A \to \alpha B \beta$, где $B \in \mathrm{M}$, 
добавим дугу $A \to B$.

Тогда \emph{отношение достижимости} в этом графе соответствует выводимости:
если $B$ достижим из $A$ в графе, то существует $\omega$, такое что 
$A \Rightarrow^* \omega$ и $\omega$ содержит $B$.
\end{digress}

\bigskip
\TODO{Хорошая работа. \\
Семь удовлетворительных атомов. \\
Ошибки форматирования, нет картинок. }

\end{frame}



% СЛАЙД 4: ПРИМЕР 2 — Граф достижимости нетерминалов
% Текст слайда подготовил Швачко Никита 5030102/20202
% Дата 11.01.2026
\begin{frame}
\frametitle{1.3.1. Пример: граф достижимости нетерминалов (7/7)}

\begin{example}
Для грамматики арифметических выражений с правилами
\[
E \to E + T \mid T, \quad T \to T * F \mid F, \quad F \to (E) \mid a
\]
построим граф достижимости нетерминалов.

\NB{Дуги графа:}
\begin{itemize}
\item Из $E \to E + T$: дуги $E \to E$, $E \to T$
\item Из $E \to T$: дуга $E \to T$
\item Из $T \to T * F$: дуги $T \to T$, $T \to F$
\item Из $T \to F$: дуга $T \to F$
\item Из $F \to (E)$: дуга $F \to E$
\end{itemize}

Получаем: $E \leftrightarrow T \leftrightarrow F$, причём $F \to E$.
Все нетерминалы взаимно достижимы (граф сильно связен).
\end{example}

\begin{remark}
Взаимная достижимость означает: из любого нетерминала можно вывести цепочку, 
содержащую любой другой нетерминал.
\end{remark}

\end{frame}


\subsubsection{1.3.2. Преобразования контекстно-свободных грамматик}\label{1.3.2.}
\begin{frame}
\frametitle{1.3.2. Преобразования контекстно-свободных грамматик (1/9)}


Если $\frak{L}(G_1)=\frak{L}(G_2)$, 
то грамматики 
$G_1$ и $G_2$
называются \odmkey{эквивалентными}{Грамматика!эквивалентная} и записывают это так:
$G_1\sim G_2$.

\begin{remark}
Ясно, эквивалентность грамматик 
рефлексивна, симметрична и транзитивна,
а потому введённая терминология и обозначения
вполне оправданы.
\end{remark}

\medskip
\TODO{Эссе 1.3. Часть 4. Пример двух существенно различных, но эквивалентных
грамматик. Пример с арифметическими выражениями.
}

\medskip
Эквивалентные грамматики
порождают один язык, и с этой точки зрения
неразличимы, 
но эквивалентные грамматики
могут иметь различные свойства,
которые при применении, 
например, алгоритмов синтаксического 
анализа, могут давать различные практические результаты.   

\medskip
\TODO{ Эссе 1.3. Часть 4. Ещё пример с арифметическими выражениями.
Разные результаты вычислений, приоритет операций. 
}

Известно несколько частных приёмов,
которые позволяют придать
контекстно-свободной грамматике
желательные свойства, не меняя языка.
Эти приёмы называются
\odmkey{эквивалентными преобразованиями}{Преобразование!эквивалентное!грамматик}
грамматик.
\end{frame}


\begin{frame}
\frametitle{1.3.2. Преобразования контекстно-свободных грамматик (2/9)}

Первым приёмом является устранение бесполезных символов из грамматики.

\smallskip
Символ $X$ называется 
\odmkey{полезным}{Символ!полезный} в грамматике 
$G=\langle{N,T,R,S}\rangle$, если существует некоторый 
\odmkey{терминальный}{Вывод!терминальный}
вывод (по правилам грамматики $G$) вида $S \Rightarrow^* \alpha X \beta \Rightarrow^* \omega$, 
где $\omega \in T^*$.

\begin{remark}
Символ $X$ может быть как терминальным,
так и нетерминальным.
Выводимая цепочка $ \alpha X \beta$ может быть первой или последней в порождении.
\end{remark}    

\smallskip
Символ $X$ называется \odmkey{бесполезным}{Символ!бесполезный}, если он не является полезным.
Бесполезные символы не нужны в грамматике, 
и от них можно и нужно избавиться.
Для выявления бесполезных символов определим свойства, характерные для полезных символов:

\begin{enumerate}
\item Символ $X$ называется  \odmkey{порождающим}{Символ!порождающий}, если $\E{\omega \in T^*}{X \Rightarrow^*\omega}$. 
\item Символ $X$ называется  \odmkey{достижимым}{Символ!достижимый}, если 
$\E{\alpha, \beta\in (N\cup T)^*}{S \Rightarrow^* \alpha X \beta}$ .
\end{enumerate}

Полезный символ является одновременно порождающим и достижимым. Таким образом, если удалить из грамматики все непорождающие и все недостижимые символы, то останутся  полезные символы и только они.


\end{frame}


\begin{frame}
\frametitle{1.3.2. Преобразования контекстно-свободных грамматик (3/9)}

\begin{example}
Рассмотрим грамматику 
$\langle\{S,A\}, \{b\}, \{S \rightarrow A|b,  B \rightarrow b\},S\rangle$.
Все символы, кроме $A$, являются порождающими: $b$ порождает себя, $S$ порождает $b$, $B$ порождает $b$.
Все символы, кроме $B$, являются достижимыми.
\end{example}

\smallskip
Зададим алгоритм индуктивного построения множества порождающих символов.
\begin{enumerate}
\item Каждый терминальный символ порождающий.
\item Если все символы в правой части $\alpha$ правила 
 $A\rightarrow \alpha$ порождающие, то символ $A$~--- порождающий.
\end{enumerate}

\smallskip
Зададим аналогичное индуктивное построение множества достижимых символов.
\begin{enumerate}
\item Начальный символ $S$ достижимый.
\item Если символ $A$ достижимый, то все символы в правой части $\alpha$ правила 
 $A\rightarrow \alpha$ достижимые.
\end{enumerate}

Таким образом, построив множество порождающих символов и
множество достижимых символов, и взяв пересечение этих множеств, получаем множество полезных символов данной грамматики. 
\begin{example}
$\langle\{S,A\}, \{b\}, \{S \rightarrow A|b,  B \rightarrow b\},S\rangle
\sim \langle\{S\}, \{b\}, \{S \rightarrow b\},S\rangle$.
\end{example}

\end{frame}


\begin{frame}
\frametitle{1.3.2. Преобразования контекстно-свободных грамматик (4/9)}

Из определений следует, что  любой {\em терминальный} 
вывод $S \Rightarrow^* \omega$, где $\omega\in T^*$, 
содержит только {\em полезные} символы, а значит бесполезные символы, хотя и могут присутствовать в
грамматике, но не влияют на язык, определяемый грамматикой. 

\begin{theorem}
Если $G_1$~--- произвольная контекстно-свободная 
грамматика, а грамматика $G_2$ получена из грамматики
$G_1$ удалением всех бесполезных символов и всех содержащих их правил, то $\frak{L}(G_1)=\frak{L}(G_2)$. 
\end{theorem}

\begin{remark}
Возможен вырожденный случай, когда начальный 
символ $S$ является бесполезным
в грамматике $G_1$, то есть непорождающим,
и язык, порождаемый грамматикой $G_1$~--- пустое множество. В таком случае 
грамматика $G_2$ оказывается пустой, $\langle\emptyset,
\emptyset, \emptyset, \quad\rangle$,
но заключение теоремы остаётся справедливым,
поскольку язык, определяемый пустой 
грамматикой, естественно считать пустым.
Поэтому, не ограничивая общности,
можно считать, что начальные символы
в обеих грамматиках совпадают.  
\end{remark}

\begin{proof}
Обозначим $L_1\Assign\frak{L}(G_1), L_2\Assign\frak{L}(G_2)$. 
Пусть $\alpha\in L_2$, т.е. $S\Rightarrow^*_{G_2} \alpha$.
Тогда $S\Rightarrow^*_{G_1} \alpha$, поскольку всякое
правило $G_2$ является правилом $G_1$ по построению $G_2$.
Отсюда $\alpha\in L_1$ и $L_2\subset L_1$.
Пусть $\alpha\in L_1$, т.е. $S\Rightarrow^*_{G_1} \alpha$.
Тогда $S\Rightarrow^*_{G_2} \alpha$, поскольку терминальный вывод не может содержать бесполезных символов.
Отсюда $\alpha\in L_2$ и $L_1\subset L_2$.
\end{proof}

\end{frame}


\begin{frame}
\frametitle{1.3.2. Преобразования контекстно-свободных грамматик (5/9)}

Вторым приёмом построения эквивалентных грамматик 
является \odmkey{устранение}{Устранение!правила} 
некоторых правил за счёт {\em подстановки} 
одних правил в другие правила.

\smallskip
\begin{theorem}
Пусть КСГ $G_1=\langle N, T, R, S\rangle$,
$A\in N$, $(A\to \alpha_1 B \alpha_2)\in R$,
$B \to \beta_1 |\dots|\beta_k$.
Построим КСГ $G_2=\langle N, T, R_2, S\rangle$,
где $R_2=R\setminus \{A\to \alpha_1 B \alpha_2\}
\cup\{A\to \alpha_1\beta_1\alpha_2|\dots|\alpha_1\beta_k\alpha_2\}$. 
Тогда $\frak{L}(G_1)=\frak{L}(G_2)$.
\end{theorem}

\begin{proof}
Обозначим $L_1\Assign\frak{L}(G_1), L_2\Assign\frak{L}(G_2)$. 
Пусть $\omega\in L_2$, т.е. $S\Rightarrow^*_{G_2} \omega$.
Тогда заменим в этом выводе 
каждое применение нового правила вида
$A\to\alpha_1\beta_i\alpha_2$ последовательным
применением двух старых правил $A\to\alpha_1 B\alpha_2$, 
$B\to\beta_i$ и получим вывод 
$S\Rightarrow^*_{G_2} \omega$.
Отсюда $\omega\in L_1$ и $L_2\subset L_1$.
Пусть теперь $\omega\in L_1$, т.е. $S\Rightarrow^*_{G_1} \omega$.
Тогда заменим в этом выводе 
каждое применение удаленного правила 
$A\to\alpha_1 B \alpha_2$ и соответствующего
(не обязательного соседнего!) 
применения правила $B\to\beta_i$
применением нового правила $A\to\alpha_1 \beta_i\alpha_2$, 
 и получим вывод 
$S\Rightarrow^*_{G_2} \omega$.
Отсюда $\alpha\in L_2$ и $L_1\subset L_2$.
\end{proof}

\begin{corollary} Пусть в КСГ 
$G_1=\langle N, T, R , S \rangle$
есть некоторый вывод $A\Rightarrow^*_{G_1}\alpha$,
где \\$A\in N, \alpha\in (N\cup T)^*$.
Построим КСГ $G_2=\langle N, T, R_2, S \rangle$,
где $R_2=R\cup\{A\to \alpha\}$.
Тогда $\frak{L}(G_1)=\frak{L}(G_2)$.
\end{corollary}

\begin{remark}
В исчислениях любой вывод можно использовать как 
\odmkey{производное}{Правило!производное} правило.
\end{remark}

\end{frame}


\begin{frame}
\frametitle{1.3.2. Преобразования контекстно-свободных грамматик (6/9)}


Во многих случаях желательно избавляться от 
\odmkey{цепных}{Правило!цепное} правил $A\to B$ 
(см. \DMref{п.~1.1.6.}{1.1.6.}).
Цепные правила можно удалять по одному на основании предыдущей теоремы,
а можно удалить все разом, 
пользуясь следующим индуктивным построением.

\begin{enumerate}
\item Для каждого нетерминала $A$ построим
множество 
$\cal{C}(A)\Assign\Set{B\in N}{A\Rightarrow^* B}$\\
нетерминалов, выводимых из нетерминала $A$ 
цепными правилами из заданного множества правил $R_1$ (при этом $A\in \cal{C}(A)$).
\item Для каждого нетерминала $B$ построим
множество $\cal{B}(B)\Assign\Set{A\in N}{B\in \cal{C}(A)}$\\
нетерминалов, из которых выводим нетерминал $B$ 
цепными правилами из заданного множества правил $R_1$ (при этом $B\in \cal{B}(B)$) 
\item Для каждого {\em нецепного}
правила $A\to\alpha$ из заданного множества правил $R_1$ 
добавим в новое множество правил $R_2$
множество правил вида $B\to \alpha$ для каждого нетерминала 
$B\in\cal{B}(A)$.
\end{enumerate}

\begin{example}
Устранение цепных 
правил в грамматике арифметических выражений. 
\begin{columns}
\begin{column}{3cm}
$E\to E+T\mid T$

$T\to T*F\mid F$

$F\to (E)\mid a$
\end{column}
\begin{column}{6cm}
$\cal{C}(E)\Assign \{E, T, F\},\cal{B}(E)\Assign \{E\}$

$\cal{C}(T)\Assign \{T, F\},\cal{B}(E)\Assign \{E, T\}$

$\cal{C}(E)\Assign \{F\},\cal{B}(E)\Assign \{E, T, F\}$

\end{column}
\begin{column}{6cm}
$E\to E+T\mid T*F \mid (E) \mid a $

$T\to T*F\mid (E) \mid a $

$F\to (E)\mid a$
\end{column}
\end{columns}
\end{example}
\end{frame}


\begin{frame}
\frametitle{1.3.2. Преобразования контекстно-свободных грамматик (7/9)}
\begin{theorem}
Пусть $G_1=\langle N, T, R_1, S\rangle$ 
содержит цепные правила, 
а из $G_2=\langle N, T, R_2, S\rangle$ 
цепные правила удалены указанным построением. 
Тогда $\frak{L}(G_1)=\frak{L}(G_2)$.
\end{theorem}

\begin{proof}
Обозначим $L_1\Assign\frak{L}(G_1), L_2\Assign\frak{L}(G_2)$. 
Пусть $\omega\in L_2$, т.е. $S\Rightarrow^*_{G_2} \omega$.
Тогда если в этом выводе применяется правило
$(B\to\alpha)\in R_2$, то это значит, что  
$\E{A\in G_1}{B\Rightarrow_{G_1}^* A \And (A\to\alpha)\in R}$ (возможно, что $B=A$). 
Заменим правило из $R_2$ последовательностью 
правил из $R_1$ и получим вывод 
$S\Rightarrow^*_{G_1} \omega$.
Отсюда $\omega\in L_1$ и $L_2\subset L_1$.
Пусть теперь $\omega\in L_1$, т.е. $S\Rightarrow^*_{G_1} \omega$.
Тогда заменим в этом выводе 
каждое применение 
последовательности цепных правил, которая заканчивается 
нецепным правилом
$A\to\dots\to B\to\alpha$ применением
правила 
$A\to\alpha$ и 
 и получим вывод 
$S\Rightarrow^*_{G_2} \omega$.
Отсюда $\omega\in L_2$ и $L_1\subset L_2$.
\end{proof}

\medskip
Если в начальном правиле $S \rightarrow \alpha$ аксиома 
$S$ входит в правую часть, $S\in \alpha$ то правило называется \odmkey{$S$-правилом}{Правило!$S$-правилом}.
Покажем, что случай $S$-правил является частным случаем цепных правил, и к нему можно применить алгоритмы избавления от цепных правил.

Действительно, для устранения $S$-правил 
достаточно ввести новое {\em цепное}
правило \\$S_1\to S$ и объявить $S_1$ 
новой аксиомой грамматики. Множество 
выводимых терминальных цепочек 
при этом, очевидно, не изменится.
\end{frame}


\begin{frame}
\frametitle{1.3.2. Преобразования контекстно-свободных грамматик (8/9)}

К нежелательным часто относят $\varepsilon$-правила,
или \odmkey{укорачивающие}{Правило!укорачивающее} правила вида $A\to\varepsilon$.

\begin{example}
Алгоритм устранения цепных правил может оказаться
неприменимым в случае наличия $\varepsilon$-правил.
\end{example}

Устранить укорачивающие правила возможно следующим
построением, во многом следующим идеям предыдущих построений. 
 
\begin{enumerate}
\item Построим
множество 
$\cal{E}\Assign\Set{A\in N}{A\Rightarrow^+ \varepsilon}$
нетерминалов, из которых выводимо
пустое слово.
\item Каждое {\em неукорачивающее} 
правило из множества $R_1$ представим в виде\\ 
$A\to \alpha_0 B_1 \alpha_1 \dots
 \alpha_{k-1}B_k\alpha_k,$
где $\A{i\in 1..k}{B_i\in \cal{E}}$, 
все символы в $\alpha_i$, напротив,  
не входят в $\cal{E}$ и $k\ge 0$.
\item Для каждого неукорачивающего правила,
содержащего $k$ символов из $\cal{E}$,
в множество $R_2$ включаются  $2^k$
правил вида
$A\to \alpha_0 X_1 \alpha_1 \dots \alpha_{k-1}X_k\alpha_k$,
где $X_i$ означает либо ничего, пустое слово, либо
символ $B_i$.
Такие правила соответствуют  
всем возможным подмножествам множества 
$\{\Tuple{B}{1}{k}\}$.
\item Если $S\in\cal{E}$, то добавлятеся правило
$S\to \varepsilon$ для возможности порождения пустого слова.
\end{enumerate}

\end{frame}


\begin{frame}
\frametitle{1.3.2. Преобразования контекстно-свободных грамматик (9/9)}
 
Третий приём, добавление нового нетерминала, очень прост: 
часть цепочки, составляющей правую  часть  правила  исходной грамматики можно обозначить новым нетерминалом, введя дополнительное правило, определяющее  новый  нетерминал.

\begin{example} Пусть в исходной грамматике $G$ одно из  правил  имеет вид $A \rightarrow abBc$.  Обозначив $bB$ через новый  нетерминал  $C$,  заменим это правило двумя:  $A \rightarrow aCc, C \rightarrow bB$.
\end{example}


Грамматика, в которой нет бесполезных символов,
а также $\varepsilon$-правил
(за исключением разве что одного правила 
$S\to\varepsilon)$, $S$-правил и цепных правил 
называется
\odmkey{приведённой}{Грамматика!приведённая}.

\begin{theorem}
Для любой контекстно-свободной грамматики $G_1$
существует такая приведённая контекстно-свободная грамматика $G_2$,
что  $\frak{L}(G_2)=\frak{L}(G_1)$.
\end{theorem}

\begin{proof}
Существование требуемой грамматики
обеспечивается ранее перечисленными алгоритмами
устранения нежелательных элментов.
Важно соблюдать порядок преобразований. 
\begin{enumerate}
\item Устранить $\varepsilon$-правила.
\item Устранить $S$-правила.
\item Устранить цепные правила.
\item Устранить бесполезные символы. 
\end{enumerate}
\end{proof}


\end{frame}

\subsubsection{1.3.3. Нормальная форма Хомского}\label{1.3.3.}
\begin{frame}
\frametitle{1.3.3. Нормальная форма Хомского}
%
%\TODO{Нормальная форма Хомского}
%
%
%\TODO{Грицаенко / Мельник. }
%
%\bigskip
%\TODO{Материал получен давно, $\approx 15$ атомов.
%Возможно, требуется ревизия.
%{\color{blue} Варисов, Самутичев}} 

Говорят, что КСГ находится в 
\odmkey{нормальной форме Хомского}{Нормальная форма!Хомского}, если все правила этой грамматики
имеют одну из двух следующих форм: 
либо $A\to BC$, либо $A\to a$, \\где $A, B, C \in N, a\in T$. 


\begin{theorem}
Для любой контекстно-свободной грамматики $G_1$
существует такая контекст\-но-свободная грамматика $G_2$
в нормальной форме Хомского,
что  $\frak{L}(G_2)=\frak{L}(G_1)-\varepsilon$.
\end{theorem}
\begin{proof}
\TODO{Эссе 1.3. Часть 5. }
\end{proof}

\begin{example}
\TODO{Эссе 1.3. Часть 5. Уупрощенная грамматика выражений в привычной форме\\ и в форме Хомского
}
\end{example}

\medskip
\begin{remark}
Если грамматика имеет нормальную форму Хомского,
то дерево разбора в такой грамматике
является двоичным деревом. Для 
двоичных деревьев известно множество
полезных и эффективных представлений и алгоритмов.
\end{remark}

\end{frame}
\subsubsection{1.3.4. Лемма о разрастании}\label{1.3.3.}
\begin{frame}
\frametitle{1.3.4. Лемма о разрастании}


%\TODO{Материал получен, $\approx 10$ атомов.
%{\color{blue} Василевский, Счастливцев\\Митрофанова / Суриков}} 


\TODO{Эссе 1.3. Часть 6. Лемма о разрастании. (теорема Огдена). \\
Очень важно неформально объяснить, что лемма о разрастании относится к языкам, хотя формулируется и доказывается с помощью грамматик. \\
Можно сначала лемму о разрастании для регулярных языков.\\
Доказательство лучше не прямое а через нормальную форму Хомского. \\
 Обязательно пример: 
как с помощью леммы проверяется, что язык 
контекстно-свободный или нет.}
\end{frame}

% Текст слайда подготовил Гребенкин Егор 
% Дата 10.01.2026
\begin{frame}
\frametitle{1.3.3. Нормальная форма Хомского (2/4)}

\begin{proof}
Пусть $G_1$ — приведённая грамматика (см. \DMref{п.~1.3.2}{1.3.2.}), не содержащая \\$\varepsilon$-правил.
Построим $G_2$ в два этапа.
\begin{enumerate}
    \item \odmkey{Устранение смешанных правил}{Устранение!смешанных правил}. 
    Для каждого терминала $a\in \Sigma$, входящего в правила длины более 1, введём новый нетерминал $X_a$ и правило $X_a \to a$.
    В исходных правилах заменим все вхождения $a$ на $X_a$.
    Теперь правила имеют вид $A \to a$ или $A \to B_1 \dots B_k$, где $B_i \in \mathrm{M}$.
    \item \odmkey{Бинаризация длинных правил}{Бинаризация!правил}.
    Для каждого правила вида $A \to B_1 B_2 \dots B_k$, где $k \ge 3$, введём новые нетерминалы $D_1, \dots, D_{k-2}$ и заменим исходное правило на цепочку:
    $$
    A \to B_1 D_1, \quad D_1 \to B_2 D_2, \quad \dots, \quad D_{k-2} \to B_{k-1} B_k.
    $$
\end{enumerate}
Полученная грамматика $G_2$ содержит только правила вида $A \to BC$ и $A \to a$.
Язык не изменился, так как каждое преобразование является эквивалентным.
\end{proof}
\end{frame}
% Текст слайда подготовил Гребенкин Егор 
% Дата 10.01.2026
\begin{frame}
\frametitle{1.3.3. Нормальная форма Хомского (3/4)}

\begin{example}
Рассмотрим приведённую грамматику арифметических выражений (без скобок):
$E \to E+T \mid T*F \mid a$.
\begin{enumerate}
    \item Устраняем терминалы в длинных правилах. Вводим $P \to +$ и $M \to *$.
    Правила принимают вид:
    $E \to EPT \mid TMF \mid a$.
    \item Разбиваем длинные правила.
    Правило $E \to EPT$ заменяем на $E \to E Z_1, Z_1 \to PT$.
    Правило $E \to TMF$ заменяем на $E \to T Z_2, Z_2 \to MF$.
\end{enumerate}
Результирующая грамматика в НФХ:
$$
\begin{aligned}
E &\to E Z_1 \mid T Z_2 \mid a, \\
Z_1 &\to P T, \quad P \to +, \\
Z_2 &\to M F, \quad M \to *.
\end{aligned}
$$
\end{example}
\end{frame}

% Текст слайда подготовил Гребенкин Егор 
% Дата 10.01.2026
\begin{frame}
\frametitle{1.3.3. Нормальная форма Хомского (4/4)}

\begin{remark}
Важнейшим свойством грамматик в нормальной форме Хомского является структура дерева разбора.
Поскольку каждое правило (кроме правил вида $A\to a$) содержит ровно два нетерминала в правой части, дерево разбора любой цепочки длины $n \ge 2$ является двоичным.
\end{remark}

\begin{corollary}
Если $G$ — грамматика в нормальной форме Хомского, и $S \Rightarrow^* \alpha$, где $\alpha \in \Sigma^*$ и $|\alpha|=n$, то глубина дерева разбора не менее $\log_2 n$.
\end{corollary}

Это свойство ограничивает <<ширину>> дерева относительно его высоты,
 что используется при доказательстве свойств контекстно-свободных языков, в частности, леммы о разрастании.
 
 \bigskip
 \TODO{Необходимо слить два текста в один без повторов и в разумном порядке, распределить по слайдам, ваправить форматы и обозначения.}
\end{frame}

% Текст слайда подготовил Гребенкин Егор 
% Дата 10.01.2026
\subsubsection{1.3.4. Лемма о разрастании (Лемма Огдена)}\label{1.3.4.}
\begin{frame}
\frametitle{1.3.4. Лемма о разрастании (1/4)}

Для контекстно-свободных языков существует аналог леммы о накачке для регулярных языков. Более сильный вариант известен как \odmkey{лемма Огдена}{Лемма!Огдена}.

\begin{theorem}
Для любого контекстно-свободного языка $L$ существует такая константа $n \in \mathbb{N}$, что любую цепочку $z \in L$, в которой отмечено не менее $n$ позиций, можно представить в виде $z = u v w x y$, так что:
\begin{enumerate}
    \item цепочка $vwx$ содержит \textbf{не более} $n$ отмеченных позиций;
    \item цепочка $vx$ содержит \textbf{хотя бы одну} отмеченную позицию;
    \item $\A{i \ge 0}{u v^i w x^i y \in L}$.
\end{enumerate}
\end{theorem}

\begin{remark}
Классическая лемма о разрастании для КС-языков (\odmkey{Bar-Hillel lemma}{Лемма!Бар-Хиллела}) является частным случаем леммы Огдена, где в цепочке отмечены все позиции.
\end{remark}
\end{frame}

% Текст слайда подготовил Гребенкин Егор 
% Дата 10.01.2026
\begin{frame}
\frametitle{1.3.4. Лемма о разрастании (2/4)}

\begin{proof}
Пусть $L=\frak{L}(G)$, где $G=\langle \mathrm{M}, \Sigma, \Pi, S\rangle$ — грамматика в нормальной форме Хомского. Пусть $k=|\mathrm{M}|$. Положим $n = 2^{k+1}$.
Рассмотрим дерево разбора цепочки $z$. Назовём узел дерева \odmkey{ветвящимся}{Узел!ветвящийся}, если оба его поддерева содержат отмеченные позиции.
Построим путь от корня к листу, каждый раз выбирая потомка, содержащего больше отмеченных позиций.
Так как отмеченных позиций $\ge 2^{k+1}$, на этом пути встретится не менее $k+1$ ветвящихся узлов.
\end{proof}
\end{frame}

% Текст слайда подготовил Гребенкин Егор 
% Дата 10.01.2026
\begin{frame}
\frametitle{1.3.4. Лемма о разрастании (3/4)}

\begin{proof} (Окончание)
На выбранном пути находится более $k$ нетерминалов. \\Согласно принципу Дирихле, найдётся нетерминал $A$, который встречается на пути дважды (как $A_1$ и $A_2$, где $A_2$ — потомок $A_1$).
Это позволяет разбить вывод на части:
$$
S \Rightarrow^* u A_1 y, \quad A_1 \Rightarrow^* v A_2 x, \quad A_2 \Rightarrow^* w.
$$
В силу того, что $A_1$ — ветвящийся узел, в поддеревьях, соответствующих $v$ и $x$, содержится как минимум одна отмеченная позиция, следовательно $vx \neq \varepsilon$.
Повторяя вывод $A \Rightarrow^* v A x$ произвольное число раз (или пропуская его), получаем цепочки $u v^i w x^i y$, принадлежащие языку $L$.
\end{proof}
\end{frame}

% Текст слайда подготовил Гребенкин Егор 
% Дата 10.01.2026
\begin{frame}
\frametitle{1.3.4. Лемма о разрастании (4/4)}

\begin{example}
Докажем, что язык $L = \Set{a^m b^m c^m}{m \ge 1}$ не является контекстно-свободным.
Предположим противное. Пусть $n$ — константа леммы Огдена.
Рассмотрим цепочку $z = a^n b^n c^n$. Отметим в ней все позиции с символом $a$.
Разложим $z = u v w x y$.
Согласно лемме, $vx$ обязан содержать отмеченные символы (то есть $a$).
В то же время, $vwx$ содержит не более $n$ отмеченных символов.
Если мы <<накачаем>> цепочку ($i=2$), количество символов $a$ увеличится.
Однако, так как $v$ и $x$ не могут содержать символы $c$ (из-за структуры цепочки и ограничения на длину $vwx$), количество $c$ останется неизменным.
Получим цепочку с разным количеством $a$ и $c$, которая не принадлежит $L$.
Противоречие.
\end{example}

\bigskip
\TODO{Хороший материал, сделано в точности то, что требовалось. \\
Требуются следующие улучшения: \\
Перемещение материала по слайдам с обычной плотностью.\\
Примечание про Огдена и наведение порядка с именованием. \\
Греческие буквы для цепочек и переносы формул. \\
Семь хороших атомов.  }

\end{frame}

\subsubsection{1.3.5. Нормальная форма Грейбах}\label{1.3.5.}
\begin{frame}
\frametitle{1.3.5. Нормальная форма Грейбах (1/2)}
%\TODO{Нормальная форма Грейбах. }
%
%\bigskip
%\TODO{Грицаенко / Мельник. Очень много.  
%Нужна прополка. Можно взять из Ахо второй вариант}
%
%\bigskip
%\TODO{Материал получен давно, $\approx 5$ атомов.
%Возможно, требуется ревизия.
%{\color{blue} Варисов, Самутичев}} 

Говорят, что КСГ находится в 
\odmkey{нормальной форме Грейбах}{Нормальная форма!Грейбах}, если все правила этой грамматики
имеют следующую форму: 
$A\to a\alpha$, где $A \in N, a\in T, \alpha\in N^*$. 

Построение нормальной формы Грейбах 
основано на устранении левой рекурсии.

Нетерминал $A\in N$ называется 
\odmkey{рекурсивным}{Нетерминал!рекурсивный},
если $A\Rightarrow^+\alpha A \beta$. 
Если при этом $\alpha =\varepsilon$, то
нетерминал называется 
\odmkey{леворекурсивным}{Нетерминал!леворекурсивный},
а если $\beta =\varepsilon$, то
нетерминал называется 
\odmkey{праворекурсивным}{Нетерминал!праворекурсивный}.

\medskip
Грамматика, в которой существуют леворекурсивные
(праворекурсивные) нетерминалы,
называется 
\odmkey{леворекурсивной}{Грамматика!леворекурсивная} (соответственно, \odmkey{праворекурсивной}{Грамматика!праворекурсивная}).


\begin{theorem}
Для любой приведённой контекстно-свободной грамматики $G_1$
существует не леворекурсивная контекстно-свободная грамматика $G_2$,
такая что  $\frak{L}(G_2)=\frak{L}(G_1)$.
\end{theorem}

\begin{proof}
\TODO{Эссе 1.3. Часть 7. \\Алгоритм устранения левой рекурсии
\\желательно связать с уравнениями с регулярными коэффициентами 
}
\end{proof}

\begin{example}
\TODO{Эссе 1.3. Часть 7. \\Упрощенная грамматика выражений в привычной форме\\ и с устранением левой рекурсии }
\end{example}

\end{frame}


\begin{frame}
\frametitle{1.3.5. Нормальная форма Грейбах (2/2)}
\begin{theorem}
Для любой не леворекурсивной контекстно-свободной грамматики $G_1$
существует контекстно-свободная грамматика $G_2$ в нормальной форме Грейбах,
такая что  $\frak{L}(G_2)=\frak{L}(G_1)$.
\end{theorem}

\begin{proof}
\TODO{Эссе 1.3. Часть 8. \\ Алгоритм приведения к нормальной форме Грейбах 
}
\end{proof}

\begin{example}
\TODO{Эссе 1.3. Часть 8. Упрощенная грамматика выражений в привычной форме\\ и в нормальной форме Грейбах}
\end{example}

\bigskip
\TODO{{\color{red} Замечания ФАН 20 января}\\
Удовлетворительно, большая часть материала пригодна, но необходимо связать с другими частями и исправить огромное количество ошибок. \\
Существенная дырка: а зачем это вообще нужно?\\
Зеленов, долг = 15 баллов
 }

\end{frame}

% Текст слайда подготовил Зеленов Никита 12.01.2026

\subsubsection{1.3.5. Нормальная форма Грейбах}\label{1.3.5.}

\begin{frame}
\frametitle{1.3.5. Нормальная форма Грейбах (1/6)}

Говорят, что КСГ находится в 
\odmkey{нормальной форме Грейбах}{Нормальная форма!Грейбах}, если все правила этой грамматики
имеют следующую форму: 
$A\to a\alpha$, где $A \in N, a\in T, \alpha\in N^*$. 

Построение нормальной формы Грейбах 
основано на устранении левой рекурсии.

Нетерминал $A\in N$ называется 
\odmkey{рекурсивным}{Нетерминал!рекурсивный},
если $A\Rightarrow^+\alpha A \beta$. 
Если при этом $\alpha =\varepsilon$, то
нетерминал называется 
\odmkey{леворекурсивным}{Нетерминал!леворекурсивный},
а если $\beta =\varepsilon$, то
нетерминал называется 
\odmkey{праворекурсивным}{Нетерминал!праворекурсивный}.

\medskip
Грамматика, в которой существуют леворекурсивные
(праворекурсивные) нетерминалы,
называется 
\odmkey{леворекурсивной}{Грамматика!леворекурсивная} (соответственно, \odmkey{праворекурсивной}{Грамматика!праворекурсивная}).

\begin{theorem}
Для любой приведённой контекстно-свободной грамматики $G_1$
существует не леворекурсивная контекстно-свободная грамматика $G_2$,
такая что  $\frak{L}(G_2)=\frak{L}(G_1)$.
\end{theorem}
\begin{proof}
Доказательство основано на двух стандартных преобразованиях и
на алгоритме последовательного устранения косвенной и непосредственной левой рекурсии.
\end{proof}

\begin{lemma}[Устранение правила]
Пусть в грамматике есть правило $A\to B\alpha$, а все $B$-правила имеют вид
$B\to \beta_1\mid \dots \mid \beta_k$.
Тогда замена правила $A\to B\alpha$ на правила
$A\to \beta_1\alpha \mid \dots \mid \beta_k\alpha$
не меняет язык грамматики.
\end{lemma}
\end{frame}


\begin{frame}
\frametitle{1.3.5. Нормальная форма Грейбах (2/6)}


\begin{proof}
Правило $A\to B\alpha$ допускает ровно те выводы, где сначала выводится некоторая цепочка из $B$,
а затем приписывается $\alpha$.
Замена на $A\to \beta_i\alpha$ непосредственно воспроизводит все эти случаи. 
\end{proof}

\begin{lemma}[Устранение непосредственной левой рекурсии]
Пусть все $A$-правила имеют вид
$A\to A\alpha_1\mid \dots \mid A\alpha_k \mid \beta_1\mid \dots \mid \beta_m$,
где ни одна $\beta_i$ не начинается с $A$.
Тогда при введении нового нетерминала $A'$ и замене этих правил на
\[
A\to \beta_1A' \mid \dots \mid \beta_mA',\qquad
A'\to \alpha_1A' \mid \dots \mid \alpha_kA' \mid \varepsilon
\]
язык грамматики не меняется.
\end{lemma}

\begin{proof}
В исходной грамматике из $A$ выводятся в точности цепочки
$\beta_i\alpha_{j_1}\alpha_{j_2}\cdots \alpha_{j_\ell}$, где $\ell\ge 0$.
В новой грамматике:
$A\Rightarrow \beta_iA'\Rightarrow \beta_i\alpha_{j_1}A'\Rightarrow \cdots
\Rightarrow \beta_i\alpha_{j_1}\cdots \alpha_{j_\ell}A'\Rightarrow
\beta_i\alpha_{j_1}\cdots \alpha_{j_\ell}$.
Обратно, любое применение $A'$ порождает ровно такую же последовательность $\alpha$-блоков
(или $\varepsilon$), значит множества выводимых цепочек совпадают.
\end{proof}

\smallskip



\end{frame}


\begin{frame}
\frametitle{1.3.5. Нормальная форма Грейбах (3/6)}
\begin{remark}
Для устранения \emph{косвенной} левой рекурсии упорядочим нетерминалы
$N=\{A_1,\dots,A_n\}$ и последовательно для $i=1,\dots,n$ заменяем
каждое правило $A_i\to A_j\alpha$ при $j<i$ по лемме об устранении правила,
после чего устраняем непосредственную левую рекурсию для $A_i$
по предыдущей лемме. Тогда в результате не остаётся правил,
из которых может начаться левый вывод $A_i\Rightarrow^+ A_i\gamma$.
\end{remark}
\begin{digress}
Связь с уравнениями: непосредственная левая рекурсия
$A\to A\alpha_1\mid\dots\mid A\alpha_k\mid \beta_1\mid\dots\mid \beta_m$
соответствует ``линейному'' уравнению языков $X=XE+B$,
где $E=\alpha_1+\dots+\alpha_k$, $B=\beta_1+\dots+\beta_m$.
Преобразование с $A'$ реализует фактор $E^*$: из $A'$ выводятся любые повторения $\alpha$-блоков.
\end{digress}

\begin{example}[Упрощённая грамматика выражений и устранение левой рекурсии]
Пусть $T=\{x,+,\ast,(,)\}$.
Обычная (леворекурсивная) грамматика арифметических выражений:
\[
\begin{array}{l}
E\to E+T \mid T,\\
T\to T\ast F \mid F,\\
F\to (E)\mid x.
\end{array}
\]
\end{example}
\end{frame}


\begin{frame}
\frametitle{1.3.5. Нормальная форма Грейбах (4/6)}

Устраняя левую рекурсию, получаем эквивалентную грамматику:
\[
\begin{array}{l}
E\to TE' ,\\
E'\to +TE'\mid \varepsilon,\\
T\to FT',\\
T'\to \ast FT'\mid \varepsilon,\\
F\to (E)\mid x.
\end{array}
\]


\begin{remark}
Полученная грамматика удобна для нисходящего разбора:
ни один нетерминал не является леворекурсивным.
\end{remark}

\begin{theorem}
Для любой не леворекурсивной контекстно-свободной грамматики $G_1$
существует контекстно-свободная грамматика $G_2$ в нормальной форме Грейбах,
такая что  $\frak{L}(G_2)=\frak{L}(G_1)$.
\end{theorem}

\begin{digress}
Обычно перед построением НФГ грамматику упрощают (удаляют бесполезные символы,
$\varepsilon$-правила и цепные правила), а затем приводят к форме,
где каждое правило имеет вид $A\to a$ или $A\to BC$ (при $\varepsilon\notin L$).
Эти преобразования сохраняют язык.
\end{digress}

\end{frame}


\begin{frame}
\frametitle{1.3.5. Нормальная форма Грейбах (5/6)}
\begin{proof}
Опишем конструкцию (для случая $\varepsilon\notin L$), которая после предварительного упрощения
строит НФГ из грамматики, у которой правила имеют вид $A\to a$ и $B\to CD$.

Введём новые нетерминалы $(A\backslash B)$ для всех $A,B\in N$ и положим
$N' = N \cup \{(A\backslash B)\mid A,B\in N\}$.
Смысл: $(C\backslash F)$ порождает такие $\gamma\in N^*$, что $F\Rightarrow^* C\gamma$.

Зададим правила:
\[
(A\backslash A)\to \varepsilon\quad (A\in N),
\]
\[
(C\backslash E)\to A(A\backslash D)(B\backslash E)
\quad \text{для всех }(B\to CD)\in R,\ A\in N,\ E\in N,
\]
\[
A\to a \quad \text{для всех }(A\to a)\in R,
\qquad
S\to a(A\backslash S)\ \text{для всех }(A\to a)\in R.
\]
Теперь заменим в правилах вида $(C\backslash E)\to A(A\backslash D)(B\backslash E)$
первый нетерминал $A$ на терминал $a$, используя все варианты правил $A\to a$.
Получим правила
\[
(C\backslash E)\to a(A\backslash D)(B\backslash E),
\]
то есть каждое правило начинается с терминала, а далее идут только нетерминалы:
получена нормальная форма Грейбах.



\end{proof}
\end{frame}


\begin{frame}
\frametitle{1.3.5. Нормальная форма Грейбах (6/6)}
Корректность (вкратце) доказывается индукцией по длине $\gamma\in N^*$:
\begin{enumerate}
\item Если $F\Rightarrow^* C\gamma$ в исходной грамматике, то $(C\backslash E)\Rightarrow^* \gamma(F\backslash E)$ в новой.
\item Для всех $C,F$ верно: $F\Rightarrow^* C\gamma$ тогда и только тогда, когда $(C\backslash F)\Rightarrow^* \gamma$.
\end{enumerate}
Из (2) следует равенство языков, в частности для стартового символа.
\begin{example}[Упрощённая грамматика выражений в НФГ]
Ниже приведён один из возможных вариантов НФГ для языка выражений из примера (3/6)
(используем вспомогательный нетерминал $R\to )$ для закрывающей скобки).
Все правила имеют вид $A\to a\alpha$.
\small
\[
\begin{array}{l}
R\to ),\\
F\to x \mid (ER,\\
M\to \ast FM \mid \ast F,\\
P\to +TP \mid +T,\\[2mm]
T\to x \mid xM \mid (ER \mid (ERM,\\[2mm]
E\to x \mid xM \mid xP \mid xMP \mid (ER \mid (ERM \mid (ERP \mid (ERMP.
\end{array}
\]
\end{example}

\begin{remark}
Грамматика получена из нелеворекурсивной формы подстановкой правил для $T$ в правила для $E$
так, чтобы в каждой продукции первым символом был терминал.
\end{remark}
\end{frame}




\subsubsection{1.3.6. Алгоритмические свойства грамматик}\label{1.3.6.}
\begin{frame}
\frametitle{1.3.6. Алгоритмические свойства грамматик (1/Х)}

\TODO{Повторное Определение термина неоднозначные грамматики и двусмысленные грамматки.  Два развёрнутых примера. Первый пример Зубковой 1.1.5. Второй пример листочки Чепулис. {\color{blue}  Войнова Алёна}}
\end{frame}


\begin{frame}
\frametitle{1.3.6. Алгоритмические свойства грамматик (2/Х)}


\TODO{Замкнутость, детерминированность, неоднозначность. \\
Теоретический параграф про разрешимое и неразрешимое.\\
Является ли проблема однозначности грамматики алгоритмически разрешимой?}

\begin{theorem}
Проблема распознавания однозначности 
контекстно-свободной грамматики алгоритмически неразрешима.
\end{theorem}

\bigskip
\TODO{\color{blue} Параграф повис в воздухе. {\color{blue} Яков Караев, сделан доклад. Нужны хорошие атомы.}}

\medskip
К сожалению, в общем случае
проблема эквивалентности двух
контекстно-свободных грамматик,
равно как и проблема определения
неоднозначности контекстно-свободной
грамматики оказываются алгоритмически неразрешимыми. 

\medskip
\TODO{разобраться и доложить!}

\medskip
\TODO{Проблемы: Неоднозначность. Пустота языка. }

\medskip
\TODO{{\color{red} Замечание ФАН 21 января.} 
Параграф трудный, но очень важный и полезный для самообразования. Помогу с источниками. }

\end{frame}

\subsubsection{1.3.7. Грамматики $LL(k)$ и $LR(k)$ }\label{1.3.7.}

% Текст слайда подготовила Королева Дарья 5030102/20201 
% Дата 24.12.25
% полностью переписал ФАН 25.12.25
\begin{frame}
\frametitle{1.3.7. Грамматики $LL(k)$ и $LR(k)$ (1/9)}

Как показано в \DMref{п. 1.4.4}{1.4.4.}, 
наилучшие известные алгоритмы синтаксического анализа \\
контекстно-свободных языков без ограничений на грамматику  
имеют асимптотическую оценку трудоёмкости 
разбора $ \cal{O}(n^3)$, где $n$~--- длина входной цепочки.  

\smallskip
В этом параграфе рассматриваются такие подклассы КСГ, 
для которых возможно построить эффективные анализаторы, 
расходующие на обработку входной цепочки длины \(n\) время и память 
порядка $\cal{O}(n)$. 

\smallskip
За эту эффективность приходится расплачиваться тем, что ни один из классов грамматик,
для которых можно построить такие эффективные анализаторы,
{\em не} порождает {\em все} контекстно-свободные языки. 
Однако эти ограниченные классы грамматик оказываются 
{\em достаточными} для описания синтаксиса современных языков программирования.

\smallskip
\begin{digress}
По традиции рассматриваемые здесь подклассы КСГ
принято обозначать $LL(k)$ и $LR(k)$. 
В этих обозначениях первая буква означает направление просмотра
входной цепочки ($L$~--- слева направо), 
вторая буква означает ориентацию вывода
($L$~--- левосторонний вывод, $R$~--- правосторонний вывод, 
\DMref{п.~1.1.5}{1.1.5.}),
а $k$~--- число символов, на которые анализатор должен <<заглянуть вперёд>>
для того, чтобы разбор оставался безвозвратным. 
\end{digress}

\end{frame}

% Текст слайда подготовила Королева Дарья 5030102/20201 
% Дата 24.12.25
% ревизия ФАН 28.12.25
% Изменения внесены 06.01.26 Королева Дарья Сергеевна
\begin{frame}
\frametitle{1.3.7. Грамматики $LL(k)$ и $LR(k)$ (2/9)}
Пусть задана однозначная КСГ $G=\langle \mathrm{M}, \Sigma, \Pi, S \rangle$ и её язык 
$L = \frak{L}(G)$. Рассмотрим функцию, по традиции обозначаемую
$\mathrm{FIRST}^G_k$, доставляющую множество всех 
возможных префиксов терминального порождения слова $\alpha\in(\mathrm{M}\cup\Sigma)^*$:      
$$
    \mathrm{FIRST}_k^G(\alpha) \Def
    \Set{\omega\in\Sigma^*}{(|\omega| < k \And \omega\in\frak{L}(\alpha))
    \Or (|\omega| = k \And \E{\beta\in\Sigma^*}{\omega\beta\in\frak{L}(\alpha) })}.
$$
\begin{remark}
Для доказательства теорем необходимо рассматривать 
функцию $\mathrm{FIRST}^G_k$ как определённую на множестве всех
цепочек, но для реализации простых алгоритмов разбора 
достаточно определить эту функцию только для нетерминальных символов:\\
$\Map{\mathrm{FIRST}^G_k}{\mathrm{M}}{2^{\Sigma^*}}$. 
Если индексы $G$ и $k$ ясны из контекста, их опускают.   
\end{remark}

\smallskip
КСГ $G=\langle \mathrm{M}, \Sigma, \Pi, S \rangle$ 
называется \odmkey{$LL(k)$-грамматикой}{Грамматика!$LL(k)$},
 если из существования двух левых выводов
$ S \Rightarrow_l^* \omega A\alpha \Rightarrow_l \omega\beta\alpha \Rightarrow^* \omega
\delta $ и
$ S \Rightarrow_l^* \omega A\alpha \Rightarrow_l \omega\gamma\alpha \Rightarrow^* \omega\lambda $,
для которых \( \mathrm{FIRST}_k(\delta) = \mathrm{FIRST}_k(\lambda) \), вытекает, 
что \( \beta = \gamma \).

\begin{remark} $LL(k)$ условие означает, что при левом разборе
цепочки \( \omega A\alpha \), где цепочка $\omega$ уже разобрана и $A$ первый слева нетерминал, существует {\em не более одного правила} для \( A \), позволяющего продолжить вывод с данными \( k \) терминалами.
\end{remark}
\end{frame}

% Текст слайда подготовила Королева Дарья 5030102/20201 
% Дата 24.12.25
% Изменения внесены 06.01.26 Королева Дарья Сергеевна
% Изменения внесены 15.02.26 Королева Дарья Сергеевна
\begin{frame}
\frametitle{1.3.7. Грамматики $LL(k)$ и $LR(k)$ (3/9)}
\begin{example}
Рассмотрим грамматику \( G_1 \) с правилами 
\(S \rightarrow aAS \mid b, A \rightarrow a \mid bSA \) 
и покажем, что $G_1$ 
является $LL(1)$-грамматикой. 
Чтобы убедиться в этом, надо рассмотреть две цепочки вывода:
 \[
 S \Rightarrow_l^* \omega S\alpha \Rightarrow_l \omega\beta\alpha \Rightarrow^* \omega\delta \quad
 S \Rightarrow_l^* \omega S\alpha \Rightarrow_l \omega\gamma\alpha \Rightarrow^* \omega\lambda ,\]
где \(\mathrm{FIRST}_1(\delta) = \mathrm{FIRST}_1(\lambda)\).
 Eсли \( \delta \) и \( \lambda \) начинаются символом \( a \), то в выводе участвует \( S \to aAS \) и \( \beta = \gamma = aAS \). Eсли же \( \delta \) и \( \lambda \) начинаются \( b \), то применяется \( S \to b \) и \( \beta = \gamma = b \).
% Откуда $\beta=\gamma$.
\end{example}

\begin{theorem}
Приведенная КСГ $G = \langle \mathrm{M}, \Sigma, \Pi, S \rangle$ является $LL(k)$-грамматикой тогда и только тогда, когда для двух различных правил \( A \to \beta \) и \( A \to \gamma \) пересечение\\ \( \mathrm{FIRST}_k(\beta\alpha) \cap \mathrm{FIRST}_k(\gamma\alpha) = \emptyset \) при всех таких \( wA\alpha \), что \( S \Rightarrow_l^* wA\alpha \).
\end{theorem}
\begin{proof}
\proofsection{Необходимость.}
Допустим, что \( \omega, A, \alpha, \beta \) и \( \gamma \) 
удовлетворяют условиям теоремы, а \\
\( \operatorname{FIRST}_k(\beta\alpha) \cap \operatorname{FIRST}_k(\gamma\alpha) \) 
содержит \( \delta \). 
Тогда по определению \(\operatorname{FIRST}\) для некоторых \( \lambda \) и \( \mu \) найдутся выводы 
\(
S \Rightarrow_i^* \omega A\alpha \Rightarrow_i^* \omega \beta\alpha \Rightarrow_i^* \omega \delta \lambda
\)
и 
\(
S \Rightarrow_i^* \omega A \alpha \Rightarrow_i^* \omega\gamma\alpha \Rightarrow_i^* \omega \delta \mu.
\)
Если \( |\delta| < k \), то \( \lambda = \mu = \varepsilon \). Так как \( \beta \neq \gamma \), то \( G \) не \( LL(k) \)-грамматика.
\end{proof}
\begin{remark}
Здесь использован тот факт, что \( \mathrm{M} \) не содержит бесполезных нетерминалов
\end{remark}
\end{frame}

% Текст слайда подготовила Королева Дарья 5030102/20201 
% Дата 24.12.25
% Изменения внесены 06.01.26 Королева Дарья Сергеевна
% Изменения внесены 15.02.26 Королева Дарья Сергеевна
\begin{frame}
\frametitle{1.3.7. Грамматики $LL(k)$ и $LR(k)$ (4/9)}
\begin{proof} (Окончание)
\proofsection{Достаточность.}
Допустим, что \( G \) не $LL(k)$-грамматика. Тогда найдутся такие два вывода 
\(
S \Rightarrow_i^* \omega A\alpha \Rightarrow_l \omega \beta \alpha \Rightarrow_i^* \omega \delta
\)
и
\(
S \Rightarrow_i^* \omega A\alpha \Rightarrow_l \omega \gamma \alpha \Rightarrow_i^* \omega \lambda,
\)
что цепочки \( \delta \) и \( \lambda \) совпадают в первых \( k \) позициях, но \( \beta \neq \gamma \). Поэтому \( A \to \beta \) и \( A \to \gamma \)~---
 различные правила из \( \Pi \) и каждое из множеств \( \operatorname{FIRST}_k(\beta \alpha) \) и \( \operatorname{FIRST}_k(\gamma \alpha) \) содержит цепочку\\ \( \operatorname{FIRST}_k(\delta) \) = \( \operatorname{FIRST}_k(\lambda) \).
\end{proof}


\begin{example} Грамматика $G$, состоящая из двух правил $S \to aS \mid a$, 
не является $LL(1)$-грамматикой, 
так как $\operatorname{FIRST}_1(aS) = \operatorname{FIRST}_1(a) = a$. 
Интуитивно это можно объяснить так: видя при разборе цепочки, начинающейся символом $a$, только этот первый символ, мы не знаем, какое из правил $S \to aS$ или $S \to a$ надо применить к $S$.

С другой стороны, $G$~--- это $LL(2)$-грамматика. В самом деле, если $S \Rightarrow^*_l \omega A\alpha$, то $A = S$ и $\alpha = \varepsilon$. Так как для $S$ даны только два указанных правила, то $\beta = aS$ и $\gamma = a$. Поскольку $\operatorname{FIRST}_2(aS) = aa$ и $\operatorname{FIRST}_2(a) = a$, то по теореме $G$ будет $LL(2)$-грамматикой.
\end{example} 

\begin{digress} Можно определить функцию 
$\operatorname{FOLLOW}_k(\beta)$, доставляющую
множество первых $k$ терминальных символов, которые могут следовать непосредственно за цепочкой $\beta$ в выводе из начального символа:
$
\operatorname{FOLLOW}_k(\beta) \Def \Set{\omega\in \Sigma^*}{ S \Rightarrow^*
 \alpha \beta \gamma \And \omega \in \operatorname{FIRST}_k(\gamma) }
$.
\end{digress}


\end{frame}

% Текст слайда подготовила Королева Дарья 5030102/20201 
% Дата 15.03.26
\begin{frame}
\frametitle{1.3.7. Грамматики $LL(k)$ и $LR(k)$ (5/9)}
Пусть $\langle \mathrm{M}, \Sigma, \Pi, S \rangle$ -- контекстно-свободная грамматика. 
\odmkey{Пополненной грамматикой}{Грамматика!Пополненная}, полученной из $G$, называется грамматика 
$
G' = \langle \mathrm{M} \cup \{ S' \}, \Sigma, \Pi \cup \{ S' \to S \}, S' \rangle,
$
где $S'$ --- новый начальный символ, не принадлежащий $\mathrm{M}$, $S' \to S$ --- нулевое правило грамматики $G'$, 
а остальные правила занумерованы числами $1, 2, \ldots, p$. 
\begin{remark} Это начальное правило добавляется для того, чтобы свёртку, в которой 
используется нулевое правило, можно было интерпретировать как сигнал о том, 
что цепочка допускается.
\end{remark}

Пусть $\langle \mathrm{M}, \Sigma, \Pi, S \rangle$ --- КСГ, 
а $G' = \langle \mathrm{M}', \Sigma, \Pi', S'\rangle $ --- полученная из неё пополненная грамматика. Тогда 
$G$ называется \odmkey{LR($k$)-грамматикой}{Грамматика!$LR(k)$} для $k \ge 0$, 
если из условий\\
$
S' \Rightarrow^*_{r} \alpha A \omega \Rightarrow_{r} \alpha \beta \omega,\quad
S' \Rightarrow^*_{r} \gamma B \delta \Rightarrow_{r} \gamma \beta \lambda, \quad
\operatorname{FIRST}_k(\omega) = \operatorname{FIRST}_k(\lambda)
$
следует, что $\alpha A \lambda = \gamma B \delta$, то есть $\alpha = \gamma$, $A = B$ и $\delta = \lambda$.
\begin{example}
Грамматика с правилами \(S \to Sa \mid a \) не  является  \(LR(0)\)-грамматикой . 
По определению \(LR(0)\)-грамматики: 
\(
S' \Rightarrow_{r}S, \space
S' \Rightarrow_{r}S \Rightarrow_{r}Sa, \space
\operatorname{FIRST}_0(\varepsilon) = \operatorname{FIRST}_0(a) = \varepsilon
\).
В первом выводе нетерминал $A \Assign S'$, во втором --- нетерминал $B \Assign S$. Tак как $A \neq B$, анализатор, видя в стеке символ $S$, не может детерминированно решить: свернуть его в $S'$ или
оставить его в стеке и ждать символ $a$ для применения правила \(S \to a \).
\end{example}
\end{frame}

% Текст слайда подготовила Королева Дарья 5030102/20201 
% Дата 15.03.26
\begin{frame}
\frametitle{1.3.7. Грамматики $LL(k)$ и $LR(k)$ (6/9)}
Рассмотрим некоторые следствия определения 
$LR(k)$-грамматики, необходимые для конструирования
эффективного синтаксического анализатора.
Пусть $\langle \mathrm{M}, \Sigma, \Pi, S \rangle$ --- некоторая КСГ и
$S\Rightarrow^*_r \alpha A \omega \Rightarrow_r \alpha \beta \omega
\Rightarrow_r^*\gamma\omega \in \Sigma^*$~---
правый терминальный вывод.
В этом случае говорят, что цепочку $\alpha \beta \omega$ 
можно \odmkey{свернуть слева}{Свёртка!слева} по правилу
$A\to\beta$, а вхождение цепочки $\beta$ в цепочку 
$\alpha \beta \omega$ называется \odmkey{основой}{Основа},
а любой префикс цепочки $\alpha\beta$ называется 
\odmkey{активным}{Префикс!активный}.
При этом структура
\(
[A \to \beta_1 \cdot \beta_2,\, \gamma]
\)
называется \odmkey{$LR(k)$-ситуацией}{Ситуация!$LR(k)$},  
если $A \to \beta_1 \beta_2$ --- правило из $\Pi$, 
а $\gamma \in \Sigma^{*}$. 
Говорят, что $LR(k)$-ситуация $[A \to \beta_1 \cdot \beta_2,\, \gamma]$ 
\odmkey{допустима для активного префикса}{Ситуация!Допустимая} $\alpha \beta_1$, если существует вывод
\(
S' \Rightarrow_r^* \alpha A \omega \Rightarrow_r \alpha \beta_1 \beta_2 \omega,
\)
такой, что $\gamma = \operatorname{FIRST}_k(\omega)$.    
\begin{example}
    Рассмотрим грамматику \(G\) с правилами 
    \(S \to C \mid D, C \to aC \mid b, D \to aD\mid c\). 
    Ситуация \([C \to a \cdot C, \varepsilon]\) допустима для \(aaa\), 
    поскольку \(S \Rightarrow_r^* aaC \Rightarrow_r aaaC\), 
    т.~е. в этом примере 
    \(\alpha = aa\) и \(\omega = \varepsilon\)
\end{example} 

\smallskip
Определим функцию $\operatorname{EFF}_k^G(\alpha)$ так:
\begin{enumerate}
    \item если $\alpha$ начинается терминалом, то
    \(
        \mathrm{EFF}_k(\alpha) = \mathrm{FIRST}_k(\alpha);
    \)
    \item если $\alpha$ начинается нетерминалом, то
    \(
        \mathrm{EFF}_k(\alpha) = 
        \{\, \omega \mid 
        \omega \in \mathrm{FIRST}_k(\alpha)\) 
         и существует вывод 
        $\alpha \Rightarrow_r^* \beta \Rightarrow_r \omega \delta$,
        где  $\beta \neq A \omega \delta$  для $ A \in \mathrm{M} \,\}$.
\end{enumerate}

\end{frame}

% Текст слайда подготовила Королева Дарья 5030102/20201 
% Дата 15.03.26
\begin{frame}
\frametitle{1.3.7. Грамматики $LL(k)$ и $LR(k)$ (7/9)}
\begin{theorem}
Пусть $\langle \mathrm{M}, \Sigma, \Pi, S \rangle$ --- КС-грамматика.  
Тогда $G$ является $LR(k)$-грамматикой тогда и только тогда, когда для каждой цепочки $u \in \Sigma^*$ выполняется условие:
если $\alpha \beta$ --- активный префикс правовыводимой цепочки в пополненной грамматике $G'$ и если $LR(k)$-ситуация $[A \to \beta \cdot, \gamma]$ допустима для $\alpha \beta$,
то не существует другой $LR(k)$-ситуации 
$[A_1 \to \beta_1 \cdot \beta_2, \delta],$
допустимой для той же $\alpha \beta$, где $\gamma \in \mathrm{EFF}_k(\beta_2 \delta)$.
\end{theorem}
\TODO{с этого места необходимо продолжить переработку}
\begin{proof}
\proofsection{Необходимость.}
Пусть $[A \to \beta \cdot, u]$ и $[A_1 \to \beta_1 \cdot \beta_2, v]$ —
различные ситуации, допустимые для $\alpha \beta$.  
Тогда существуют выводы:
\[
\begin{aligned}
S' &\Rightarrow_r^* \alpha A \omega \Rightarrow_r \alpha \beta \omega,\\
S' &\Rightarrow_r^* \alpha_1 A_1 \delta \Rightarrow_r \alpha_1 \beta_1 \beta_2 \delta,
\end{aligned}
\]
где $\mathrm{FIRST}_k(\omega)=u$, $\mathrm{FIRST}_k(\delta)=v$ и $\alpha \beta = \alpha_1 \beta_1$. Кроме того, $\beta_2 \delta \Rightarrow^*_r u \lambda$ для некоторой цепочки $\lambda$ и в этом (возможно, пустом) выводе терминал левом конце никогда не заменяется пустой цепочкой. 
\end{proof}
\end{frame}

% Текст слайда подготовила Королева Дарья 5030102/20201 
% Дата 15.03.26
\begin{frame}
\frametitle{1.3.7. Грамматики $LL(k)$ и $LR(k)$ (8/9)}
\begin{proof} (Продолжение)
Рассмотрим 3 случая:
\begin{enumerate}
    \item $\beta_2 = \varepsilon$.
    Тогда $u=v$, и выводы имеют вид
    \(
    S' \Rightarrow_r^* \alpha A \omega \Rightarrow_r \alpha \beta \omega,\)
    \(S' \Rightarrow_r^* \alpha A_1 \delta \Rightarrow_r \alpha \beta_1 \delta.
    \)
    Так как ситуации различны, то либо $A \ne A_1$, либо $\beta \ne \beta_1$,
    и нарушается $LR(k)$-условие.
\item $\beta_2 = \mu$, $\mu \in \Sigma^+$.
Тогда \(
S' \Rightarrow_r^* \alpha A \omega \Rightarrow_r \alpha \beta \omega,\)
\(S' \Rightarrow_r^* \alpha A_1 \delta \Rightarrow_r \alpha \beta_1 \mu \delta.
\)
По условию $u \in \mathrm{EFF}_k(\beta_2 v) = \mathrm{EFF}_k(\mu v)$
Так как $\beta_2 = \mu$ состоит только из терминалов:
$\mathrm{EFF}_k(\mu v) = \mathrm{FIRST}_k(\mu v)$
Таким образом  $u \in \mathrm{FIRST}_k(\mu v)$
Противоречие аналогично
\item $\beta_2$ содержит нетерминал.
Пусть $\beta_2 \Rightarrow_r^* u_1 B u_3 \Rightarrow_r u_1 u_2 u_3$, где $u_1 u_2 \ne \varepsilon$.
Тогда:
\[
\begin{aligned}
S' &\Rightarrow_r^* \alpha A \omega \Rightarrow_r \alpha \beta \omega,\\
S' &\Rightarrow_r^* \alpha_1 A_1 \delta \Rightarrow_r \alpha_1 \beta_1 \beta_2 \delta 
     \Rightarrow_r \alpha_1 \beta_1 u_1 B u_3 \delta \Rightarrow_r \alpha_1 \beta_1 u_1 u_2 u_3 \delta.
\end{aligned}
\]
В $LR(k)$-грамматике требуется, чтобы
\(
\alpha A u_1 u_2 u_3 \delta = \alpha_1 \beta_1 u_1 B u_3 \delta,
\) т.е. \\
\(
\alpha A u_1 u_2 = \alpha_1 \beta_1 u_1 B,
\)
\(
A u_1 u_2 = \beta u_1 B,
\)
но это невозможно, так как $u_1 u_2 \ne \varepsilon$.
\end{enumerate}

\vspace{-12pt}
\end{proof}

\begin{remark}
именно здесь используется условие, что \(u \in \mathrm{EFF}_k(\beta_2, v).\)
Если бы в формулировке использовать только множество $\mathrm{FIRST}_k$,
цепочка $u_1 u_2$ могла бы быть пустой, и доказательство бы не прошло.
\end{remark}

\end{frame}

% Текст слайда подготовила Королева Дарья 5030102/20201 
% Дата 15.03.26
\begin{frame}
\frametitle{1.3.7. Грамматики $LL(k)$ и $LR(k)$ (9/9)}
\begin{proof} (Окончание)
\proofsection{Достаточность.}
Пусть $G$ --- не $LR(k$)-грамматика.  
Тогда существуют два вывода:
\[
\begin{aligned}
\quad & S' \Rightarrow_r^* \alpha A \omega \Rightarrow_r \alpha \beta \omega,\\
\quad & S' \Rightarrow^*_{r} \gamma B \delta \Rightarrow_{r} \gamma \beta \lambda= \alpha \beta \lambda,
\end{aligned}
\]
где $\mathrm{FIRST}_k(\omega) = \mathrm{FIRST}_k(\lambda) = u$,
но $\alpha A \lambda \ne \gamma B \delta$.
Выберем $\alpha \beta$ минимальной длины.  
Пусть $\alpha_1 A_1 y_1$ — последняя открытая цепочка в выводе:
\(
S' \Rightarrow_r^* \gamma B \delta \Rightarrow_r^* \alpha_1 A_1 y_1
\Rightarrow_r \alpha_1 \beta_1 \beta_2 y_1 \Rightarrow_r^* \alpha_1 \beta_1 \lambda.
\)
Тогда $\alpha_1 \beta_1 = \alpha \beta$, и
$u = \mathrm{FIRST}_k(\lambda) \in \mathrm{EFF}_k(\beta_2, \lambda)$. 

Ситуация \([A_1 \to \beta_1 \cdot \beta_2, v]\) допустима для \(\alpha \beta\), где \(v = \mathrm{FIRST}_k(y_1)\).
В то же время ситуация \([A \to \beta \cdot , u]\) допустима для \(\alpha \beta\), поэтому осталось доказать, что \(A_1 \to \beta_1 \cdot \beta_2\) отличается от \(A \to \beta \). 

Пусть \(A_1 \to \beta_1 \cdot \beta_2\) совпадает с \(A \to \beta \). Тогда 
\(
S' \Rightarrow_r^*  \alpha_1 A \lambda
\Rightarrow_r \alpha_1 \beta \lambda
\), где \(\alpha_1 \beta = \alpha \beta \), поэтому \(\alpha_1 = \alpha\) и \(\alpha A \lambda = \alpha B \delta \), что противоречит предположению о том, что \(G\) не $LR(k)$-грамматика 
\end{proof}
\end{frame}

% Текст слайда подготовила Королева Дарья 5030102/20201 
% Дата 15.03.26
\begin{frame}{Специальные классы грамматик I (10/10)}
\begin{remark}
     Можно дать неформальное, но наглядное определение \(LR(k)\)-грамматики в терминах деревьев выводов. Грамматика \(G\) будет \(LR(k)\)-грамматикой, если, просмотрев только часть кроны дерева вывода в этой грамматике, расположенную слева от данной внутренней вершины, и часть кроны, выведенную из нее, а также следующие \(k\) терминальных символов, можно установить, какое правило было применено к этой вершине.
\end{remark}

\TODO{ Этому замечанию надобно найти место где-то в другом месте. Слайд лишний. } 
 
\end{frame}

\subsubsection{1.3.8. Разделённые грамматики}\label{1.3.8.}
% Текст слайда подготовила Королева Дарья 5030102/20201 
% Дата 17.02.26
\begin{frame}
\frametitle{1.3.8. Разделённые грамматики (1/5)}
Введение понятия $LL(1)$-грамматики — это один из возможных способов наложить на грамматику такие условия, чтобы шаг предсказания очередного правила вывода можно было сделать детерминированным, используя только текущий входной символ.

\smallskip
Существует обобщающий подход разбора, соответствующий {\em разделённой} грамматике. 
Для разделенной грамматики шаг предсказания тривиален. В связи с этим современное языкотворчество имеет тенденцию к использованию преимущественно разделённых грамматик,
ввиду простоты их реализации и возможности элементарной проверки условий разделенности по записи правил.
В то время как проверка выполнения свойства 
$LL(k)$ требует применения специального алгоритма.
Более того, проблема определения того, существует ли такое число $k$,
что грамматика является $LL(k)$-грамматикой,
является алгоритмически неразрешимой.  
\end{frame}

\begin{frame}
\frametitle{1.3.8. Вставной слайд в разделённые грамматики}
Этот слайд не является частью материала. УДАЛИТЬ после понимания.

\smallskip
Предыдущий и особенно последующий тексты содержат опечатки,
пропуски определений, ошибки в обозначениях и т.~д.
Всё это надобно устранить обычным порядком.
Здесь дополнительно имеются ошибки в ПОСЛЕДОВАТЕЛЬНОСТИ изложения.
Что такое подход разбора? Откуда взялся шаг предсказания?
Как работает МП-преобразователь?  
Изложение построено верно по существу,
но неверно методически.
ПРИЛЕЖНЫМ студентам (для которых пишется этот учебник!)
затруднено изучение материала.
В этом учебнике алгоритмы анализа
изложены в \DMref{1.4.}{1.4.},
автоматы-распознаватели [должны быть] описаны в главе 2.
Нельзя забегать вперёд без предупреждения: студенты тоже люди.

\smallskip
Как решить проблему?
Первый вариант: насытить текст ссылками вперёд,
убедившись, что ссылки ведут туда, куда надо.

\smallskip
Второй вариант: 
изложить свойства разделённых грамматик чисто грамматически, в частности,
ДОКАЗАТЬ первое замечание и все недоказанные утверждения далее,
 а рассказы про разбор убрать.
Второй подход предпочтельнее, немного труднее,
существенно полезнее для автора (надобно разобраться, в чём дело)
и поэтому намного ЦЕННЕЕ в баллах. 
Хотя количество слайдов СОКРАТИТСЯ, но оценка ВОЗРАСТЁТ! 

\end{frame}

% Текст слайда подготовила Королева Дарья 5030102/20201 
% Дата 17.02.26
\begin{frame}
\frametitle{1.3.8. Разделённые грамматики (2/5)}

КСГ $G=\langle \mathrm{M}, \Sigma, \Pi, S \rangle$ называется \odmkey{простой}{Грамматика!простая}
(или \odmkey{разделённой}{Грамматика!разделённая}),
если для любого нетерминала $A \in\mathrm{M}$
правые части всех правил для этого нетерминала
начинаются с различных терминальных символов.
Другими словами, для любой пары $a\in \Sigma$, $A\in \mathrm{M}$
существует не более одного правила вида $A\to a\alpha$. 
\begin{remark}
Очевидно, что простая грамматика не содержит $\varepsilon$-правил
и является \\$LL(1)$-грамматикой.
\end{remark}
Для разделенной грамматики шаг предсказания максимально прост: если \( a \) — текущий входной символ, \( A \) — верхний символ магазина и в грамматике есть правило \( A \to a \alpha \), то в магазине надо заменить \( A \) на \( \alpha \) и сдвинуть входную головку. 

\smallskip
Обобщающий подход разбора: к \( A \)-правилам разделенной грамматики разрешается добавить \( e \)-правило \( A \to e \), а на шаге предсказания для текущих символов \( a \) и \( A \) разрешается применить правило \( A \to e \) (т.е. удалить \( A \) из магазина), если в грамматике нет правила вида \( A \to a \alpha{} \).

\begin{remark}
 Добавление \( e \)-правил к разделенным грамматикам позволяет выделить новые классы грамматик, позволяющие описывать более сложные языковые конструкции, сохраняя при этом эффективность разбора $\cal{O}(n)$
\end{remark}

\end{frame}

% Текст слайда подготовила Королева Дарья 5030102/20201 
% Дата 17.02.26
\begin{frame}
\frametitle{1.3.8. Разделённые грамматики (3/5)}

КСГ $G=\langle \mathrm{M}, \Sigma, \Pi, S \rangle$ в \odmkey{слабой нормальной форме Грейбах}{Грамматика!в слабой нормальной форме Грейбах} \DMref{п. 1.3.5}{1.3.5.}, у которой правые части правил вида \( A \to a\alpha \) и \( A \to b\beta \) начинаются разными символами, т. е. \( a \neq b \), то она называется \odmkey{псевдоразделенной}{Псевдоразделенная!Грамматика}.

Для формализации выбора правила по текущему символу для псевдоразделенной грамматики строится простой МП-распознаватель \( M_G = \langle \Sigma, \mathrm{M}, \delta, S \rangle \), где функция переходов \( \delta \) определяется следующим образом:
\begin{enumerate}
  \item Если в стеке нетерминал \( A \), а на входе символ \( a \), и существует правило \( A \to a\alpha \), то \( A \) заменяется на \( a\alpha \) в стеке без сдвига входной головки. \(\delta(a, A) = (a, 1)\)
  \item Если правила для символа \( a \) нет, но существует \( e \)-правило \( A \to e \), то нетерминал \( A \) удаляется из стека (также без сдвига головки). \(\delta(b, S) = (e, 0)\)
\end{enumerate}
Автомат служит для детерминизации шага предсказания в левом выводе. Он позволяет сопоставить нетерминал на вершине стека с текущим входным символом без преждевременного сдвига входной головки
\end{frame}

% Текст слайда подготовила Королева Дарья 5030102/20201 
% Дата 17.02.26
\begin{frame}
\frametitle{1.3.8. Разделённые грамматики (4/5)}

\begin{example}
Пусть \( G_1 \) — псевдоразделенная грамматика с правилами
\(
S \to aSA \mid e,  A \to a \mid b.
\)

Тогда \( M_{G_1} = (\{a, b\}, \{S, A\}, \delta, S) \), где
\[
\begin{aligned}
\delta(a, S) &= (SA, 1)  \\
\delta(b, S) &= (e, 0) \\
\delta(\$, S) &= (e, 0) \\
\delta(a, A) &= (e, 1) \\
\delta(b, A) &= (e, 1)
\end{aligned}
\]
\end{example}


\begin{remark}
МП-распознаватель ограничивается анализом одного текущего символа. \\  
Обобщением этого подхода является $k$-предсказывающий алгоритм, где решение принимается на основе аванцепочки длиной $k$. Конфигурация алгоритма представляет собой необработанный вход, содержимое магазина и выходную ленту (номера правил). Работа алгоритма координируется управляющей таблицей, которая фактически реализует функцию \( \delta \) МП-распознавателя. 
\end{remark}
\end{frame}

% Текст слайда подготовила Королева Дарья 5030102/20201 
% Дата 17.02.26
\begin{frame}{1.3.8. Разделённые грамматики (5/5)}
Псевдоразделенная грамматика G называется \odmkey{слаборазделенной}{Слаборазделенная! грамматика}, если $\frak{L}(G) = \frak{L}(M_G)$ для соответствующего простого МП-преобразователя $M_G$

\begin{example}
    Грамматика $G$ с правилами \(S \to a S A \mid e, A \to b \mid e\) слаборазделенная

    Порождаемый язык \(\frak{L} (G) = \{ a^n b^n \mid  0 \leq k \leq n \} \)

    Логика простого МП-распознавателя \( M_G \)

\begin{enumerate}
\item \( \delta(S, a) = (SA, 1) \). При чтении \( n \) символов \( a \) в стеке накапливается \( n \) \( A\).
\item \( \delta(S, b) = (e, 0) \). При встрече первого \( b \) нетерминал \( S \) удаляется из стека без сдвига входа.
\item \( \delta(A, b) = (e, 1) \). Символы \( b \) последовательно удаляют \( A \) из стека.
\item \( \delta(A, \$) = (e, 0) \). Избыточные нетерминалы \( A \) (если \( k < n \)) стираются на маркере конца строки.
\end{enumerate}

\( \frak{L}(G) = \frak{L}(M_G) \), грамматика является слаборазделенной.

\begin{remark}
    Если цепочка начинается с \( b \), стек опустеет раньше времени. Если количество \( b \) превысит \( n \), для лишних символов в стеке не останется нетерминалов \( A\).
\end{remark}
\end{example}
\begin{remark}
    Не существует алгоритма, который по произвольной КСГ определяет, является ли она слаборазделенной.
\end{remark}
\end{frame}



\subsubsection{1.3.9. Классификация контекстно-свободных грамматик}\label{1.3.9.}
\begin{frame}
\frametitle{1.3.9. Классификация контекстно-свободных грамматик}
\TODO{Итоговый обзор на тему: в каких случаях какие грамматики целесообразно применять. 
Здесь же про грамматики в регулярной форме.}

\TODO{Общая диаграмма онтология всех классов грамматик. 
Здесь же про класс КС-языков, что 
язык $\{a^n b^n c^n\}$ НС-язык, но не КС-язык,
и при этом $\{a^n b^n\}$ КС-язык, но 
не регулярный язык.}

\bigskip
\TODO{\color{blue} Параграф повис в воздухе. Кто???}
\end{frame}

% Текст слайда подготовил Швачко Никита 5030102/20202
% Дата 11.01.2026
\begin{frame}
\frametitle{1.3.9. Классификация контекстно-свободных грамматик (2/5)}

В курсе рассмотрены различные классы грамматик. 
Выбор подходящего класса зависит от задачи:

\medskip
\begin{tabular}{p{45mm}|p{100mm}}
\textbf{Класс грамматики} & \textbf{Когда применять} \\
\hline
\odmkey{Регулярные грамматики}{Грамматика!регулярная} & 
Лексический анализ, токенизация, поиск по образцу. 
Разбор за $\mathcal{O}(n)$, минимальные ресурсы. \\
\hline
\odmkey{$LL(k)$-грамматики}{Грамматика!$LL(k)$} & 
Ручное написание парсера (рекурсивный спуск). 
Простая реализация, понятная структура. \\
\hline
\odmkey{$LR(k)$-грамматики}{LR(k)!Грамматика} & 
Автоматическая генерация парсера (yacc, bison). 
Мощнее $LL(k)$, но сложнее для понимания. \\
\hline
Общие КСГ & 
Когда важна естественность описания, а не эффективность разбора.
Алгоритмы CYK, Эрли: $\mathcal{O}(n^3)$. \\
\end{tabular}

\end{frame}



% СЛАЙД 2: ИЕРАРХИЯ КЛАССОВ ГРАММАТИК
% Текст слайда подготовил Швачко Никита 5030102/20202
% Дата 11.01.2026
\begin{frame}
\frametitle{1.3.9. Классификация контекстно-свободных грамматик (3/5)}

\odmkey{Иерархия классов грамматик}{Иерархия!грамматик} по Хомскому и их свойства:

\begin{center}
\begin{tabular}{|c|c|c|c|}
\hline
\textbf{Тип} & \textbf{Название} & \textbf{Распознаватель} & \textbf{Сложность} \\
\hline
0 & С фразовой структурой & Машина Тьюринга & Полуразрешимо \\
\hline
1 & Контекстно-зависимые (НС) & ЛО-автомат & $\mathcal{O}(c^n)$ \\
\hline
2 & Контекстно-свободные & МП-автомат & $\mathcal{O}(n^3)$ \\
\hline
3 & Регулярные & Конечный автомат & $\mathcal{O}(n)$ \\
\hline
\end{tabular}
\end{center}

\begin{remark}
Внутри класса КСГ (тип 2) существует своя иерархия подклассов:
\[
\text{Регулярные} \subset LL(1) \subset LL(k) \subset \text{Однозначные КСГ} \subset \text{КСГ}.
\]
При этом для каждого фиксированного $k$: $LL(k) \subset LR(k)$.
\end{remark}

\end{frame}



% СЛАЙД 3: ПРИМЕРЫ ЯЗЫКОВ РАЗНЫХ КЛАССОВ
% Текст слайда подготовил Швачко Никита 5030102/20202
% Дата 11.01.2026
\begin{frame}
\frametitle{1.3.9. Классификация контекстно-свободных грамматик (4/5)}

\NB{Примеры языков, разделяющие классы:}

\begin{itemize}
\item \textbf{Регулярный язык:} $L_1 = \{a^n \mid n \ge 0\} = a^*$~--- распознаётся конечным автоматом.

\item \textbf{КС-язык, но не регулярный:} $L_2 = \{a^n b^n \mid n \ge 0\}$~--- 
требует МП-автомат (нужен стек для подсчёта $a$). 
Доказательство нерегулярности~--- лемма о разрастании.

\item \textbf{НС-язык, но не КС-язык:} $L_3 = \{a^n b^n c^n \mid n \ge 0\}$~--- 
требует линейно-ограниченный автомат. 
Доказательство~--- лемма Огдена (см. п.~1.3.4).

\item \textbf{Рекурсивно перечислимый, но не КЗ:} 
язык $\{w \mid \text{машина Тьюринга } M_w \text{ останавливается}\}$~--- 
проблема остановки неразрешима линейно ограниченным автоматом.
\end{itemize}

\begin{digress}
Каждый следующий класс строго содержит предыдущий. 
Это обосновывает необходимость изучения всей иерархии Хомского.
\end{digress}

\end{frame}



% СЛАЙД 4: ГРАММАТИКИ В РЕГУЛЯРНОЙ ФОРМЕ
% Текст слайда подготовил Швачко Никита 5030102/20202
% Дата 11.01.2026
\begin{frame}
\frametitle{1.3.9. Классификация контекстно-свободных грамматик (5/5)}

\odmkey{Грамматика в регулярной форме}{Грамматика!в регулярной форме}~--- 
это КСГ, в которой правые части правил записаны с использованием 
регулярных операций (объединение, конкатенация, итерация).

\begin{example}
Вместо правил $A \to aA \mid a$ можно записать $A \to a^+$ или $A \to aa^*$.
\end{example}

\NB{Преимущества регулярной формы:}
\begin{itemize}
\item Более компактная запись грамматики
\item Естественное представление повторяющихся конструкций
\item Используется в РБНФ (\DMref{п.~1.1.8}{1.1.8.}) и современных парсер-генераторах
\end{itemize}

\begin{remark}
Грамматика в регулярной форме эквивалентна обычной КСГ~--- 
регулярные операции можно раскрыть, введя дополнительные нетерминалы.
Например, $A \to \alpha^*$ заменяется на $A \to \varepsilon \mid \alpha A$.
\end{remark}

\bigskip
\TODO{Отличный ПЛАН! \\
Теперь каждый пункт раскрыть, обосновать, насытить ссылками и примерами. \\
Интересно, кто такие современные парсер-генераторы?}

\end{frame}


\subsection{1.4. Методы синтаксического анализа}\label{1.4.}
\begin{frame}
\frametitle{1.4. Методы синтаксического анализа}
В самом общем случае \odmkey{синтаксический анализ}{Синтаксический анализ} 
(или \odmkey{разбор}{Разбор}) --- это
сопоставление заданному выражению некоторого языка специальной
структуры, характерной для языков этого класса.
\begin{example}
Синтаксический анализ предложения на русском языке по типу, 
членам предложения, частям речи и т.~д. 
\end{example}
\begin{remark} Иногда употребляют транслитерацию английского слова <<парсинг>>. 
\end{remark}

\smallskip
Применительно к языкам предметной области 
можно ограничиться языками, которые
порождаются контекстно-свободными грамматиками.
В таком случае результатом синтаксического анализа
является дерево разбора. 

\smallskip
\begin{remark}
Если грамматика неоднозначна, то требуется
найти {\em все} деревья разбора. 
\end{remark} 
    
\begin{digress}
Задача синтаксического анализа имеет больше 
практической значимости по сравнению с задачей
распознавания принадлежности слова языку.
Дело в том, что в результате синтаксического
анализа получается больше информации:
кроме ответа <<да>> мы получаем дерево разбора,
которое можно использовать для дальнейшей работы,
а кроме ответа <<нет>> мы получаем сообщения о синтаксических ошибках,
которые также можно использовать для дальнейшей работы. 
\end{digress}
\end{frame}


\subsubsection{1.4.1. Постановка задачи синтаксического анализа}
\begin{frame}
\frametitle{1.4.1. Постановка задачи синтаксического анализа}
\TODO{Синтез и анализ на основе грамматики. Что требуется от парсера. Разбор с возвратами и без возвратов. \\Что является результатом разбора для 
грамматик разных типов: регулярное выражение, дерево разбора, семантические падежи Филлмора 
и т.д. в отступлениях.}

\bigskip

\TODO{Синтаксический разбор как логический вывод.
Разрешимость и полуразрешимость.
Разбор перебором. Требования к эффективности.
Безвозвратный разбор и разбор с возвратами.
Ограничения на алгоритмы разбора для грамматик разных классов.}

\bigskip

\TODO{06.04.26. В целом  последующее написано плохо. Повторы и утроения, и при этом пропущена обработка ошибок, ретрансляция, предобработка.}

\end{frame}

\begin{frame}
	\frametitle{1.4.1. Постановка задачи синтаксического анализа (1/5)}
	\odmkey{Синтаксический анализ}{Синтаксический анализ} --- по заданным $G$ и $w\in\Sigma^*$ решить задачу распознавания $w\in\frak{L}(S)$ и, при необходимости, построить структурное свидетельство.
	
	Пусть грамматика
	\[
		G=\langle \mathrm{M},\Sigma,\Pi,S\rangle,\qquad 
		\frak{L}(S)=\{\,w\in\Sigma^*\mid S\Rightarrow^* w\,\}.
	\]
	
	\odmkey{Синтез}{Синтез}:
	\[
		\text{построить } w\in\frak{L}(S)\ \text{посредством вывода } S\Rightarrow^* w.
	\]
	
	\odmkey{Анализ}{Анализ}:
	\[
		\text{по } w \text{ установить } (w\in\frak{L}(S)?)\ \text{и/или построить свидетельство } S\Rightarrow^* w.
	\]

	
	Задача синтаксического анализа включает два варианта:
	\begin{enumerate}
		\item \odmkey{Распознавание (membership)}{Распознавание}:
		\[
			\text{вход: }(G,w),\qquad \text{выход: }(w\in\frak{L}(S)?).
		\]
		\item \odmkey{Построение структуры (parsing)}{Построение структуры}:
		\[
			\text{вход: }(G,w),\qquad \text{выход: структура, кодирующая вывод } S\Rightarrow^* w.
		\]
	\end{enumerate}
\end{frame}


% Слайд подготовил Морщинин Никита 5030102/20201
% Дата 11.01.2026
\begin{frame}
	\frametitle{1.4.1. Постановка задачи синтаксического анализа (2/5)}
	\odmkey{Результат}{Результат} (как свидетельство) зависит от класса грамматики (иерархия Хомского) 
    
    и постановки задачи анализа.
	
	\begin{itemize}
		\item Тип~3 (регулярные)
		Распознавание: принимающий прогон конечного автомата на $w$ (цепочка состояний/переходов). Парсинг по регулярному выражению: дерево
        
        разбора регулярного выражения.
		
		\item Тип~2 (КС-грамматики)
		Парсинг: дерево разбора $T_w$, удовлетворяющее $S\Rightarrow^* w$; при неоднозначности --- множество деревьев $\{T_w\}$.
		
		\item Тип~1 (КЗ-грамматики)
		Результат парсинга задаётся как вывод $S\Rightarrow^* w$ с указанием применённых правил и позиций применения, удовлетворяющих контекстным 
        
        условиям правил.
		
		\item Тип~0 (неограниченные) Результат парсинга (при успехе) --- некоторый вывод $S\Rightarrow^* w$; универсальной гарантии завершения для всех $w$ нет.
	\end{itemize}
	
	\begin{digress}
		В задачах обработки естественного языка построенная синтаксическая структура может служить входом для семантического этапа, например для назначения \odmkey{семантических ролей}{семантических падежей Филлмора}; это не является частью синтаксического распознавания.
	\end{digress}
\end{frame}


% Слайд подготовил Морщинин Никита 5030102/20201
% Дата 11.01.2026
\begin{frame}
	\frametitle{1.4.1. Постановка задачи синтаксического анализа (3/5)}
	\odmkey{Требования к парсеру}{Требования} формулируются для фиксированных $G$ и входа $w$ 

    
    относительно языка $\frak{L}(S)$:
	\begin{itemize}
		\item \odmkey{Корректность}{Корректность}: если парсер принимает $w$, то $w\in\frak{L}(S)$; при parsing выдаваемая структура корректно кодирует вывод $S\Rightarrow^* w$.
		\item \odmkey{Полнота}{Полнота}: для всех $w\in\frak{L}(S)$ парсер принимает $w$ (и строит структуру, если это 
        
        требуется постановкой).
		\item \odmkey{Детерминизм}{Детерминизм}: выбор действия однозначен и не использует перебор альтернатив.
		\item \odmkey{Эффективность}{Эффективность}: ресурсы ограничены функциями от $n=\Length(w)$, приемлемыми для практики.
		\item \odmkey{Диагностика}{Диагностика}: при $w\notin\frak{L}(S)$ формируется сообщение об ошибке и её локализации.
	\end{itemize}
	
	\medskip
	
	\odmkey{Неоднозначность}{Неоднозначность}: $\exists w\in\frak{L}(S)$, для которого существует более одного дерева разбора. В задачах компиляции неоднозначность устраняется переписыванием грамматики или фиксированными правилами разрешения конфликтов, поскольку разным деревьям могут соответствовать разные интерпретации.
\end{frame}


% Слайд подготовил Морщинин Никита 5030102/20201
% Дата 11.01.2026
\begin{frame}
	\frametitle{1.4.1. Постановка задачи синтаксического анализа (4/5)}
	\odmkey{Разбор как логический вывод}{Разбор как вывод}.
	Правила $\Pi$ трактуются как правила вывода; утверждение $w\in\frak{L}(S)$ эквивалентно существованию доказательства:
	\[
		w\in\frak{L}(S)\ \Longleftrightarrow\ \exists\ \text{вывод } S\Rightarrow^* w.
	\]
	Парсер реализует алгоритм поиска такого доказательства; дерево разбора является 
    
    внешним представлением найденного доказательства.
	
	\bigskip
	
	\odmkey{Разрешимость и полуразрешимость}{Разрешимость} задачи распознавания:
	\begin{itemize}
		\item Тип~3: разрешима (конечный автомат завершает вычисление на любом входе).
		\item Тип~2: разрешима (существуют общие алгоритмы разбора, завершающиеся на любом $(G,w)$).
		\item Тип~0: полуразрешима: при $w\in\frak{L}(S)$ перебор выводов находит свидетельство; при $w\notin\frak{L}(S)$ алгоритм может не закончить работу.
	\end{itemize}
\end{frame}


% Слайд подготовил Морщинин Никита 5030102/20201
% Дата 11.01.2026
\begin{frame}
	\frametitle{1.4.1. Постановка задачи синтаксического анализа (5/5)}
	\odmkey{Разбор с возвратами}{Разбор с возвратами} (backtracking) --- перебор альтернатив применения правил с возвратом и попыткой применить следующее правило при неудаче; в худшем случае число 
    
    частичных разборов экспоненциально по $n=|w|$, что нарушает требования эффективности.
	
	\medskip
	
	\odmkey{Безвозвратный разбор}{Безвозвратный разбор} (deterministic parsing) --- выбор действия однозначен по текущей конфигурации и ограниченному предпросмотру входа; перебор отсутствует, типично
    
    достигается линейная сложность по $n$.
	
	\begin{remark}
		Ограничения на алгоритмы разбора задаются классом грамматик: общий разбор применим к более широким классам при более высокой асимптотике; линейные
        
        безвозвратные методы применимы лишь к подклассам, обеспечивающим детерминизм выбора действий.
	\end{remark}
	
\bigskip
\TODO{Январские тезисы. \\
Работать, работать и работать.}	
	
\end{frame}

\subsubsection{1.4.2. Нисходящий разбор}
\begin{frame}
\frametitle{1.4.2. Нисходящий разбор}

При нисходящем разборе результат синтаксического анализа 
--- дерево разбора --- строится, начиная с корня, то есть с аксиомы КСГ.

\medskip 
\TODO{Возможно, несколько параграфов.\\
Связь нисходящего разбора с левым выводом.\\
Рекурсивный спуск. Пример с арифметическими выражениями. \\
$LL(k)$-грамматики. Отступление про ANTL-R.
\\
\color{red}$LL(*)???$}

%\bigskip
\TODO{\color{blue} Игнатьев--Караев--Кириллова, очень неровный материал, доклады сделаны.}
\end{frame}

% Текст слайда подготовил Чинь Тхи Тхи Хоай - 5030102/20202
% 1.4.2 Нисходящий разбор
% Дата 14.12.2025
\begin{frame}
\frametitle{1.4.2 Нисходящий разбор (1/7)}
Пусть $D$ --- дерево вывода в грамматике $G$. \TODO{Это всё не в тему.}

\odmkey{Сечение дерева}{Сечение дерева} --- множество вершин, пересекающее \odmkey{все}{все} пути от корня к листьям так, чтобы на каждом пути была ровно \odmkey{одна}{одна} вершина из этого множества.

\odmkey{Крона сечения}{Крона сечения} --- последовательность меток вершин из этого сечения, записанных слева направо.

\begin{example}        
        Дерево вывода с кроной $abaSbS$. 
        \includegraphics[width=0.25\textwidth]{1.4.2.png} 
\end{example}

\begin{lemma}
    Пусть $S = \alpha_0, \alpha_1, \dots, \alpha_n$ - вывод цепочки $\alpha_n$ из $S$ в КС-грамматике $G = (N, \Sigma, P, S)$. Тогда в $G$ можно построить дерево вывода $D$, для которого $\alpha_n$ --- крона, а $\{\alpha_0, \alpha_1, \dots, \alpha_{n-1}\}$ --- некоторые из крон сечений.
\end{lemma}
\end{frame}


\begin{frame}
\frametitle{1.4.2 Нисходящий разбор (2/7)}
\begin{proof}
    Построим последовательность деревьев вывода $D_i$, где $0 \leq i \leq n$, таких что $\alpha_i$ --- крона дерева $D_i$. Пусть $D_0$ --- дерево, состоящее из единственной вершины, помеченной $S$.
    
    Допустим, что
    \[
    \alpha_i = \beta_i A \gamma_i
    \]
    и после применения правила
    \[
    A \to X_1 X_2 \dots X_k
    \]
    к выделенному вхождению $A$ получается
    \[
    \alpha_{i+1} = \beta_i X_1 X_2 \dots X_k \gamma_i.
    \]

    Тогда дерево $D_{i+1}$ получается из $D_i$ добавлением к листу, помеченному выделенным вхождением $A$ (оно является $(|\beta_i| + 1)$-м символом кроны $D_i$), новых потомков, помеченных $X_1, X_2, \dots, X_k$ соответственно.
    
    Очевидно, что крона дерева $D_{i+1}$ будет $\alpha_{i+1}$. Дерево $D_n$ будет искомым деревом вывода $D$. \TODO{Эти леммы с доказательствамим нужны, но не здесь, а в другом месте.}
\end{proof}
\end{frame}


\begin{frame}
\frametitle{1.4.2 Нисходящий разбор (3/7)}
\begin{lemma}
    Пусть $D$ --- дерево вывода в КС-грамматике $G = (N, \Sigma, P, S)$ с кроной $\alpha$. Тогда $S \Rightarrow^* \alpha$.
\end{lemma}
\begin{proof}
    Рассмотрим последовательность сечений $(C_0,\dots,C_n)$, где
    \begin{enumerate}
        \item $C_0$ содержит только корень дерева $D$,
        \item $C_{i+1}$ для $0 \leq i < n$ получается из $C_i$ заменой одной нетерминальной вершины её прямыми потомками,
        \item $C_n$ --- крона дерева $D$.
    \end{enumerate}
    
    Если $\alpha_i$ --- крона сечения $C_i$, то $(\alpha_0, \alpha_1, \dots, \alpha_n)$ --- вывод цепочки $\alpha_n$ из $\alpha_0$ в $G$.
\end{proof}

Если в доказательстве леммы сечение $C_{i+1}$ получается из $C_i$ заменой самой левой нетерминальной вершины её прямыми потомками, то вывод $(\alpha_0, \alpha_1, \dots, \alpha_n)$ называется \odmkey{левым выводом}{левым выводом} цепочки $\alpha_n$ из $\alpha_0$ в грамматике $G$. \TODO{Это в другом месте. Ссылка.}

Если $S = \alpha_0, \alpha_1, \dots, \alpha_n = w$ --- левый вывод терминальной цепочки $w$, то каждая цепочка $\alpha_i$ $(0 \leq i < n)$ имеет вид:
\[
\alpha_i = x_i A_i \beta_i,\text{ где } x_i \in \Sigma^*, A_i \in N, \text{ и } \beta_i \in (N \cup \Sigma)^*.
\]    
и $\alpha_{i+1}$ получается из $\alpha_i$ заменой самого левого нетерминала $A_i$ правой частью правила.
\end{frame}


\begin{frame}
\frametitle{1.4.2 Нисходящий разбор (4/7)}
Для некоторой КС-грамматики $G$ цепочка $w \in L(G)$ считается \odmkey{разобранной}{разобранной} (или \odmkey{проанализированной}{проанализированной}), если известно одно (а возможно, и все) из её деревьев вывода.

\begin{definition} \TODO{Ужас! Пожалуйста, не надо!}
    Пусть $G = (N, \Sigma, P, S)$ --- КС-грамматика, правила которой занумерованы числами $1, 2, \dots, p$. \\
    Пусть $\alpha \in (N \cup \Sigma)^*$. \odmkey{Левым разбором}{Левым разбором} цепочки $\alpha$ называется последовательность правил, применённых при левом выводе цепочки $\alpha$ из $S$.
\end{definition}

Пусть $\pi = \{i_1, \dots, i_n\}$ --- левый разбор цепочки $w \in L(G)$, где $G = (N, \Sigma, P, S)$ --- КС-грамматика. Зная $\pi$, можно построить дерево разбора цепочки $w$ следующим нисходящим образом.

Начнём с корня, помеченного $S$. Тогда $i_1$ даёт правило, которое надо применить к $S$. Допустим, что $i_1$ --- номер правила
\[
S \Rightarrow X_1 X_2 \dots X_k.
\]
Присоединим к вершине, помеченной $S$, потомков $X_1, X_2, \dots, X_k$ и пометим их этими символами.
\end{frame}

\begin{frame}
\frametitle{1.4.2 Нисходящий разбор (5/7)}
Если $X_i$ --- первый слева нетерминал в цепочке $X_1 X_2 \dots X_k$, то первыми $i-1$ символами цепочки $w$ должны быть $X_1 X_2 \dots X_{i-1}$. Правило с номером $i_2$, например, должно тогда иметь вид
\[
X_i \Rightarrow Y_1 Y_2 \dots Y_\ell,
\]
и можно продолжить построение дерева разбора цепочки $w$, применяя ту же процедуру к вершине, помеченной $X_i$.

Продолжая в том же духе, можно построить всё дерево разбора цепочки $w$, соответствующее \odmkey{левому разбору}{левому разбору} этой цепочки.

Связь с левым выводом:

Если на каждом шаге выбирается \odmkey{самый левый}{самый левый} нетерминал текущей строки $\alpha_i$, то последовательность подстановок
\[
S = \alpha_0 \Rightarrow \alpha_1 \Rightarrow \dots \Rightarrow \alpha_k = w
\]
является \odmkey{левым выводом}{левым выводом} в грамматике $G$. Следовательно, нисходящий разбор моделирует левый вывод.

\TODO{Это первая фраза в тему. Но она непонятна. Что значит моделирует?}
\end{frame}

\begin{frame}
\frametitle{1.4.2 Нисходящий разбор (6/7)}
\odmkey{Рекурсивный вывод}{Рекурсивный вывод} --- метод доказательства принадлежности цепочки $w$ языку $L(A)$:

если для правила $A \to X_1 X_2 \dots X_k$ выполнено $w_i \in L(X_i)$ для всех $i$, то
\[
w = w_1 w_2 \dots w_k \in L(A).
\]

Пусть $G = (N, \Sigma, P, E)$, где
\[
N = \{E, E', T, T', F\}, 
\]
\[
\Sigma = \{+, *, (, ), \text{num}\},
\]
\[E \text{--- стартовый символ.}
\]

Правила $P$:
\begin{align*}
    E  &\to T E' \\
    E' &\to + T E' \mid \varepsilon \\
    T  &\to F T' \\
    T' &\to * F T' \mid \varepsilon \\
    F  &\to (E) \mid \text{num}
\end{align*}
\end{frame}

\begin{frame}
\frametitle{1.4.2 Нисходящий разбор (7/7)}


\begin{columns}[t]
\begin{column}{0.3\textwidth}
\scriptsize
\begin{algorithmic}
\STATE \textbf{Вход:}
\STATE КС-грамматика $G=\langle N,\Sigma,P,E\rangle$
\STATE поток токенов $a_1\dots a_m \in \Sigma,\ \$$
\STATE \textbf{Выход:} accept / reject

\STATE
\STATE \textbf{parse(tokens):}
\STATE продвинуться к первому токену
\STATE parseE(tokens)
\IF{next\_token = \$}
    \STATE ACCEPT
\ELSE
    \STATE REJECT
\ENDIF

\STATE
\STATE \textbf{parseE(tokens):} \COMMENT{$E \to T \, E'$}
    \STATE parseT(tokens)
    \WHILE{next\_token = $+$}
        \STATE consume($+$)
        \STATE parseT(tokens)
    \ENDWHILE
\end{algorithmic}
\end{column}

\begin{column}{0.3\textwidth}
\scriptsize
\begin{algorithmic}
\STATE \textbf{parseT(tokens):} \COMMENT{$T \to F \, T'$}
    \STATE parseF(tokens)
    \WHILE{next\_token = $*$}
        \STATE consume($*$)
        \STATE parseF(tokens)
    \ENDWHILE

\STATE
\STATE \textbf{parseF(tokens):} \COMMENT{$F \to (E)\mid\text{num}$}
    \IF{next\_token = num}
        \STATE consume(num)
    \ELSIF{next\_token = $($}
        \STATE consume($($)
        \STATE parseE(tokens)
        \STATE consume($)$)
    \ENDIF
\end{algorithmic}
\TODO{Ровно два слайда в тему, но оформлено неверно. 
1) Определение рекурсивного спуска. 2) Пример с арифметическими выражениями.
3) Анализатор арифметических выражений.\\
Посмотрите другие слайды!}
\end{column}
\end{columns}

\end{frame}

\begin{frame}
\frametitle{1.4.2. Нисходящий разбор (???)}

\TODO{Вторая попытка. \\
Стало намного лучше, но всё подлежит переделыванию. Неправильные обозначения, 
нет ссылок на уже введенный материал, взятые в интернете безапелляционные и при этом НЕВЕРНЫЕ оценочные суждения. \\ 
Очень плохо аргументированы примеры: можно перейти к прямой польской записи, а можно в инфиксной записи избавиться от левой рекурсии. \\
Напрасно Вы убрали рекурсивный спуск --- с него надо начинать! \\
Совершенно не используется идея простой или разделённой грамматики и её последствия 
для современного языкотворчества и реализации DSL. \\
Атомов много, а хороших нет. }
\end{frame}

% Текст слайда подготовил Чинь Тхи Тхи Хоай - 5030102/20202
% 1.4.2 Нисходящий разбор
% Дата 10.01.2026
\begin{frame}
\frametitle{1.4.2. Нисходящий разбор (1/7)}

\odmkey{Нисходящий разбор}{Нисходящий разбор} — метод синтаксического анализа, при котором дерево вывода строится от корня (аксиомы грамматики) к листьям (входной цепочке). Алгоритмически это соответствует поиску \odmkey{левого вывода}{левого вывода} $S \Rightarrow_{lm}^* \omega$ для входной строки $\omega$.

Пусть задана КС-грамматика $G = \langle N, T, R, S \rangle$. На каждом шаге анализатор рассматривает текущую сентенциальную форму $\alpha A \gamma$, где $\alpha \in T^*$ — обработанный префикс, $A \in N$ — левый нетерминал, подлежащий замене. Требуется выбрать продукцию $A \to \beta$ так, чтобы существовало продолжение вывода $\alpha \beta \gamma \Rightarrow_{lm}^* \omega$.

Основная проблема — недетерминизм выбора продукции. В общем случае требуется перебор (backtracking), что неэффективно. Практический интерес представляют грамматики, допускающие детерминированный выбор на основе ограниченного предпросмотра (lookahead).

\begin{digress}
Левая рекурсия делает грамматику непригодной для нисходящего разбора — процедура анализа зацикливается. Методы устранения левой рекурсии рассмотрены в п. 1.3.4.
\end{digress}
\end{frame}

\begin{frame}
\frametitle{1.4.2. Нисходящий разбор (2/7)}
Ключевая идея — возможность выбора продукции $A \to \beta$ на основе следующего символа входной цепочки. Если для любого нетерминала $A$ и любого входного символа $a$ существует не более одной применимой продукции, разбор может выполняться без возвратов.

Для формализации вводятся два множества:
\begin{itemize}
    \item $\textrm{FIRST}(\alpha)$ — терминалы, которые могут начинать вывод из цепочки $\alpha$
    \item $\textrm{FOLLOW}(A)$ — терминалы, которые могут следовать непосредственно за $A$ в выводе из $S$
\end{itemize}

Грамматика $G$ называется \odmkey{LL(1)-грамматикой}{LL(1)-грамматикой}, если для каждого нетерминала $A$ и любых двух различных продукций $A \to \alpha$ и $A \to \beta$ выполняется:
\[
\textrm{FIRST}(\alpha) \cap \textrm{FIRST}(\beta) = \varnothing,
\]
а если $\varepsilon \in \textrm{FIRST}(\alpha)$, то дополнительно:
\[
\textrm{FIRST}(\beta) \cap \textrm{FOLLOW}(A) = \varnothing.
\]
\end{frame}

\begin{frame}
\frametitle{1.4.2. Нисходящий разбор (3/7)}
\begin{example}
Рассмотрим грамматику арифметических выражений:
\[
\begin{array}{rcl}
E & \to & T E' \\
E' & \to & +\,T\,E' \mid \varepsilon \\
T & \to & F\,T' \\
T' & \to & *\,F\,T' \mid \varepsilon \\
F & \to & (\,E\,) \mid \mathbf{id}
\end{array}
\]

Вычислим необходимые множества (опуская детали):
\begin{align*}
&\textrm{FIRST}(E') = \{ +, \varepsilon \}, \quad \textrm{FOLLOW}(E') = \{ ), \$ \} \\
&\textrm{FIRST}(T') = \{ *, \varepsilon \}, \quad \textrm{FOLLOW}(T') = \{ +, ), \$ \}
\end{align*}

Проверим условие LL(1) для $E'$: 
$\textrm{FIRST}(+\,T\,E') = \{ +\}$, 
$\textrm{FIRST}(\varepsilon) = \varnothing$,
и $\{+\} \cap \{), \$\} = \varnothing$. Условие выполняется.
\end{example}
\end{frame}

\begin{frame}
\frametitle{1.4.2. Нисходящий разбор (4/7)}
Для LL(1)-грамматики строится таблица $M[A, a]$, где $A \in N$, $a \in T \cup \{\$\}$.

\textbf{Алгоритм построения:}
Для каждой продукции $A \to \alpha$:
\begin{enumerate}
    \item Для каждого $a \in \textrm{FIRST}(\alpha)$ положить $M[A, a] = A \to \alpha$.
    \item Если $\varepsilon \in \textrm{FIRST}(\alpha)$, то для каждого $b \in \textrm{FOLLOW}(A)$ положить $M[A, b] = A \to \alpha$.
\end{enumerate}

\begin{center}
\small
\begin{tabular}{|c|c|c|c|c|c|c|}
\hline
 & \textbf{id} & \textbf{+} & \textbf{*} & \textbf{(} & \textbf{)} & \textbf{\$} \\ \hline
$E$  & $E \to T E'$ &  &  & $E \to T E'$ &  &  \\ \hline
$E'$ &  & $E' \to +\,T\,E'$ &  &  & $E' \to \varepsilon$ & $E' \to \varepsilon$ \\ \hline
$T$  & $T \to F\,T'$ &  &  & $T \to F\,T'$ &  &  \\ \hline
$T'$ &  & $T' \to \varepsilon$ & $T' \to *\,F\,T'$ &  & $T' \to \varepsilon$ & $T' \to \varepsilon$ \\ \hline
$F$  & $F \to \mathbf{id}$ &  &  & $F \to (\,E\,)$ &  &  \\ \hline
\end{tabular}
\end{center}

\begin{remark}
Все ячейки либо содержат ровно одну продукцию, либо пусты — грамматика действительно LL(1).
\end{remark}

\end{frame}

\begin{frame}
\frametitle{1.4.2. Нисходящий разбор (5/7)}

Алгоритм разбора по таблице:
Используется стек. Начальное состояние: стек содержит $[S, \$]$.

На каждом шаге пусть $X$ — верхний символ стека, $a$ — текущий входной символ.

\begin{enumerate}
    \item Если $X \in T$ и $X = a$, то снять $X$ со стека и продвинуть вход.
    \item Если $X \in T$ и $X \ne a$, то ошибка.
    \item Если $X \in N$, то найти в таблице $M[X, a] = X \to Y_1 \dots Y_k$.
        \begin{itemize}
            \item Если продукция найдена, снять $X$ и положить на стек символы $Y_k, \dots, Y_1$ (чтобы $Y_1$ оказался сверху). Для $\varepsilon$-продукции просто снять $X$.
            \item Если не найдено — ошибка.
        \end{itemize}
    \item Если $X = \$$ и $a = \$$, разбор успешно завершён.
\end{enumerate}

\begin{remark}
Этот алгоритм работает за линейное время $O(|\omega|)$ и использует память $O(|\omega|)$ для стека.
\end{remark}

\end{frame}

\begin{frame}
\frametitle{1.4.2. Нисходящий разбор (6/7)}

Пример разбора цепочки $\mathbf{id} + \mathbf{id}$

\begin{center}
\small
\begin{tabular}{|c|c|l|}
\hline
\textbf{Стек} & \textbf{Вход} & \textbf{Действие} \\ \hline
$E\$$ & $\mathbf{id} + \mathbf{id}\$$ & $M[E, \mathbf{id}] = E \to T E'$ \\ \hline
$T\,E'\$$ & $\mathbf{id} + \mathbf{id}\$$ & $M[T, \mathbf{id}] = T \to F\,T'$ \\ \hline
$F\,T'\,E'\$$ & $\mathbf{id} + \mathbf{id}\$$ & $M[F, \mathbf{id}] = F \to \mathbf{id}$ \\ \hline
$\mathbf{id}\,T'\,E'\$$ & $\mathbf{id} + \mathbf{id}\$$ & совпадение $\mathbf{id}$ \\ \hline
$T'\,E'\$$ & $+ \mathbf{id}\$$ & $M[T', +] = T' \to \varepsilon$ \\ \hline
$E'\$$ & $+ \mathbf{id}\$$ & $M[E', +] = E' \to +\,T\,E'$ \\ \hline
$+\,T\,E'\$$ & $+ \mathbf{id}\$$ & совпадение $+$ \\ \hline
$T\,E'\$$ & $\mathbf{id}\$$ & $M[T, \mathbf{id}] = T \to F\,T'$ \\ \hline
$F\,T'\,E'\$$ & $\mathbf{id}\$$ & $M[F, \mathbf{id}] = F \to \mathbf{id}$ \\ \hline
$\mathbf{id}\,T'\,E'\$$ & $\mathbf{id}\$$ & совпадение $\mathbf{id}$ \\ \hline
$T'\,E'\$$ & $\$$ & $M[T', \$] = T' \to \varepsilon$ \\ \hline
$E'\$$ & $\$$ & $M[E', \$] = E' \to \varepsilon$ \\ \hline
$\$$ & $\$$ & успех \\ \hline
\end{tabular}
\end{center}

\begin{remark}
Применённые продукции образуют левый вывод.
\end{remark}

\end{frame}

\begin{frame}
\frametitle{1.4.2. Нисходящий разбор (7/7)}

\begin{itemize}
    \item Каждая LL(1)-грамматика однозначна.
    \item Существуют однозначные грамматики, не являющиеся LL($k$) ни для какого $k$.
    \item Проверка принадлежности грамматики классу LL(1) разрешима.
    \item Класс языков, порождаемых LL(1)-грамматиками, строго включает регулярные языки и строго содержится в классе детерминированных КС-языков.
\end{itemize}


\begin{itemize}
    \item \textbf{LL($k$)}: использование $k$ символов предпросмотра расширяет класс грамматик, но таблицы имеют размер $\cal{O}(|N| \cdot |T|^k)$.
    \item \textbf{LL($*$)}: современные инструменты (ANTLR) используют динамический анализ с помощью конечных автоматов.
    \item \textbf{Рекурсивный спуск}: альтернативная реализация, где каждый нетерминал соответствует процедуре. Удобен для ручного кодирования и встраивания семантических действий.
\end{itemize}

\begin{digress}
Нисходящий разбор широко применяется благодаря своей наглядности и простоте интеграции семантических действий. Однако он менее мощен, чем восходящие методы (LR), и требует устранения левой рекурсии.
\end{digress}

\end{frame}

% Текст слайда подготовил Чинь Тхи Тхи Хоай - 5030102/20202
% 1.4.2. Нисходящий разбор
% Дата 31.03.2026
\begin{frame}
\frametitle{1.4.2. Нисходящий разбор (1/7)}

\odmkey{Нисходящим разбором}{Нисходящий разбор} называется такое построение дерева
разбора, которое начинается с корня, то есть с аксиомы грамматики, и заканчивается
кроной, то есть входной цепочкой. Тем самым нисходящий разбор соответствует поиску
левостороннего вывода, о котором говорилось в \DMref{п.~1.3.1}{1.3.1.}

Как отмечалось в \DMref{п.~1.4.1}{1.4.1.}, при синтаксическом анализе требуется по
входной цепочке восстановить её синтаксическую структуру. Для нисходящего разбора
это означает, что анализатор, начиная с аксиомы, должен шаг за шагом предсказывать,
какое правило вывода следует применить дальше.

В общем случае такой выбор может требовать перебора. Поэтому практически важны прежде
всего такие формы грамматик, для которых очередное действие анализатора можно
выбирать без возвратов.

\end{frame}

\begin{frame}
\frametitle{1.4.2. Нисходящий разбор (2/7)}

Самая естественная форма нисходящего разбора --- \odmkey{рекурсивный спуск}
{Рекурсивный спуск}. При нём каждому нетерминалу $A$ ставится в соответствие
процедура $P_A$, которая должна распознать тот фрагмент входной цепочки, который
выводим из $A$.

Если выбрано правило
\[
A \to X_1 \dots X_n,
\]
то процедура $P_A$ последовательно пытается распознать символы $X_1$, \dots, $X_n$:
терминал --- сравнением с текущим символом входной цепочки, а нетерминал ---
вызовом соответствующей процедуры.

Разбор считается успешным, если процедура для аксиомы завершилась без ошибки и при
этом вся входная цепочка исчерпана.

\begin{remark}
Левая рекурсия делает рекурсивный спуск непригодным: процедура вызывает саму себя,
не продвигаясь по входной цепочке. Преобразования КС-грамматик, в том числе
устранение левой рекурсии, рассматривались в \DMref{п.~1.3.2}{1.3.2.}
\end{remark}

\end{frame}

\begin{frame}
\frametitle{1.4.2. Нисходящий разбор (3/7)}

\begin{example}
Проще всего идея рекурсивного спуска видна на выражениях в прямой польской записи:
\[
E \to a \mid +EE \mid *EE.
\]
Здесь $a$ обозначает атомарное выражение.

В самом деле, если текущий символ есть $a$, то выражение распознано. Если текущий
символ есть $+$, то после него надо распознать ещё два выражения; то же верно и для
символа $*$.

Таким образом, форма записи выражения здесь прямо подсказывает структуру процедуры
разбора. Этот пример удобен именно потому, что в нём идея нисходящего разбора
видна без дополнительных преобразований грамматики.
\end{example}

\end{frame}

\begin{frame}
\frametitle{1.4.2. Нисходящий разбор (4/7)}

Однако в языках программирования и DSL пользователю обычно нужна не польская,
а более привычная инфиксная запись. В этом случае рекурсивный спуск применяют
не к исходной грамматике, а к эквивалентной грамматике, полученной средствами
\DMref{п.~1.3.2}{1.3.2.}

\begin{example}
Вместо грамматики
\[
E \to E + T \mid T,\qquad
T \to T * F \mid F,\qquad
F \to (E) \mid a
\]
берут, например, эквивалентную грамматику
\[
E \to T E',\qquad
E' \to + T E' \mid \varepsilon,
\]
\[
T \to F T',\qquad
T' \to * F T' \mid \varepsilon,
\]
\[
F \to (E) \mid a.
\]
После такого преобразования рекурсивный спуск становится возможным.
\end{example}

Тем самым привычная для пользователя форма записи может быть согласована с удобной
для реализации структурой грамматики.

\end{frame}

\begin{frame}
\frametitle{1.4.2. Нисходящий разбор (5/7)}

Посмотрим, как работает рекурсивный спуск для цепочки
\[
a+a*a.
\]
Процедура $P_E$ вызывает $P_T$, затем $P_{E'}$. Первая из них распознаёт начальный
символ $a$. После этого процедура $P_{E'}$ видит символ $+$, читает его, вызывает
$P_T$ и тем самым переходит к разбору правой части суммы.

На этом шаге процедура $P_T$ сначала распознаёт символ $a$, а затем процедура
$P_{T'}$ видит символ $*$, читает его и распознаёт следующий символ $a$.
После этого и $P_{T'}$, и $P_{E'}$ завершаются по $\varepsilon$-ветви.

\begin{remark}
Тем самым порядок вызовов процедур непосредственно воспроизводит структуру дерева
разбора. Именно в этом состоит главная наглядность рекурсивного спуска.
\end{remark}

\end{frame}

\begin{frame}
\frametitle{1.4.2. Нисходящий разбор (6/7)}

Если выбор правила удаётся делать по ограниченному предпросмотру входной цепочки,
то говорят о \odmkey{предсказывающем разборе}{Предсказывающий разбор}. Связанные
с этим классы грамматик LL$(k)$ были введены в \DMref{п.~1.3.7}{1.3.7.}, поэтому здесь
мы не повторяем соответствующих определений.

С точки зрения реализации особенно важен тот случай, когда выбор очередной
альтернативы прост и прозрачен. Именно таким случаем являются
\odmkey{разделённые грамматики}{Грамматика!разделённая}, рассмотренные в
\DMref{п.~1.3.8}{1.3.8.}

Для разделённой грамматики устройство нисходящего разбора видно уже по самой записи
правил: различным альтернативам соответствуют различные начальные терминалы, и
поэтому анализатору не приходится выполнять сложное предсказание.

\end{frame}

\begin{frame}
\frametitle{1.4.2. Нисходящий разбор (7/7)}

Итак, нисходящий разбор естественно рассматривать в следующей последовательности:
сначала рекурсивный спуск как прямое программирование грамматики, затем --- как
предсказывающий разбор для подходящих классов грамматик, и, наконец, как особенно
простую и удобную технику для разделённых грамматик.

\begin{digress}
С практической точки зрения это важно для языкотворчества и реализации DSL.
Язык предметной области приходится не только один раз определить, но и многократно
изменять. Чем проще устроен синтаксический анализ, тем легче развивать язык,
встраивать семантические действия и локализовать ошибки.
\end{digress}

Поэтому нисходящий разбор интересен не только как частный метод синтаксического
анализа, но и как ориентир при выборе формы грамматики в прикладной работе.

\end{frame}

\subsubsection{1.4.3. Восходящий разбор}
\label{1.4.3.}
\begin{frame}
\frametitle{1.4.3. Восходящий разбор}

При восходящем разборе результат синтаксического анализа 
--- дерево разбора --- строится, начиная с кроны, то есть с входной цепочки.

\medskip
\TODO{Возможно, несколько параграфов.\\
Связь восходящего разбора с обращенным правым выводом.\\
Перенос--свёртка.  Пример с арифметическими выражениями.\\ 
$LR(k)$-грамматики. Связь с МП-автоматами.}
\end{frame}


\begin{frame}
\frametitle{1.4.3. Восходящий разбор(1/6)}
Если при построении вывода $\alpha$, всегда заменяется самый правый нетерминальный символ промежуточной цепочки, то вывод называется \odmkey{правым}{правым} или \odmkey{правосторонним выводом}{правосторонним! выводом} $\alpha$.
\begin{columns}[t]
\begin{column}{10.1cm}
\begin{example}
Рассмотрим КСГ $G_1$ с правилами
$E \rightarrow E + T \mid T$, 
$T \rightarrow T * F \mid F$,
$F \rightarrow (E) \mid a$.
Дерево разбора справа соответствует правостороннему выводу цепочки $a + a * a$, описанному следующим образом:
 $E \Rightarrow E + T \Rightarrow E + T * F,  \Rightarrow E + T * a \Rightarrow E + F * a \Rightarrow E + a * a \Rightarrow T + a * a \Rightarrow F + a * a \Rightarrow a + a * a.$ Этот вывод соответствует стандартной интерпретации с учётом приоритета операций (умножение выполняется раньше сложения), то есть $a + (a * a)$. Данная грамматика является неоднозначной, поэтому для цепочки $a + a * a$ существует также альтернативный левосторонний вывод, приводящий к неверной группировке $(a + a) * a$.
\end{example}

\end{column}
\begin{column}{5cm}
\vspace{-20pt}
\begin{figure}[h]
\begin{center}
\includegraphics[height=50mm]{1_4_3.jpg}
\end{center}
\end{figure}
\end{column}
\end{columns}


\end{frame}


\begin{frame}
\frametitle{1.4.3. Восходящий разбор(2/6)}

Напротив левого разбора, \odmkey{правый разбор}{Правый!разбор} цепочки \(\alpha \) — это обращение последовательности правил, примененных при правом выводе  цепочки \(\alpha \) из стартового символа \(S\).\\
Вообще \odmkey{правый разбор}{Правый!разбор} цепочки \(\alpha \) в грамматике G = $\langle \mathrm{M}, \Sigma, \Pi, S \rangle$ - это последовательность правил, с помощью которых можно свернуть цепочку \(\alpha \) начальному символу S. В терминах деревьев выводов \odmkey{правый анализ}{Правый!анализ} цепочки \(\alpha \) представляет собой последовательность отсечений основ, сводящую дерево вывода с кроной \(\alpha \) к одной вершине, помеченной S. В сущности это равносильно процессу \odmkey{заполнения}{заполнения} дерева вывода, начинающемуся с одной только кроны и идущему от листьев от корню, то есть снизу вверх. Поэтому с процессом порождения правого разбора часто ассоциируется термин \odmkey{восходящий анализ}{Восходящий!анализ} или \odmkey{восходящий разбор}{Восходящий!разбор}.


\end{frame}


\begin{frame}
\frametitle{1.4.3. Восходящий разбор (3/6)}
Метод \odmkey{перенос-свёртка}{Перенос!свёртка} представляет собой фундаментальную технику восходящего синтаксического анализа в теории компиляторов. Большинство восходящих анализаторов используют стек и два основных действия:
\begin{itemize}
\item \odmkey{Перенос (Shift)}{Перенос!Shift} - перемещение следующего входного символа в стек.

\item \odmkey{Свёртка (Reduce)}{Свёртка!Reduce} - если на вершине стека находится последовательность символов, совпадающая с правой частью правила \(A\rightarrow \beta \), она заменяется на нетерминал \(A\).
\end{itemize}
Процесс продолжается до тех пор, пока стек не будет содержать только стартовый символ грамматики, а входной поток не исчерпан (accept).
\end{frame}


\begin{frame}
\frametitle{1.4.3. Восходящий разбор(4/6)}


\begin{example}
Вернемся к приведенному выше примеру, но уже с использованием метода перенос-свёртка. Ниже приведена последовательность шагов перенос-свёртка парсера (восходящий разбор).
\begin{table}[h]
\centering
\scriptsize  % Thu nhỏ chữ một mức
\begin{tabular}{|c|c|c|c|c|}
\hline
\textbf{Шаг} & \textbf{Стек} & \textbf{Остаток входа} & \textbf{Действие} & \textbf{Правило} \\
\hline
0  & $\epsilon$& a + a * a \$& Shift     & —                          \\
\hline
1  & a & + a * a \$ & Reduce & F $\to$ a\\
\hline
2  & F & + a * a \$ & Reduce & T $\to$ F \\
\hline
3  & T & + a * a \$ & Reduce & E $\to$ T\\
\hline
4  & E& + a * a \$& Shift & — \\
\hline
5  & E + & a * a \$ & Shift& —\\
\hline
6  & E + a & * a \$  & Reduce& F $\to$ a\\
\hline
7  & E + F & * a \$  & Reduce & T $\to$ F \\
\hline
8  & E + T & * a \$  & Shift & —  \\
\hline
9  & E + T *  & a \$ & Shift & — \\
\hline
10 & E + T * a& \$ & Reduce    & F $\to$ a \\
\hline
11 & E + T * F  & \$ & Reduce    & T $\to$ T * F\\
\hline
12 & E + T & \$ & Reduce    & E $\to$ E + T  \\
\hline
13 & E  & \$ & Accept    & —  \\
\hline
\end{tabular}
\label{tab:shift-reduce-example}
\end{table}
\end{example}
\end{frame}


\begin{frame}
\frametitle{1.4.3. Восходящий разбор(5/6)}

\begin{remark}
Однако «чистый» перенос-свёртка (без дополнительных механизмов) часто сталкивается с конфликтами:
\begin{itemize}
\item \odmkey{Перенос-свёртка}{Перенос!перенос}: в текущей ситуации возможно как выполнение переноса, так и свёртки.
\item \odmkey{Свёртка-свёртка}{Свёртка!свёртка}:на вершине стека одновременно подходят правые части нескольких разных правил.
\end{itemize}
Поэтому перенос-свёртка работает эффективно и детерминированно только при наличии механизма однозначного разрешения конфликтов. \\
\end{remark}
\odmkey{$LR(k)$}{LR(k)} как раз и является усовершенствованной и стандартизированной версией алгоритма перенос-свёртка, которая гарантирует детерминизм за счёт использования предпросмотра (lookahead) на k символов. Он сохраняет базовую механику shift и reduce, но добавляет предварительно построенную таблицу разбора, полностью устраняющую конфликты для грамматик класса $LR(k)$. На практике подавляющее большинство компиляторов использует подклассы с $k=1$: $SLR(1)$, $LALR(1)$ (yacc, bison) или полные $LR(1)$.

\end{frame}


\begin{frame}
\frametitle{1.4.3. Восходящий разбор(6/6)}
Восходящий разбор по $LR(k$) (в частности, $LR(1)$) напрямую реализуется на ДМП-автомате с помощью алгоритма перенос-свёртка : магазин (стек) хранит последовательность состояний и символов, соответствующих префиксу правого разбора; таблица ACTION определяет действия перенос  или свёртка  на основе текущего состояния и просмотра k символов вперёд; таблица GOTO управляет переходами после свёртки. Отсутствие конфликтов в таблицах гарантирует детерминизм. Таким образом, существование $LR(k)$-грамматики эквивалентно возможности эффективного восходящего синтаксического анализа языка детерминированным магазинным автоматом.

\bigskip
\TODO{Структурно материал раздела построен верно.\\
Источник взят какой-то очень поверхностный.\\
Все абзацы, начиная с первого, содержат неточности и неверные утверждения, текст надобно весь переработать.\\
Восемь атомов с оценкой удовлетворительно.  
} 

\end{frame}


\subsubsection{1.4.4. Специальные методы синтаксического анализа}
% Текст слайда подготовил Чинь Тхи Тхи Хоай - 5030102/20202
% 1.4.4. Специальные методы синтаксического анализа
% Дата 06.01.2026
% Ревизия ФАН 15.01.26
\begin{frame}
\frametitle{1.4.4. Специальные методы синтаксического анализа (1/9)}
Методы синтаксического анализа, рассмотренные в предыдущих параграфах 
(\DMref{п.~1.4.2}{1.4.2} и \DMref{1.4.3}{1.4.3.}), ориентированы на использование контекстно-свободных грамматик, обладающих дополнительными структурными свойствами, такими как детерминированность, отсутствие левой рекурсии или принадлежность к классам $LL(k)$ и $LR(k)$. Эти ограничения позволяют строить эффективные алгоритмы разбора, однако существенно сужают класс допустимых грамматик.

Для анализа языков, заданных контекстно-свободными грамматиками с другими
ограничениями или без ограничений разработаны специальные методы, 
основанные на иных математических идеях. В разделе рассматриваются три 
знаменитых алгоритма:
\begin{enumerate}
    \item \odmkey{Алгоритм Дейкстры}{Алгоритм!Дейкстры} 
    (\odmkey{алгоритм сортировочной станции}{Алгоритм!сортировочной станции}) 
    для анализа арифметических выражений;
    \item \odmkey{Алгоритм Эрли}{Алгоритм!Эрли} для произвольных КС-грамматик;
    \item \odmkey{Алгоритм Кока--Янгера--Касами}{Алгоритм!Кока--Янгера--Касами} (\odmkey{алгоритм CYK}{Алгоритм!CYK}) для грамматик в нормальной форме
    Хомского.
\end{enumerate}

Эти алгоритмы демонстрируют поучительные компромиссы между общностью, эффективностью и сложностью реализации. \TODO{именные сноски}
\end{frame}

\begin{frame}
\frametitle{1.4.4. Специальные методы синтаксического анализа (2/9)}
\odmkey{Алгоритм Дейкстры}{Алгоритм Дейкстры} (shunting-yard algorithm) решает задачу синтаксического анализа арифметических выражений за один проход ($\cal{O}(n)$). Его основная идея — использование двух стеков: \odmkey{стека операндов}{стека операндов} $V$ и \odmkey{стека операторов}{стека операторов} $O$. \TODO{оператор операция, буква О}

\textbf{Правила алгоритма:} 
Для каждого символа $x$ во входной строке: \TODO{псевдокод}
\begin{enumerate}
\item Если $x$ - операнд ($a, b, c, \dots$), то поместить $x$ в стек $V$.
\item Если $x$ - оператор ($+, *, \dots$), то пока
$O \neq \varepsilon$ и $\mathrm{Pr}(\mathrm{Top}(O)) \ge \mathrm{Pr}(x)$,
выполнять \odmkey{свёртку}{свёртку}. Затем поместить $x$ в $O$.
\item Если $x = \texttt{(}$, то поместить его в $O$.
\item Если $x = \texttt{)}$, то выполнять свёртки, пока
$\mathrm{Top}(O) \neq \texttt{(}$, после чего удалить \texttt{(} из $O$.
\item В конце строки выполнить все оставшиеся свёртки.
Результат - верхний элемент стека $V$.
\end{enumerate}

\textbf{Операция свёртки:}
\begin{enumerate}
\item Извлечь оператор $\oplus$ из $O$.
\item Извлечь операнды $y$ и $x$ из $V$.
\item Поместить в $V$ результат $x \oplus y$
(или поддерево разбора $(\oplus, x, y)$).
\end{enumerate}
\end{frame}

\begin{frame}
\frametitle{1.4.4. Специальные методы синтаксического анализа (3/9)}

\begin{example}
Рассмотрим выражение $a + b * c$ и грамматику:
\[
E \to E + T \mid T, \qquad
T \to T * F \mid F, \qquad
F \to a \mid b \mid c.
\]

Алгоритм Дейкстры неявно строит правосторонний вывод.
Процесс разбора: \TODO{таблица}
\end{example}

\begin{itemize}
\item \textbf{Шаг 1:} Вход $a$. $V = [a]$, $O = [\,]$.
\item \textbf{Шаг 2:} Вход $+$. $V = [a]$, $O = [+]$.
\item \textbf{Шаг 3:} Вход $b$. $V = [a, b]$, $O = [+]$.
\item \textbf{Шаг 4:} Вход $*$. Так как $\Pr(*) > \Pr(+)$,
оператор $*$ помещается в стек.  
$V = [a, b]$, $O = [+, *]$.
\item \textbf{Шаг 5:} Вход $c$. $V = [a, b, c]$, $O = [+, *]$.
\item \textbf{Шаг 6:} Конец строки. Выполняются свёртки:
\begin{enumerate}
\item Для $*$: $V = [a, (b * c)]$, $O = [+]$.
\item Для $+$: $V = [(a + (b * c))]$, $O = [\,]$.
\end{enumerate}
\end{itemize}

Итоговое дерево разбора соответствует выводу: \TODO{картинка}
\[
E \Rightarrow E + T
\Rightarrow E + T * F
\Rightarrow E + T * c
\Rightarrow E + b * c
\Rightarrow a + b * c.
\]

\end{frame}

\begin{frame}
\frametitle{1.4.4. Специальные методы синтаксического анализа (4/9)}
\odmkey{Алгоритм Эрли}{Алгоритм Эрли} (1970) представляет собой универсальный метод синтаксического анализа для произвольных КС-грамматик. Его основная идея — динамическое построение множества \odmkey{ситуаций}{ситуаций} $S_j$ для каждой позиции $j$ во входной строке $w = a_1 \dots a_n$.

\odmkey{Ситуацией}{Ситуация} называется тройка $[A \to \alpha \cdot \beta, i]$, где:
\TODO{Ситуация ? Конфигурация}
\begin{itemize}
    \item $A \to \alpha\beta \in R$ — правило грамматики;
    \item Точка $\cdot$ указывает текущую позицию в правиле;
    \item $i$ — позиция во входной строке, где началось применение этого правила.
\end{itemize}

Алгоритм инициализирует $S_0$ и последовательно строит множества $S_1, \dots, S_n$ с помощью трёх операций:
\begin{enumerate}
    \item \textbf{Предсказание:} Если $[A \to \alpha \cdot B\beta, i] \in S_j$, добавить $[B \to \cdot \gamma, j]$ для всех правил $B \to \gamma \in R$.
    \item \textbf{Сканирование:} Если $[A \to \alpha \cdot a\beta, i] \in S_j$ и $a = w_{j+1}$, добавить $[A \to \alpha a \cdot \beta, i]$ в $S_{j+1}$.
    \item \textbf{Завершение:} Если $[A \to \gamma \cdot, i] \in S_j$, то для всех $[B \to \alpha \cdot A\beta, k] \in S_i$ добавить $[B \to \alpha A \cdot \beta, k]$ в $S_j$.
\end{enumerate}

Строка допускается, если в $S_n$ присутствует ситуация $[S' \to S \cdot, 0]$, где $S'$ — новый стартовый нетерминал, добавленный к грамматике. \TODO{Слова = Агафонов/Курочкин}
\end{frame}

\begin{frame}
\frametitle{1.4.4. Специальные методы синтаксического анализа (5/9)}
\begin{theorem}
Алгоритм Эрли корректен и полон: $w \in \mathfrak{L}(G)$ тогда и только тогда, когда $[S' \to S \cdot, 0] \in S_n$.
\end{theorem}
\begin{proof}
Индукция по длине вывода. Корректность: каждая добавляемая ситуация соответствует возможному выводу в грамматике. Полнота: для любого вывода $S \Rightarrow^* w$ в грамматике соответствующая ситуация будет построена алгоритмом.
\end{proof}

\begin{theorem}[О сложности]
В худшем случае алгоритм Эрли работает за время $\cal{O}(n^3)$ и использует память $\cal{O}(n^2)$ для грамматики фиксированного размера.
\end{theorem}
\begin{proof}
Каждая ситуация $[A \to \alpha \cdot \beta, i]$ определяется нетерминалом $A$, позицией в правиле, начальной позицией $i$ и текущей позицией $j$. Число возможных ситуаций есть $\cal{O}(n^2 \cdot |G|)$. Обработка одной ситуации (операции предсказания и завершения) может, в худшем случае, породить $\cal{O}(n)$ новых ситуаций (из-за цикла по $k$ в завершении). Итоговая верхняя оценка времени работы: $\cal{O}(n^2 \cdot |G|) \cdot \cal{O}(n) = \cal{O}(n^3 \cdot |G|)$. Для фиксированной грамматики $|G|$ — константа, что даёт $\cal{O}(n^3)$.
\end{proof}

\begin{example}
Для грамматики $E \to E + T \mid T,\; T \to a$ и строки $a+a$ алгоритм последовательно строит множества ситуаций $S_0, S_1, S_2, S_3$. В $S_3$ появляется ситуация $[E' \to E \cdot, 0]$, что подтверждает принадлежность строки языку.
\TODO{Переписать слайд. Развернуть пример.}
\end{example}
\end{frame}

\begin{frame}
\frametitle{1.4.4. Специальные методы синтаксического анализа (6/9)}
\odmkey{Алгоритм Кока--Янгера--Касами}{Алгоритм!Кока--Янгера--Касами} (CYK) предназначен для синтаксического анализа КС-грамматик, приведённых к {\em нормальной форме Хомского} (НФХ, см. \DMref{п.~1.3.3}{1.3.3.}).
В НФХ каждое правило имеет вид $A \to BC$ или $A \to a$.

Пусть $w = a_1 \dots a_n$ — входная строка. Построим треугольную таблицу $P$ размером $n \times n$, где $P[i,j]$ — множество нетерминалов, порождающих подстроку $a_i \dots a_j$.
\TODO{псевдокод}


\begin{enumerate}
    \item \textbf{Инициализация} для подстрок длины 1:
        \[
        P[i,i] = \{ A \mid A \to a_i \in R \}.
        \]
    \item \textbf{Рекуррентное правило} для длины $l > 1$:
        \[
        P[i,j] = \bigcup_{k=i}^{j-1} \{ A \mid A \to BC \in R,\ B \in P[i,k],\ C \in P[k+1,j] \}.
        \]
\end{enumerate}

Строка принадлежит языку, если $S \in P[1,n]$.

\begin{remark}
    Для корректной работы алгоритма грамматика должна быть приведена к НФХ и не содержать $\varepsilon$-правил (кроме, возможно, $S \to \varepsilon$).
\end{remark}
\end{frame}

\begin{frame}
\frametitle{1.4.4. Специальные методы синтаксического анализа (7/9)}

\begin{theorem}
    Алгоритм CYK корректен и полон для грамматик в НФХ:
    $w \in \mathfrak{L}(G)$ тогда и только тогда, когда $S \in P[1,n]$.
\end{theorem}

\begin{proof}
Индукция по длине подстроки $l = j-i+1$.

\textbf{База:} $l=1$. По определению НФХ, $A \Rightarrow^* a_i$ тогда и только тогда, когда $A \to a_i \in R$, что соответствует инициализации.

\textbf{Шаг:} Если $A \Rightarrow^* a_i \dots a_j$ для $l>1$, то существует правило $A \to BC$ и разбиение подстроки такое, что $B \Rightarrow^* a_i \dots a_k$ и $C \Rightarrow^* a_{k+1} \dots a_j$. По предположению индукции, $B \in P[i,k]$ и $C \in P[k+1,j]$, следовательно, алгоритм добавит $A$ в $P[i,j]$.
\end{proof}

\begin{example}
    Рассмотрим грамматику в НФХ и строку $w = baaba$:
    \[
    \begin{array}{rcl}
    S & \to & AB \mid BC,\\
    A & \to & BA \mid a,\\
    B & \to & CC \mid b,\\
    C & \to & AB \mid a.
    \end{array}
    \]
    \textbf{Инициализация ($l=1$):}
    $P[1,1]=\{B\}$, $P[2,2]=\{A,C\}$, $P[3,3]=\{A,C\}$, $P[4,4]=\{B\}$, $P[5,5]=\{A,C\}$.
\end{example}
\end{frame}

\begin{frame}
\frametitle{1.4.4. Специальные методы синтаксического анализа (8/9)}
\begin{example} (\textit{Продолжение})
Для $l > 1$:
\begin{itemize}
    \item \textbf{$l=2$:} Например, $P[1,2] = \{A, S\}$, т.к. $A \to BA$ ($B \in P[1,1]$, $A \in P[2,2]$) и $S \to BC$ ($B \in P[1,1]$, $C \in P[2,2]$).
    \item \textbf{$l=3$:} $P[2,4] = \{B\}$, т.к. $B \to CC$ ($C \in P[2,3]$, $C \in P[4,4]$? Нет, требуется аккуратная проверка. Фактически, $P[2,4]$ получается из разбиений $k=2$ и $k=3$).
    \item \textbf{$l=5$:} $P[1,5] = \{A, S, C\}$.
\end{itemize}

\vspace{2mm}
Итоговая треугольная таблица $P$:
\[
\begin{array}{c|ccccc}
  & 1 & 2 & 3 & 4 & 5\\
\hline
1 & \{B\} & \{A,S\} & \varnothing & \varnothing & \{A,S,C\}\\
2 &       & \{A,C\} & \{B\} & \{B\} & \{A,S,C\}\\
3 &       &         & \{A,C\} & \{S,C\} & \{B\}\\
4 &       &         &         & \{B\} & \{A,S\}\\
5 &       &         &         &         & \{A,C\}\\
\end{array}
\]

Так как $S \in P[1,5]$, строка $baaba$ принадлежит языку.
\end{example}
\end{frame}

\begin{frame}
\frametitle{1.4.4. Специальные методы синтаксического анализа (9/9)}
\begin{digress}
Алгоритм CYK допускает расширения для решения смежных задач:

\begin{enumerate}
    \item Построение деревьев разбора:
    
    Если в $P[i,j]$ хранить не только нетерминалы, но и соответствующие им поддеревья, то можно восстановить все возможные деревья разбора.
    
    \item Вероятностные КС-грамматики (ВКСГ):
    
    Для грамматик с вероятностями правил $p(A \to \alpha)$ алгоритм вычисляет наиболее вероятный разбор:
    \[
    V[i,j,A] = \max_{\substack{A \to BC \in R \\ i \leq k < j}} p(A \to BC) \cdot V[i,k,B] \cdot V[k+1,j,C]
    \]
    Вероятность строки: $V[1,n,S]$.
\end{enumerate}

\begin{remark}
    Эти расширения демонстрируют гибкость алгоритма CYK, хотя основное его назначение — проверка принадлежности строки языку, порождаемому грамматикой в НФХ.
\end{remark}

\begin{example}
    В задачах обработки естественного языка ВКСГ и модифицированный алгоритм CYK используются для синтаксического анализа предложений.
\end{example}

\end{digress}
\end{frame}


\subsubsection{1.4.4. Специальные методы синтаксического анализа}
\label{1.4.4.}
% Текст слайда подготовил Чинь Тхи Тхи Хоай - 5030102/20202
% 1.4.4. Специальные методы синтаксического анализа
% Дата 06.01.2026
% Ревизия ФАН 15.01.26

\begin{frame}
\frametitle{1.4.4. Специальные методы синтаксического анализа (1/13)}

Методы синтаксического анализа, рассмотренные в предыдущих параграфах (\DMref{п.~1.4.2}{1.4.2.} и \DMref{п.~1.4.3}{1.4.3.}), ориентированы на использование контекстно-свободных грамматик, обладающих дополнительными структурными свойствами, такими как детерминированность, отсутствие левой рекурсии или принадлежность к классам $LL(k)$ и $LR(k)$. Эти ограничения позволяют строить эффективные алгоритмы разбора, однако существенно сужают класс допустимых грамматик.

Для анализа языков, заданных контекстно-свободными грамматиками с другими ограничениями или вовсе без таких ограничений, разработаны специальные методы, основанные на иных математических идеях. Рассматриваются три знаменитых алгоритма:
\begin{enumerate}
    \item \odmkey{Алгоритм Дейкстры}{Алгоритм!Дейкстры}
    (\odmkey{алгоритм сортировочной станции}{Алгоритм!сортировочной станции})%
    \footnote{Э.~Дейкстра (1930--2002)}
    для анализа арифметических выражений;
    \item \odmkey{Алгоритм Эрли}{Алгоритм!Эрли}%
    \footnote{Дж.~Эрли (1945--1980)}
    для произвольных КС-грамматик;
    \item \odmkey{Алгоритм Кока--Янгера--Касами}{Алгоритм!Кока--Янгера--Касами}
    (\odmkey{алгоритм CYK}{Алгоритм!CYK})%
    \footnote{Дж.~Кок, Д.~Янгер, Т.~Касами}
    для грамматик в нормальной форме Хомского.
\end{enumerate}

Эти алгоритмы демонстрируют поучительные компромиссы между общностью, эффективностью и сложностью реализации.
\end{frame}


\begin{frame}
\frametitle{1.4.4. Специальные методы синтаксического анализа (2/13)}

\odmkey{Алгоритм Дейкстры}{Алгоритм!Дейкстры} решает задачу синтаксического анализа
арифметических выражений за один проход, то есть за время $\cal O(n)$.
Его основная идея состоит в использовании двух стеков:
\odmkey{стека операндов}{Стек!операндов} $V$ и
\odmkey{стека операторов}{Стек!операторов} $U$.

При чтении входной цепочки $\omega$ операнды сразу помещаются в стек $V$,
а операторы --- в стек $U$, но с учётом их приоритетов. Если новый оператор
имеет меньший или равный приоритет по сравнению с верхним оператором в стеке,
то сначала выполняется \odmkey{свёртка}{Свёртка}: верхний оператор применяется
к двум верхним операндам.

Тем самым алгоритм неявно строит структуру выражения в соответствии с обычными
правилами приоритета операций и скобок.

\begin{remark}
Этот алгоритм удобен именно для выражений, синтаксис которых задаётся прежде всего
приоритетами и ассоциативностью операций, а не общей КС-грамматикой.
\end{remark}
\end{frame}

\begin{frame}
\frametitle{1.4.4. Специальные методы синтаксического анализа (3/13)}

\textbf{Правила алгоритма Дейкстры.}
Для каждого входного символа $a$ цепочки $\omega$:
\begin{enumerate}
\item Если $a$ является операндом, то поместить его в стек $V$.
\item Если $a$ является оператором, то пока
\[
U \neq \varepsilon
\qquad \text{и} \qquad
\mathrm{Pr}(\mathrm{Top}(U)) \ge \mathrm{Pr}(a),
\]
выполнять свёртку. Затем поместить $a$ в $U$.
\item Если $a=\texttt{(}$, то поместить его в $U$.
\item Если $a=\texttt{)}$, то выполнять свёртки, пока
\[
\mathrm{Top}(U)\neq \texttt{(},
\]
после чего удалить $\texttt{(}$ из $U$.
\item В конце цепочки выполнить все оставшиеся свёртки.
\end{enumerate}

\textbf{Операция свёртки.}
\begin{enumerate}
\item Извлечь оператор $\oplus$ из $U$.
\item Извлечь два верхних операнда из $V$.
\item Поместить в $V$ результат применения операции $\oplus$.
\end{enumerate}
\end{frame}

\begin{frame}
\frametitle{1.4.4. Специальные методы синтаксического анализа (4/13)}

\begin{example}
Рассмотрим выражение
\[
a+b*c
\]
и грамматику
\[
E \to E + T \mid T,\qquad
T \to T * F \mid F,\qquad
F \to a \mid b \mid c.
\]

Алгоритм Дейкстры восстанавливает структуру этого выражения с учётом приоритетов
операций. Ход разбора удобно представить таблицей.
\end{example}

\[
\begin{array}{|c|c|c|c|}
\hline
\text{Шаг} & \text{Входной символ} & V & U \\
\hline
1 & a & [a] & [\,] \\
2 & + & [a] & [+] \\
3 & b & [a,b] & [+] \\
4 & * & [a,b] & [+,*] \\
5 & c & [a,b,c] & [+,*] \\
\hline
\end{array}
\]

После чтения всей цепочки выполняются две свёртки:
сначала по оператору $*$, затем по оператору $+$.
\end{frame}

\begin{frame}
\frametitle{1.4.4. Специальные методы синтаксического анализа (5/13)}

\begin{example}
(\textit{Продолжение})
После конца входной цепочки выполняются оставшиеся свёртки:
\begin{enumerate}
    \item после свёртки по $*$ получаем
    \[
    V=[a,(b*c)],\qquad U=[+];
    \]
    \item после свёртки по $+$ получаем
    \[
    V=[(a+(b*c))],\qquad U=[\,].
    \]
\end{enumerate}

Следовательно, итоговая структура выражения имеет вид
\[
a+(b*c),
\]
то есть сначала вычисляется произведение, а затем сумма.
\end{example}

\begin{remark}
Таким образом, алгоритм Дейкстры полезен в тех случаях, когда синтаксис выражений
определяется в первую очередь приоритетами операций и скобочной структурой.
\end{remark}
\end{frame}

\begin{frame}
\frametitle{1.4.4. Специальные методы синтаксического анализа (6/13)}

\odmkey{Алгоритм Эрли}{Алгоритм!Эрли} представляет собой универсальный метод
синтаксического анализа для произвольных КС-грамматик. Его основная идея состоит
в динамическом построении множеств \odmkey{ситуаций}{Ситуация} $S_j$
для каждой позиции $j$ во входной цепочке
\[
\omega=a_1\dots a_n.
\]

\odmkey{Ситуацией}{Ситуация} называется запись вида
\[
[A \to \alpha \cdot \beta,\ i],
\]
где:
\begin{itemize}
    \item $A \to \alpha\beta \in \Pi$ --- правило грамматики;
    \item точка $\cdot$ указывает текущую позицию в правиле;
    \item $i$ --- позиция во входной цепочке, где началось распознавание
    нетерминала $A$.
\end{itemize}

Смысл ситуации таков: к позиции $j$ уже распознана часть $\alpha$, а часть $\beta$
ещё предстоит распознать.
\end{frame}

\begin{frame}
\frametitle{1.4.4. Специальные методы синтаксического анализа (7/13)}

Алгоритм Эрли инициализирует множество $S_0$ и последовательно строит множества
$S_1,\dots,S_n$ с помощью трёх операций.

\begin{enumerate}
    \item \textbf{Предсказание.}
    Если в $S_j$ имеется ситуация
    $[A \to \alpha \cdot B\beta,\ i]$,
    то нетерминал $B$ должен быть распознан, начиная с позиции $j$.
    Поэтому для каждого правила $B \to \gamma \in \Pi$
    в $S_j$ добавляется ситуация
    $[B \to \cdot \gamma,\ j]$.

    \item \textbf{Сканирование.}
    Если в $S_j$ имеется ситуация
    $[A \to \alpha \cdot a\beta,\ i]$
    и очередной символ входной цепочки равен $a_{j+1}=a$,
    то символ $a$ считается прочитанным, и в $S_{j+1}$ добавляется
    ситуация
    $[A \to \alpha a \cdot \beta,\ i]$.

    \item \textbf{Завершение.}
    Если в $S_j$ имеется завершённая ситуация
    $[A \to \gamma \cdot,\ i]$,
    то распознавание нетерминала $A$, начатое в позиции $i$,
    успешно закончено в позиции $j$.
    Поэтому для всякой ситуации
    $[B \to \alpha \cdot A\beta,\ k]\in S_i$
    в $S_j$ добавляется ситуация
    $[B \to \alpha A \cdot \beta,\ k]$.
\end{enumerate}

\end{frame}

\begin{frame}
\frametitle{1.4.4. Специальные методы синтаксического анализа (8/13)}

Цепочка $\omega$ допускается тогда и только тогда, когда в $S_n$ присутствует
ситуация
\[
[S' \to S \cdot,\ 0],
\]
где $S'$ --- новый стартовый нетерминал, добавленный к грамматике.

\begin{theorem}
Алгоритм Эрли корректен и полон:
\[
\omega \in \mathfrak{L}(G)
\quad \Longleftrightarrow \quad
[S' \to S \cdot,\ 0] \in S_n.
\]
\end{theorem}

\begin{proof}
Индукция по длине вывода. Корректность состоит в том, что каждая добавляемая
ситуация соответствует некоторому возможному фрагменту вывода. Полнота состоит
в том, что для любого вывода
\[
S \Rightarrow^* \omega
\]
соответствующая завершающая ситуация будет построена алгоритмом.
\end{proof}
\end{frame}

\begin{frame}
\frametitle{1.4.4. Специальные методы синтаксического анализа (9/13)}

\begin{theorem}[О сложности]
В худшем случае алгоритм Эрли работает за время $\cal O(n^3)$ и использует память
$\cal O(n^2)$ для грамматики фиксированного размера.
\end{theorem}

\begin{proof}
Каждая ситуация $[A \to \alpha \cdot \beta,\ i]$ определяется правилом грамматики,
положением точки, начальной позицией $i$ и текущей позицией $j$. Число возможных
ситуаций имеет порядок $\cal O(n^2\cdot |G|)$. Обработка одной ситуации в худшем
случае может породить $\cal O(n)$ новых ситуаций. Поэтому общая верхняя оценка
времени работы имеет вид $\cal O(n^3\cdot |G|)$, а для фиксированной грамматики
сводится к $\cal O(n^3)$.
\end{proof}

\begin{example}
Для грамматики
\[
E \to E + T \mid T,\qquad T \to a
\]
и цепочки $\omega=a+a$ алгоритм строит множества ситуаций
$S_0,S_1,S_2,S_3$.
\end{example}
\end{frame}

\begin{frame}
\frametitle{1.4.4. Специальные методы синтаксического анализа (10/13)}

\begin{example}
(\textit{Продолжение})
После добавления нового стартового символа $E'$ получаем правило
\[
E' \to E.
\]
В множестве $S_0$ возникают начальные ситуации
\[
[E' \to \cdot E,0],\qquad
[E \to \cdot E+T,0],\qquad
[E \to \cdot T,0],\qquad
[T \to \cdot a,0].
\]

После чтения первого символа $a$ в множестве $S_1$ появляется ситуация
\[
[T \to a \cdot,0],
\]
а затем, после завершения, и ситуации для нетерминала $E$.

После чтения символа $+$ и второго символа $a$ в $S_3$ появляется ситуация
\[
[E' \to E \cdot,0],
\]
что и подтверждает принадлежность цепочки $\omega=a+a$ языку.
\end{example}
\end{frame}

\begin{frame}
\frametitle{1.4.4. Специальные методы синтаксического анализа (11/13)}

\odmkey{Алгоритм Кока--Янгера--Касами}{Алгоритм!Кока--Янгера--Касами}
(CYK) предназначен для синтаксического анализа КС-грамматик, приведённых к
\odmkey{нормальной форме Хомского}{Нормальная форма!Хомского}
(НФХ, см. \DMref{п.~1.3.3}{1.3.3.}). В НФХ каждое правило имеет вид
\[
A \to BC
\qquad \text{или} \qquad
A \to a.
\]

Пусть $\omega=a_1\dots a_n$ --- входная цепочка. Построим треугольную таблицу
$P$ размера $n\times n$, где $P[i,j]$ обозначает множество нетерминалов,
порождающих подцепочку $a_i\dots a_j$.

\begin{enumerate}
    \item \textbf{Инициализация:}
    \[
    P[i,i]=\{A \mid A \to a_i \in \Pi\}.
    \]

    \item \textbf{Рекуррентное правило:}
    для $i<j$
    \[
    P[i,j]=
    \bigcup_{k=i}^{j-1}
    \{A \mid A \to BC \in \Pi,\ B \in P[i,k],\ C \in P[k+1,j]\}.
    \]
\end{enumerate}

Итак, таблица $P$ последовательно накапливает сведения о том, какие нетерминалы
порождают подцепочки входной цепочки.
\end{frame}

\begin{frame}
\frametitle{1.4.4. Специальные методы синтаксического анализа (12/13)}

\begin{theorem}
Алгоритм CYK корректен и полон для грамматик в НФХ:
\[
\omega \in \mathfrak{L}(G)
\quad \Longleftrightarrow \quad
S \in P[1,n].
\]
\end{theorem}

\begin{proof}
Индукция по длине подцепочки $l=j-i+1$.

\textbf{База:} $l=1$. По определению НФХ,
\[
A \Rightarrow^* a_i
\quad \Longleftrightarrow \quad
A \to a_i \in \Pi.
\]
Это в точности соответствует инициализации таблицы.

Следовательно, для подцепочек длины $1$ содержимое таблицы $P$ определяется
корректно. Остаётся проверить, что рекуррентное правило правильно обрабатывает
подцепочки большей длины.
\end{proof}
\end{frame}

\begin{frame}
\frametitle{1.4.4. Специальные методы синтаксического анализа (13/13)}

\begin{proof}[Продолжение доказательства]
\textbf{Шаг:} пусть $l>1$ и
\[
A \Rightarrow^* a_i\dots a_j.
\]
Тогда существует правило $A \to BC$ и такое разбиение подцепочки, что
\[
B \Rightarrow^* a_i\dots a_k,
\qquad
C \Rightarrow^* a_{k+1}\dots a_j.
\]

По предположению индукции имеем
\[
B \in P[i,k],\qquad C \in P[k+1,j],
\]
а значит, алгоритм добавит $A$ в $P[i,j]$.
\end{proof}

\begin{digress}
Если в ячейках таблицы хранить не только нетерминалы, но и указания на способ
их получения, то можно восстанавливать деревья разбора. Для вероятностных
контекстно-свободных грамматик аналогичная схема приводит к поиску наиболее
вероятного разбора.
\end{digress}
\end{frame}


\subsection{1.5. Методы определения семантики}\label{1.5.}
\begin{frame}
\frametitle{1.5. Методы определения семантики}
Определить семантику языка~--- это означает задать
правила приписывания смысла всем или некоторым 
словам или предложениям языка (\DMref{п.~0.1.1}{0.1.1.}).

\smallskip
Для конечных языков задать семантику нетрудно, например, таблицей.

\begin{example} Англо-русский разговорник задаёт
семантику конечного (и очень небольшого) подмножества
выражений английского языка.
\end{example}

\smallskip

Для бесконечных языков требуются алгоритмические средства
задания потенциально бесконечных соответствий, т.~е.
правила {\em вычисления} смысла предложения.
Такие правила существенно зависят от того,
как именно определён язык и как именно задан смысл.
Универсальных правил задания семантики, применимых во всех случаях,
в настоящее время не известно.

\begin{example} Методы, рассматриваемые в этом разделе,
неприменимы для извлечения смысла (музыки) из нотной записи.
\end{example}

\end{frame}

\subsubsection{1.5.1. Постановка задачи определения семантики}\label{1.5.1.}
\begin{frame}
\frametitle{1.5.1. Постановка задачи определения семантики}
\TODO{Семантика языков, определённых грамматически. Схема. ФАН}

\vspace{5\baselineskip}

\TODO{Общая философия. Абстракция выполнения. 
Доказательство правильности.\\
Императивные и декларативные языки. \\Что требуется от семантики? ФАН}

\bigskip
\TODO{\color{blue} $\approx 5$ атомов, Петрунин}

\end{frame}


% СЛАЙД 1: СЕМАНТИКА ЯЗЫКОВ, ОПРЕДЕЛЁННЫХ ГРАММАТИЧЕСКИ
% Текст слайда подготовил Швачко Никита 5030102/20202
% Дата 11.01.2026
\begin{frame}
\frametitle{1.5.1. Постановка задачи определения семантики (2/6)}

\NB{Схема определения семантики языка, заданного грамматикой:}

\medskip
\begin{center}
\begin{tabular}{ccccc}
\textbf{Исходный текст} & $\xrightarrow{\text{Лексер}}$ & \textbf{Токены} & 
$\xrightarrow{\text{Парсер}}$ & \textbf{Дерево разбора} \\
& & & & $\downarrow$ \\
& & & & \textbf{Семантический анализ} \\
& & & & $\downarrow$ \\
& & & & \textbf{Смысл (значение)}
\end{tabular}
\end{center}

\begin{remark}
Синтаксис языка определяет \emph{форму} предложений (что корректно записано).\\
Семантика языка определяет \emph{смысл} предложений (что это означает).
\end{remark}

\begin{example}
Для арифметического выражения $2 + 3 * 4$:\\
~--- синтаксис определяет дерево разбора с учётом приоритета операций,\\
~--- семантика определяет результат вычисления: $14$.
\end{example}

\bigskip
\TODO{Здесь и далее неплохой ПЛАН. \\
Уточнить и реализовать. }

\end{frame}



% СЛАЙД 2: ОБЩАЯ ФИЛОСОФИЯ И АБСТРАКЦИЯ ВЫПОЛНЕНИЯ
% Текст слайда подготовил Швачко Никита 5030102/20202
% Дата 11.01.2026
\begin{frame}
\frametitle{1.5.1. Постановка задачи определения семантики (3/6)}

\odmkey{Абстракция выполнения}{Абстракция!выполнения}~--- 
формальное описание того, как программа изменяет состояние 
вычислительной системы, без привязки к конкретной реализации.

\begin{digress}
Зачем нужна формальная семантика?
\begin{itemize}
\item \textbf{Точная спецификация языка}~--- устранение двусмысленностей
\item \textbf{Доказательство свойств программ}~--- верификация, корректность
\item \textbf{Эквивалентность реализаций}~--- разные компиляторы дают одинаковый результат
\item \textbf{Оптимизация}~--- преобразования, сохраняющие смысл
\end{itemize}
\end{digress}

\begin{remark}
Доказательство правильности программы~--- это математическое обоснование того,
что программа соответствует своей спецификации для всех допустимых входов.
\end{remark}

\end{frame}



% СЛАЙД 3: ИМПЕРАТИВНЫЕ И ДЕКЛАРАТИВНЫЕ ЯЗЫКИ
% Текст слайда подготовил Швачко Никита 5030102/20202
% Дата 11.01.2026
\begin{frame}
\frametitle{1.5.1. Постановка задачи определения семантики (4/6)}

\NB{Два подхода к определению смысла программ:}

\medskip
\begin{tabular}{|p{65mm}|p{65mm}|}
\hline
\textbf{Императивные языки} & \textbf{Декларативные языки} \\
\hline
Описывают \emph{как} вычислять & Описывают \emph{что} вычислять \\
\hline
Последовательность команд & Система уравнений или правил \\
\hline
Изменяемое состояние (переменные) & Неизменяемые значения \\
\hline
Семантика: переходы состояний & Семантика: подстановки, редукции \\
\hline
Примеры: C, Java, Python & Примеры: Haskell, Prolog, SQL \\
\hline
\end{tabular}

\begin{remark}
Для императивных языков естественна \emph{операционная} семантика 
(описание через абстрактную машину).\\
Для декларативных языков естественна \emph{денотационная} семантика 
(описание через математические функции).
\end{remark}

\end{frame}



% СЛАЙД 4: МЕТОДЫ ОПРЕДЕЛЕНИЯ СЕМАНТИКИ (ОБЗОР)
% Текст слайда подготовил Швачко Никита 5030102/20202
% Дата 11.01.2026
\begin{frame}
\frametitle{1.5.1. Постановка задачи определения семантики (5/6)}

\NB{Основные методы определения семантики:}

\begin{enumerate}
\item \odmkey{Атрибутные грамматики}{Грамматика!атрибутная} (\DMref{п.~1.5.2}{1.5.2.})~--- 
семантика вычисляется как атрибуты узлов дерева разбора

\item \odmkey{Синтаксически управляемый перевод}{Перевод!синтаксически управляемый} (\DMref{п.~1.5.3}{1.5.3.})~--- 
семантика как перевод в другой язык (например, в байт-код)

\item \odmkey{Операционная семантика}{Семантика!операционная} (\DMref{п.~1.5.5}{1.5.5.})~--- 
описание через выполнение на абстрактной машине

\item \odmkey{Денотационная семантика}{Семантика!денотационная} (\DMref{п.~1.5.7}{1.5.7.})~--- 
сопоставление программе математической функции

\item \odmkey{Исчисления программ}{Исчисление!программ} (\DMref{п.~1.5.8}{1.5.8.})~--- 
описание через пред- и постусловия (аксиомы Хоара)
\end{enumerate}

\end{frame}


% =====================================================
% СЛАЙД 5: ПРИМЕР СЕМАНТИЧЕСКОГО АНАЛИЗА
% =====================================================
% Текст слайда подготовил Швачко Никита 5030102/20202
% Дата 11.01.2026
\begin{frame}
\frametitle{1.5.1. Постановка задачи определения семантики (6/6)}

\begin{example}
Рассмотрим семантику оператора присваивания $x := e$ в императивном языке.

\NB{Операционная семантика:}
\begin{itemize}
\item Вычислить значение выражения $e$ в текущем состоянии $\sigma$
\item Обновить состояние: $\sigma' = \sigma[x \mapsto \text{значение}(e, \sigma)]$
\end{itemize}

\NB{Денотационная семантика:}
\[
[\![ x := e ]\!](\sigma) = \sigma[x \mapsto [\![ e ]\!](\sigma)]
\]

\NB{Деривационная семантика (исчисление программ):}
\[
\Prog{P}{x \Assign e}{P \And x = e}
\]
\end{example}

\begin{remark}
Все три подхода описывают одну и ту же семантику разными способами.
Выбор метода зависит от целей: реализация, верификация или спецификация.
\end{remark}

\end{frame}



\subsubsection{1.5.2. Атрибутные грамматики}\label{1.5.2.}
\begin{frame}
\frametitle{1.5.2. Атрибутные грамматики}
\TODO{Идея Кнута. \\Синтезированные и унаследованные атрибуты.
\\Пример = школьный калькулятор.\\
Пример = символьное дифференцирование. }

\bigskip
\TODO{{\color{red}Замечания ФАН 20 января:} \\
По теме и структуре почти хорошо, по качеству выполнения совершенно никуда не годится. 
Это безобразная халтура, позорище. Делал Олег Жохов. \\
Неправильные обозначения, неверное форматирование, невнятные слова, недоделанные примеры.\\
Далее: не выполнено задание (см. выше), нет позиционирования, нет ссылки на Кнута, нет финальных выводов.\\
Что делать: полностью переделать и расширить до 9 слайдов.}
\end{frame}

% Текст слайда подготовил Жохов Олег 6.01.2026
    \begin{frame}
        \frametitle{1.5.2. Атрибутные грамматики (1/7)}
        Рассмотрим формальную КСГ $G =$ ⟨$N, T, R, S$⟩. Она хорошо описывает 
        синтаксис языка, но её недостаточно для описания \odmkey{семантики}{семантики} языка, 
        то есть значений в нём. Для описания семантики введем \odmkey{атрибуты}{атрибуты} - дополнительные данные, ассоциированные с грамматическими символами.
        
        \odmkey{Атрибутной грамматикой}{Атрибутной!грамматикой} называется четверка $AG =$ ⟨$G, A, F$⟩, где $А$ – конечное множество всех атрибутов, $A = A_0 + A_1$, 
        где $A_0$ ‐ конечное множество синтезируемых атрибутов, $A_1$ ‐ конечное множество наследуемых атрибутов,  $F$ ‐ конечное множество семантических правил. 
        
        
        \remark $A_0 \cap A_1 = \oslash$ 
        
        Рассмотрим атрибутную грамматику $AG$. С каждым символом $X \in V, V = 
        N \cup T$, связывается конечное множество атрибутов $A$, которое делится на 
        $A_0$ и $A_1$.  Будем обозначать атрибут $\alpha$ символа $X$ как $\alpha(X)$. Каждое $\alpha(X) \in A(X)$ имеет множество значений $V_\alpha$. Принадлежность атрибута $\alpha$ символу $X$ обозначим как $X.\alpha$.
        \remark  $A_1(S) = \oslash, A_0(t) = \oslash$, где $t \in T$
        \remark Символ $\alpha$ не имеет отношения к цепочкам.
    \end{frame}


% Текст слайда подготовил Жохов Олег 6.01.2026
    \begin{frame}
    \frametitle{1.5.2. Атрибутные грамматики (2/7)}
    
        Пусть $F$ состоит из набора $m$ правил, где $p$-е правило имеет вид:
        $X_{p0} \to X_{p1}X_{p2}...X_{pn_{p}}$, где $n_p \geq 0, X_{p0} \in N, X_{pj} \in N$ при $1 \leq j \leq n_p$
        \example $p$-е правило имеет вид $X_{p0} \to X_{p1}X_{p2}...X_n$
        
        \odmkey{Семантическими правилами}{Семантическими!правилами}  называются функции $f_{pja}$, определенные для всех $1 \leq p \leq m, 0 \leq j \leq n_p$ и $\alpha \in A_0(X_{pj})$, если $j = 0$; или $\alpha \in A_1(X_{pj})$, если $j > 0$.
        
        $f_{pj\alpha}: V_{\alpha1} \times V_{\alpha2} \times...\times V_{\alpha t} \to V_{\alpha}$ для $t = t(p,j,\alpha) \geq 0$, где все  $\alpha_i = \alpha_i(p,j,\alpha)$ являются атрибутами $X_{pk_i}$ при $0 \leq k_i = k_i(p,j,\alpha) \leq n_p, 1 \leq i \leq t$. 
        
        \example 
        
        $X_{p0}.\alpha_i = f_p,_{p0},_{\alpha_i}(X_{p1}.\alpha_{j_1},X_{p1}.\alpha_{j_2},...,X_{p1}.\alpha_{j_{t_1}},X_{p2}.\alpha_{j_1}, X_{p1}.\alpha_{j_2},...,X_{p1}.\alpha_{j_{j_2}},...,X_{pk}.\alpha_{j_{t_k}})$
        
        Здесь $X_{p0}.\alpha_{i} - $ $i$-й атрибут символа $X_{p0}$, где $1 \leq p \leq m, 0 \leq p_0 \leq n_p, \alpha_i \in A(X_{p0})$; $f_p,_{p0},_{\alpha_i} - $ функция вычисления атрибута $\alpha_i$ символа $X_{p_0}$ правила вывода $p$; $p_\beta - $ индексы символов правила $p$, где $1 \leq \beta \leq k \leq n_p$; $\alpha_{j_{t_{1}}},...,\alpha_{j_{t_{k}}} - $ атрибуты символов $X_{p_\beta}$, где $1 \leq \beta \leq k \leq n_p, t_\beta \leq t$  $\forall \beta \in [1,k]$
        \remark Последняя запись означает, что необязательно каждый символ $X_{p\beta}$ будет иметь одинаковое число атрибутов. $t_\beta$, варьируется от символа к символу и не превышает общего количества атрибутов своего символа.
    
    \end{frame}

% Текст слайда подготовил Жохов Олег 6.01.2026
    \begin{frame}
        \frametitle{1.5.2. Атрибутные грамматики (3/7)}
        
        \begin{figure}[h!]
            \centering
            \begin{minipage}{0.3\textwidth} 
                \centering
                \includegraphics[width=\linewidth]{1.5.2_1.png} 
                \label{fig:1}
            \end{minipage}
            \hfill 
            \begin{minipage}{0.65\textwidth}
                Семантические правила используются для сопоставления цепочкам КС языка «значения» при помощи дерева вывода. Для любого вывода терминальной цепочки $t \in S$ при помощи синтаксических правил построим обычное дерево вывода. Корнем дерева будет $S$, каждый узел помечается либо терминальным символом, либо нетерминалом $X_{p0}$, 
            \end{minipage}
        \end{figure}
        
        соответствующим применению $k$-го правила для некоторого $p$. В последнем случае у этого узла будет $n_p$ прямых потомков.
        
        Пусть теперь $X$ – метка некоторого узла дерева, $\alpha \in A(X)$ – атрибут $Х$. Если $\alpha \in A_0(X)$, то $X = X_{po}$ для некоторого $p$, если же $\alpha \in A_1(X)$, то $X = X_{pj}$ для некоторых $j$ и $p$. В обоих случаях дерево на участке этого узла имеет вид, представленный в схеме. По определению атрибут $\alpha$ имеет в этом узле значение $\upsilon$ = $f_{pj\alpha}(\upsilon_1,...,\upsilon_t)$, если в соответствующем семантическом правиле $f_{pj\alpha}: V_{\alpha1} \times V_{\alpha2} \times...\times V_{\alpha t} \to V_{\alpha}$ все атрибуты $\alpha_1,...,\alpha_t$ уже определены и имеют в узлах с метками $X_{pk_1},...,X_{pk_t}$ значения $\upsilon_1,...,\upsilon_t$ соответственно.
    \end{frame}

% Текст слайда подготовил Жохов Олег 10.01.2026
    \begin{frame}{1.5.2. Атрибутные грамматики (4/7)}
        \example Нужно дать точное определение двоичной системы записи чисел. Для этого используем КСГ $G =$ ⟨⟨$B,L,N$⟩ ⟨"$0$", "$1$", "$.$"⟩, $R$, $N$⟩, где:
        $R =$ ⟨$B \to 0, B \to 1, L \to B, L \to LB, N \to L, N \to L\cdot L$⟩
    
        Зададим следующие атрибуты. $B$: $\upsilon(B) \in Q$ - значение, $s(B) \in Q$ - масштаб; $L$: $\upsilon(L) \in Q$ - значение, $s(L) \in Q$ - масштаб, $l \in Z$ - длина; $N$: $\upsilon(B) \in Q$ - значение.
    
        \begin{table}
        \begin{tabular}{cc}
        \centering
        Синтаксические & Семантические\\
        правила & правила\\
        $B \to 0$ &  $\upsilon(B) := 0$\\
        $B \to 1$ & $\upsilon(B) := 2^{s(B)}$\\
        $L \to B$ & $\upsilon(L) := \upsilon(B), s(B) := s(L), l(L) := 1$\\
        $L_1 \to L_2B$ & $\upsilon(L_1) := \upsilon(L_2) + \upsilon(B), s(B) := s(L_1),$\\
         & $s(L_2) := s(L_1) + 1, l(L_1) := l(L_2) + 1$\\
        $N \to L$ & $\upsilon(N) := \upsilon(L), s(L) := 0$\\
        $N \to L_1 \cdot L_2$ & $\upsilon(N) := \upsilon(L_1) + \upsilon(L_2), s(L_1) := 0, s(L_2) := -l(L_2)$
        \end{tabular}
        \end{table}
    
        
    \end{frame}

% Текст слайда подготовил Жохов Олег 10.01.2026
    \begin{frame}{1.5.2. Атрибутные грамматики (5/7)}
        \remark В данной семантике индексы используются для обозначения одного и того же нетерминала в рамках одного правила. 1 - индекс родительского узла, 2 - дочернего.
        Построим дерево разбора с сопоставленными каждой вершине атриубами - \odmkey{атрибутированное}{атрибутированное} дерево разбора для представления числа 13.25 в двоичной системе счисления. 
        \includegraphics[width=\linewidth]{1.5.2_2.png}
    \end{frame}


% Текст слайда подготовил Жохов Олег 10.01.2026
    \begin{frame}{1.5.2. Атрибутные грамматики (6/7)}
        Получили $13.25_{10} = 1101.01_2$. Атрибуты $\upsilon(B), \upsilon(L), \upsilon(N), l(L)$ узлов вычисляются через знаяения атрибутов дочерних узлов. Такие атрибуты называют \odmkey{синтезированными}{синтезированными}. Атрибуты $s(B), s(L)$ узлов вычисляются через значения атрибутов узлов-предков. Это \odmkey{унаследованные}{унаследованные} атрибуты.
        Грамматики, в которых все атрибуты являются синтезированными, называют
         \odmkey{S-атрибутными грамматиками}{S-атрибутными!грамматиками}.
          Грамматики, в которых все наследуемые атрибуты зависят только от наследуемых атрибутов  предка и/или от любых атрибутов левых братьев называют 
         \odmkey{L-атрибутными грамматиками}{L-атрибутными!грамматиками}
        \remark Всякая S-атрибутная грамматика является L атрибутной, поскольку в  
        S-атрибутной грамматике нет наследуемых атрибутов.
        \odmkey{Графом зависимостей $D_p$}{ Графом зависимостей $D_p$} $\forall p \in R$ является ориентированный граф, вершинами которого являются атрибуты символов, соответствующих $p$ и проведена дуга из $X_m.\alpha$ в $X_k.\alpha_j$ тогда и только тогда, когда $\forall p \in R$ $\exists f | $  $X_m.\alpha_i = f(...,X_k.\alpha_j,...)$
           \odmkey{Графом зависимостей $D(\alpha)$}{Графом зависимостей $D(\alpha)$}  выводимой цепочки $S\Rightarrow^*_{AG} \alpha$ является ориентированный граф, полученный объединением графов зависимостей правил вывода $p$, применённых в $\alpha$. Узлы, соответствующие одним и тем же атрибутам одного и того же символа в разных частях дерева, связанные правилами, отождествляются
    \end{frame}


% Текст слайда подготовил Жохов Олег 10.01.2026
    \begin{frame}{1.5.2. Атрибутные грамматики (7/7)}


        \odmkey{Граф зависимостей $D'(T)$}{Граф зависимостей $D'(T)$} - это ориентированный граф с множеством узлов $A(X)$, в котором дуга из $\alpha$ в $\alpha'$ существует тогда и только тогда, когда в полном графе 
        зависимостей $D(T)$ существует ориентированный путь из 
        узла $(X, \alpha)$ в узел $(X, \alpha')$ для дерева вывода $T$ с корнем $X$.
        
        \example Граф зависимостей для дерева из слайда 5/7.
        
        \begin{figure}[h!]
            \centering
            \begin{minipage}{0.6\textwidth}  
                \centering
                \includegraphics[scale = 0.5]{1.5.2_3.png} 
            \end{minipage}
            \hfill 
            \begin{minipage}{0.35\textwidth} 
                 Атрибутная грамматика называется \odmkey{незацикленной}{незацикленной} или \odmkey{корректной}{корректной}, если граф зависемостей $D'(T)$ не содержит циклов.
                 \remark В примере нет циклов. Грамматика корректна.
            \end{minipage}
        \end{figure}

    \end{frame}



\subsubsection{1.5.3. Синтаксически управляемый перевод}\label{1.5.3.}
\begin{frame}
\frametitle{1.5.3. Синтаксически управляемый перевод}
\TODO{Семантика перевода (трансляции). \\Грамматика анализа и грамматика синтеза.\\
Пример = перевод в обратную польскую запись.\\
См. Ахо Ульман \\
\color{blue} Бурков Александр}


\end{frame}


\subsubsection{1.5.4. Контекстные условия}\label{1.5.4.}
\begin{frame}
\frametitle{1.5.4. Контекстные условия}
\TODO{Что невыразимо в КСГ. Перенос части синтаксисиса в семантику. Проверка типов. Грамматики Ван-Вейнгаардена}

\end{frame}


\subsubsection{1.5.5. Операционная семантика}\label{1.5.5.}
\begin{frame}
\frametitle{1.5.5. Операционная семантика}
\TODO{Абстрактная машина. Перевод в команды интерпретатора. \\
Семантические процедуры, распределитель семантик.}

\begin{center}
 \includegraphics[width = 15cm]{1_5_5_1.jpg} 
\end{center}

\end{frame}

\subsubsection{1.5.6. Венский метод определения языков программирования}\label{1.5.6.}
\begin{frame}
\frametitle{1.5.6. Венский метод определения языков программирования}
\TODO{\color{blue} Колосков Александр. Доклад сделан, но материал для вклейки в слайды пока не готов}
\end{frame}

\subsubsection{1.5.7. Денотационная семантика}\label{1.5.7.}
\begin{frame}
\frametitle{1.5.7. Денотационная семантика}
\TODO{$\lambda$-исчисление.
Денотационная семантика.}

\bigskip
\TODO{\color{blue} $\approx 25$ атомов. Петрунин}
\end{frame}

\subsubsection{1.5.8. Исчисления программ}\label{1.5.8.}
\begin{frame}
\frametitle{1.5.8. Исчисления программ}
\TODO{Аксионы Хоара\\
\color{blue} Величко Арсений }

\bigskip
\TODO{\color{blue} $\approx 25$ атомов. Кто это сделал?}
\end{frame}


% СЛАЙД 1: ИСЧИСЛЕНИЕ ПРОГРАММ — ВВЕДЕНИЕ
% Текст слайда подготовил Швачко Никита 5030102/20202
% Дата 11.01.2026
\begin{frame}
\frametitle{1.5.8. Исчисления программ (2/9)}

\odmkey{Исчисление программ}{Исчисление!программ}~--- формальная система 
для доказательства свойств программ, предложенная Ч.~Хоаром (1969).

\begin{remark}
Исчисление программ~--- это частный случай исчисления предикатов, 
в котором формулы описывают связь между состоянием программы 
до и после выполнения оператора.
\end{remark}

\NB{Основная формула исчисления программ:}
\[
\Prog{P}{S}{Q}
\]
где:
\begin{itemize}
\item $P$~--- \odmkey{предусловие}{Предусловие} (условие на состояние до выполнения $S$)
\item $S$~--- оператор (фрагмент программы)
\item $Q$~--- \odmkey{постусловие}{Постусловие} (условие на состояние после выполнения $S$)
\end{itemize}

\end{frame}



% СЛАЙД 2: СЕМАНТИКА ТРИПЛЕТА ХОАРА
% Текст слайда подготовил Швачко Никита 5030102/20202
% Дата 11.01.2026
\begin{frame}
\frametitle{1.5.8. Исчисления программ (3/9)}

\NB{Содержательный смысл формулы $\Prog{P}{S}{Q}$:}

Если до выполнения оператора $S$ программа находилась в состоянии $w$, 
для которого выполнено предусловие $P(w)$, то после выполнения $S$ 
программа перейдёт в состояние $w' = S(w)$, в котором выполнено постусловие $Q(w')$.

\[
\Prog{P}{S}{Q} \Def \A{w}{P(w) \Impl \E{w'}{w' = S(w) \And Q(w')}}
\]

\begin{example}
$\Prog{x = 5}{x \Assign x + 1}{x = 6}$~--- корректная формула.

$\Prog{x > 0}{x \Assign x * 2}{x > 0}$~--- также корректная формула.
\end{example}

\begin{remark}
Различают \emph{частичную} корректность (если $S$ завершается, то $Q$ выполнено) 
и \emph{полную} корректность (дополнительно доказывается завершаемость $S$).
\end{remark}

\end{frame}



% СЛАЙД 3: АКСИОМЫ И ПРАВИЛА ХОАРА
% Текст слайда подготовил Швачко Никита 5030102/20202
% Дата 11.01.2026
\begin{frame}
\frametitle{1.5.8. Исчисления программ (4/9)}

\NB{Основные аксиомы и правила вывода:}

\begin{enumerate}
\item \textbf{Аксиома присваивания:}
$\Prog{P}{v \Assign e}{P \And v = e}$

\item \textbf{Правило последовательности:}
$\Rule{\Prog{P}{S_1}{Q} \quad \Prog{Q}{S_2}{R}}{\Prog{P}{S_1; S_2}{R}}{1}$

\item \textbf{Правило условного оператора:}
$\Rule{\Prog{P \And B}{S_1}{Q} \quad \Prog{P \And \Not B}{S_2}{Q}}
{\Prog{P}{\textbf{if } B \textbf{ then } S_1 \textbf{ else } S_2}{Q}}{2}$

\item \textbf{Правило цикла:}
$\Rule{\Prog{I \And B}{S}{I}}
{\Prog{I}{\textbf{while } B \textbf{ do } S}{I \And \Not B}}{3}$

где $I$~--- \odmkey{инвариант цикла}{Инвариант!цикла}
\end{enumerate}

\end{frame}



% СЛАЙД 4: ПРИМЕНЕНИЕ И СВЯЗЬ С ДРУГИМИ РАЗДЕЛАМИ
% Текст слайда подготовил Швачко Никита 5030102/20202
% Дата 11.01.2026
\begin{frame}
\frametitle{1.5.8. Исчисления программ (5/9)}

\NB{Применение исчисления программ:}
\begin{itemize}
\item Формальная верификация программ
\item Доказательство корректности алгоритмов
\item Автоматическая генерация тестов
\item \odmkey{Аннотирование программ}{Аннотирование!программ}~--- 
добавление утверждений в виде комментариев
\end{itemize}

\begin{digress}
На практике полностью формальное доказательство корректности 
трудоёмко и применяется редко. 
Однако неформальное обоснование с использованием идей Хоара 
(инварианты, пред- и постусловия) чрезвычайно полезно 
при разработке и отладке программ.
\end{digress}

\end{frame}


% =====================================================
% СЛАЙД 5: ДОПОЛНИТЕЛЬНЫЕ ПРАВИЛА ВЫВОДА
% По материалам раздела 3.3.2
% =====================================================
% Текст слайда подготовил Швачко Никита 5030102/20202
% Дата 11.01.2026
\begin{frame}
\frametitle{1.5.8. Исчисления программ (6/9)}

\NB{Дополнительные правила вывода:}

Правила эквивалентных преобразований условий:
\[
\Rule{P \Equiv P_1 \quad \Prog{P_1}{S}{Q_1} \quad Q_1 \Equiv Q}{\Prog{P}{S}{Q}}{4}
\]

Правила усиления предусловия и ослабления постусловия:
\[
\Rule{P \Impl P_1 \quad \Prog{P_1}{S}{Q_1} \quad Q_1 \Impl Q}{\Prog{P}{S}{Q}}{5}
\]

\begin{remark}
Правило (5) позволяет использовать более сильное предусловие 
и более слабое постусловие, что упрощает доказательства.
\end{remark}

\end{frame}


% =====================================================
% СЛАЙД 6: ПРИМЕР — ОПЕРАТОР ОБМЕНА ЗНАЧЕНИЯМИ
% По материалам раздела 3.3.2
% =====================================================
% Текст слайда подготовил Швачко Никита 5030102/20202
% Дата 11.01.2026
\begin{frame}
\frametitle{1.5.8. Исчисления программ (7/9)}

\begin{example}
Рассмотрим \odmkey{оператор обмена значениями}{Оператор!обмена значениями} $x :=: y$.

Семантика этого оператора:
\[
\Prog{P(x,y)}{x :=: y}{P(y,x)}
\]

Рассмотрим фрагмент программы сортировки трёх значений:
\begin{algorithmic}
\STATE\textbf{if $x > y$ then $x :=: y$ endif}
\STATE\textbf{if $y > z$ then $y :=: z$ endif}
\STATE\textbf{if $x > y$ then $x :=: y$ endif}
\end{algorithmic}

\NB{Цель:} Доказать, что постусловием является $x \leq y \leq z$.
\end{example}

\end{frame}


% =====================================================
% СЛАЙД 7: ДОКАЗАТЕЛЬСТВО СВОЙСТВА СОРТИРОВКИ
% По материалам раздела 3.3.2
% =====================================================
% Текст слайда подготовил Швачко Никита 5030102/20202
% Дата 11.01.2026
\begin{frame}
\frametitle{1.5.8. Исчисления программ (8/9)}

Введём процедуру $S2(a,b) \Assign \textbf{if } a > b \textbf{ then } a :=: b \textbf{ endif}$.

\NB{Свойство процедуры $S2$:}
\[
\Prog{P(a,b)}{S2(a,b)}{(P(a,b) \Or P(b,a)) \And a \leq b}
\]

\begin{proof}
\begin{tabbing}
1.\= $\Prog{P(a,b) \And a > b}{a :=: b}{P(b,a) \And b > a}$ \quad\= аксиома \\[2pt]
2.\> $P(b,a) \And b > a \Impl (P(a,b) \Or P(b,a)) \And a \leq b$ \> общезначимо \\[2pt]
3.\> $\Prog{P(a,b) \And a > b}{a :=: b}{(P(a,b) \Or P(b,a)) \And a \leq b}$ \> правило 5 \\[2pt]
4.\> $P(a,b) \And \Not(a > b) \Impl (P(a,b) \Or P(b,a)) \And a \leq b$ \> общезначимо \\[2pt]
5.\> $\Prog{P(a,b)}{S2(a,b)}{(P(a,b) \Or P(b,a)) \And a \leq b}$ \> правило 2
\end{tabbing}
\end{proof}

\end{frame}


% =====================================================
% СЛАЙД 8: ПРАКТИЧЕСКОЕ ЗНАЧЕНИЕ
% По материалам раздела 3.3.2
% =====================================================
% Текст слайда подготовил Швачко Никита 5030102/20202
% Дата 11.01.2026
\begin{frame}
\frametitle{1.5.8. Исчисления программ (9/9)}

\begin{digress}
Формальное доказательство даже простой программы (сортировка трёх элементов~--- 
частный случай пузырьковой сортировки) оказывается громоздким.
Доказательство всегда длиннее и сложнее самой программы.

Однако как программисты используют языки высокого уровня вместо машинного кода,
так и математики используют менее формальные, но более выразительные методы.
\end{digress}

\NB{Практическая ценность исчисления программ:}
\begin{itemize}
\item Неформальное обоснование с использованием инвариантов полезно при разработке
\item \odmkey{Аннотирование программ}{Аннотирование!программ}~--- 
вставка утверждений о состоянии в виде комментариев
\item Современные инструменты (Dafny, Frama-C, SPARK) автоматизируют верификацию
\end{itemize}

\begin{remark}
Подробное изложение с полным доказательством~--- 
см. раздел 3.3.2 <<Исчисление программ>>.
\end{remark}

\bigskip
\TODO{хороший план\\ Как перераспределить материал между частями курса? \\
Что можно сказать про <<современные инструменты>> ? }

\end{frame}



\subsubsection{1.5.9. Преобразователи предикатов}\label{1.5.9.}
\begin{frame}
\frametitle{1.5.9. Преобразователи предикатов}
\TODO{Преобразователи предикатов Дейкстры.\\
Слабейшие предусловия. \\
Завершаемость. \\
\color{blue} Величко Арсений}

\end{frame}

% Текст слайда подготовил Черняховский Лев и Бурлака Никита 5030102/20202
% Дата 04.04.2026

\begin{frame}
\frametitle{1.5.9. Преобразователи предикатов (1/9)}

В п.~1.5.8 свойства программ задавались тройками Хоара
$\Prog{P}{S}{Q}$.
\odmkey{Преобразователь предикатов}{Преобразователь!предикатов}
решает обратную задачу:
по оператору $S$ и требуемому постусловию $Q$
он вычисляет условие на начальное состояние программы.

\smallskip
Метод предикатных преобразователей был предложен
Э.~Дейкстрой как средство формального вывода программ.

\begin{remark}
Если тройка Хоара отвечает на вопрос
<<что нужно доказать для готовой программы>>,
то преобразователь предикатов отвечает на вопрос
<<что должно быть истинно до выполнения $S$, чтобы после него было истинно $Q$?>>.
\end{remark}

\begin{example}
Для оператора $x \Assign x+1$ и постусловия $x>10$
искомое предусловие равно $x>9$.
\end{example}

\end{frame}

\begin{frame}
\frametitle{1.5.9. Преобразователи предикатов (2/9)}

Пусть $S$ - оператор, а $Q$ - предикат на заключительном состоянии.

\smallskip
\odmkey{Слабейшим предусловием}{Предусловие!слабейшее}
для $Q$ относительно $S$ называется такой предикат
$\mathrm{wp}(S,Q)$, что

\[
\Prog{\mathrm{wp}(S,Q)}{S}{Q},
\]

и, кроме того, для любого предиката $P$ из
$\Prog{P}{S}{Q}$ следует

\[
P \Impl \mathrm{wp}(S,Q).
\]

\begin{remark}
Слово <<слабейшее>> понимается в обычном логическом смысле:
предикат $P_1$ слабее предиката $P_2$, если $P_2 \Impl P_1$.
Следовательно, $\mathrm{wp}(S,Q)$ накладывает на начальное состояние
минимальные ограничения, достаточные для достижения $Q$.
\end{remark}

\end{frame}

\begin{frame}
\frametitle{1.5.9. Преобразователи предикатов (3/9)}

Если требуется доказать не только корректность результата,
но и завершаемость программы, используют предикат
$\mathrm{wp}(S,Q)$.
Если же речь идёт только о \emph{частичной корректности},
то удобно использовать
\odmkey{слабейшее либеральное предусловие}{Предусловие!либеральное}
$\mathrm{wlp}(S,Q)$.

\smallskip
Предикат $\mathrm{wlp}(S,Q)$ характеризует такие начальные состояния,
из которых любой \emph{завершившийся} запуск $S$ приводит к состоянию,
удовлетворяющему $Q$.

\begin{remark}
В терминологии п.~1.5.8
$\mathrm{wp}$ соответствует сильным формулам исчисления программ,
а $\mathrm{wlp}$ - слабым.
Если оператор $S$ всегда завершается,
то $\mathrm{wp}(S,Q)$ и $\mathrm{wlp}(S,Q)$ совпадают.
\end{remark}

\begin{example}
Для бесконечного цикла и любого нетривиального $Q$
предикат $\mathrm{wlp}$ может быть истинным,
тогда как $\mathrm{wp}$ - ложным:
частичная корректность не нарушается,
но завершения нет.
\end{example}

\end{frame}

\begin{frame}
\frametitle{1.5.9. Преобразователи предикатов (4/9)}

Для простейших операторов преобразователи предикатов
вычисляются непосредственно.

\smallskip
Для оператора, не меняющего состояние,
например $\textbf{yield }e$, имеем
\[
\mathrm{wp}(\textbf{yield }e, Q)=Q.
\]

Для присваивания общая форма аксиомы Хоара записывается так:
\[
\Prog{Q[e/x]}{x \Assign e}{Q},
\]
откуда
\[
\mathrm{wp}(x \Assign e, Q)=Q[e/x].
\]

Здесь $Q[e/x]$ обозначает результат подстановки
выражения $e$ вместо всех свободных вхождений переменной $x$
в предикат $Q$.

\begin{example}
\[
\mathrm{wp}(x \Assign x+1,\ x>10)=(x+1>10),
\]
то есть слабейшим предусловием является $x>9$.
\end{example}

\end{frame}

\begin{frame}
\frametitle{1.5.9. Преобразователи предикатов (5/9)}

Для составных операторов преобразование предикатов выполняется
\odmkey{композиционно}{Семантика!композиционная},
по структуре программы.

\begin{enumerate}
\item Для последовательности:
\[
\mathrm{wp}(S_1;S_2,Q)=\mathrm{wp}(S_1,\mathrm{wp}(S_2,Q)).
\]

\item Для условного оператора:
\[
\mathrm{wp}(\textbf{if } B \textbf{ then } S_1 \textbf{ else } S_2 \textbf{ end if},Q)=
\]
\[
(B \Impl \mathrm{wp}(S_1,Q)) \And (\Not B \Impl \mathrm{wp}(S_2,Q)).
\]
\end{enumerate}

\begin{remark}
Вычисление ведётся справа налево:
сначала обрабатывается последний оператор,
затем предпоследний и т.д.
Именно поэтому метод особенно удобен для автоматического
получения условий корректности.
\end{remark}

\end{frame}

\begin{frame}
\frametitle{1.5.9. Преобразователи предикатов (6/9)}

\begin{example}
Пусть
\[
S \Assign x \Assign x+1;\ y \Assign x+y,
\qquad
Q \Assign (y>10).
\]

Тогда слабейшее предусловие для $Q$ вычисляется так:
\[
Q_2 \Assign y>10,
\]
\[
Q_1 \Assign \mathrm{wp}(y \Assign x+y,Q_2)=x+y>10,
\]
\[
Q_0 \Assign \mathrm{wp}(x \Assign x+1,Q_1)=x+1+y>10.
\]

Следовательно,
\[
\mathrm{wp}(S,Q)=x+y>9.
\]
\end{example}

\begin{remark}
Здесь хорошо видно отличие от прямого исполнения программы:
мы не вычисляем значения переменных,
а \emph{обратно переносим} требование к конечному состоянию
на начало фрагмента.
\end{remark}

\end{frame}

\begin{frame}
\frametitle{1.5.9. Преобразователи предикатов (7/9)}

Для цикла точное выражение $\mathrm{wp}$ или $\mathrm{wlp}$
обычно не удаётся записать простой замкнутой формулой.
Поэтому цикл снабжают
\odmkey{инвариантом}{Инвариант!цикла}
$I$.

\smallskip
Для оператора
\[
\textbf{while } B \textbf{ do } S \textbf{ end while}
\]
достаточно проверить следующие условия:

\[
P \Impl I,
\]
\[
I \And B \Impl \mathrm{wlp}(S,I),
\]
\[
I \And \Not B \Impl Q.
\]

\begin{remark}
Первое условие связывает исходное предусловие с инвариантом.
Второе означает, что одна итерация цикла сохраняет инвариант.
Третье позволяет после выхода из цикла получить требуемое постусловие.
\end{remark}

\end{frame}

\begin{frame}
\frametitle{1.5.9. Преобразователи предикатов (8/9)}

Чтобы перейти от частичной корректности к полной,
нужно доказать завершаемость цикла.
Для этого вводят
\odmkey{вариант цикла}{Вариант!цикла}
$t$.

\smallskip
Вариант должен обладать двумя свойствами:

\begin{enumerate}
\item при выполнении $I \And B$ он принимает значения из
множества $\mathbb{N}$;

\item каждая итерация цикла уменьшает его значение.
\end{enumerate}

Например, второе условие можно записать так:
\[
\A{n\in\mathbb{N}}{
\Prog{I \And B \And t=n}{S}{I \And t<n}
}.
\]

\begin{remark}
В терминологии Дейкстры условие $I$ называется
\emph{инвариантным отношением},
а функция $t$ - \emph{вариантной функцией}.
Наличие варианта исключает бесконечное выполнение цикла.
\end{remark}

\end{frame}

\begin{frame}
\frametitle{1.5.9. Преобразователи предикатов (9/9)}

Преобразователи предикатов тесно связаны
с исчислением программ.

\smallskip
Для полной и частичной корректности
тройки Хоара $\Prog{P}{S}{Q}$
эквивалентны требованиям:
\[
\text{полная: } P \Impl \mathrm{wp}(S,Q),
\qquad
\text{частичная: } P \Impl \mathrm{wlp}(S,Q).
\]

\NB{Практический смысл метода:}
\begin{itemize}
\item позволяет вычислять недостающие предусловия по желаемому результату;
\item сводит доказательство корректности к проверке логических импликаций;
\item приводит к генерации условий корректности для автоматической проверки.
\end{itemize}

\begin{remark}
Тем самым в дополнение к п.~1.5.8,
где утверждения задавались и затем доказывались,
здесь появляется новый инструмент~--- их
\emph{обратное вычисление} по структуре программы.
\end{remark}

\end{frame}



\subsection{1.6. Системы программирования}\label{1.6.}
\begin{frame}
\frametitle{1.6. Системы программирования}
\TODO{Лексический и синтаксичесикй анализ. Оптимизация синтаксического анализа. \\
Препроцессоры и шаблоны. \\
Пошаговая компиляция и интерпретация. Обработка синтаксических ошибок.\\
Генерация кода. Оптимизация. \\
Административная система времени выполнения.\\
Тестирование и отладка. \\
Рефакторинг, автозаполнение. \\  
Построение моделей 
программ из кода. Потоковый анализ.
Глобальная экономия памяти.}


\bigskip
\TODO{{\color{red} Замечание ФАН 22 января} \\
Расширить, обновить, согласовать с Соломатиным.}

\end{frame}



